% Complete reviewed Unit 17; editorial master: chapters/ch17/student.md
% Markdown SHA-256: 4e832c39606a72ef832cb726f1c3722bd13ccff0db23c02c007924deb6d3a361
% One-use initial assembly followed by complete content and page review.
\newcommand{\runningtitle}{Unit 17. Exposition and praise}
% Layout inherited from the reviewed previous course; set a running title per unit.
% This is a small hand-maintained TeX layout, not a conversion pipeline.
\RequirePackage{fix-cm}
\documentclass[11pt,twoside]{article}
\usepackage{fontspec}
\setmainfont{Cambria}[
  Renderer=HarfBuzz,
  BoldFont=cambriab.ttf,
  ItalicFont=cambriai.ttf,
  BoldItalicFont=cambriaz.ttf
]
\setmonofont{lmmono10-regular.otf}[BoldFont=lmmonolt10-bold.otf]
\usepackage[paperwidth=165mm,paperheight=241mm,inner=21mm,outer=19mm,
  top=22mm,bottom=20mm,headheight=14pt,headsep=7mm,footskip=10mm]{geometry}
\usepackage{xcolor,array,booktabs,longtable,tabularx}
\usepackage{enumitem,fancyhdr}
\usepackage[unicode,hidelinks]{hyperref}
\definecolor{ink}{HTML}{000000}
\definecolor{muted}{HTML}{4C5558}
\setlength{\parindent}{0pt}
\setlength{\parskip}{3.5pt}
\AtBeginDocument{\fontsize{11}{14.3}\selectfont}
\setlength{\tabcolsep}{4pt}
\renewcommand{\arraystretch}{1.10}
\setlist[itemize]{leftmargin=1.3em,label=\textbullet,itemsep=2pt,topsep=3pt,parsep=0pt}
\setlist[enumerate]{leftmargin=1.8em,itemsep=3pt,topsep=4pt,parsep=0pt}
\widowpenalty=10000
\clubpenalty=10000
\emergencystretch=1em
\raggedbottom
\pagestyle{fancy}
\fancyhf{}
\providecommand{\runningtitle}{Old English}
\fancyhead[LE,RO]{\small\itshape \runningtitle}
\fancyhead[LO,RE]{\small Old English}
\fancyfoot[LE,RO]{\small\thepage}
\renewcommand{\headrulewidth}{0.3pt}
\newcommand{\booktitle}[2]{
  {\fontsize{9}{11}\selectfont\color{muted}#1\par}
  \vspace{3pt}
  {\fontsize{20}{24}\selectfont\color{ink}#2\par}
  \vspace{6pt}
}
\newcommand{\source}[1]{\par{\footnotesize\color{muted}#1\par}}
\newcommand{\readingtext}[1]{
  \par\addvspace{5pt}
  \begingroup\leftskip=1.5em
    \fontsize{12}{16}\selectfont #1\par
  \endgroup\addvspace{5pt}
}
\makeatletter
\renewcommand{\section}{\@startsection{section}{1}{\z@}
  {-2.4ex \@plus -0.6ex \@minus -0.2ex}{0.8ex \@plus 0.2ex}
  {\normalfont\fontsize{12}{15}\selectfont\bfseries\color{ink}}}
\renewcommand{\subsection}{\@startsection{subsection}{2}{\z@}
  {-1.8ex \@plus -0.5ex \@minus -0.2ex}{0.5ex \@plus 0.2ex}
  {\normalfont\fontsize{11}{14}\selectfont\bfseries\color{ink}}}
\makeatother

\hyphenpenalty=10000
\exhyphenpenalty=10000
\emergencystretch=3em
\newcommand{\itempara}[2]{\par\addvspace{3pt}\noindent\hangindent=1.8em\hangafter=1\makebox[1.8em][l]{#1.}#2\par}
\newcommand{\reserve}[1]{\par\penalty-100\begingroup\dimen0=\pagegoal\advance\dimen0 by-\pagetotal\ifdim\dimen0<#1\newpage\fi\endgroup}
\begin{document}
\addfontfeatures{Ligatures=TeXOff}
\raggedright
\booktitle{UNIT 17}{Connected exposition, praise and figurative language}

% mdparagraph:P001:begin
A connected passage uses familiar clauses for different purposes: to report what happened, compare people or actions, praise, request, and interpret an event. Keep the grammar of each clause distinct from the meaning supplied by its setting. A word can retain its ordinary inflection while taking a new contextual sense.
% mdparagraph:P001:end

% mdparagraph:P002:begin
Work through the clause review first, then study the new systems in manageable groups before the form and composition practice. Cook’s complete joint Group XII, XIII, Exercises 37–43 follow. The final reading contains a complete musical performance, criticism, demonstration and audience response. The four translated examples and the lettered tasks are new teaching material.
% mdparagraph:P002:end

% mdparagraph:P003:begin
\textbf{How to use the help.} A dictionary form gives a word and meaning; you may still need to form the required inflection. A complete expression labelled \textbf{Supplied} stays available while you answer and requires no independent recall of its unfamiliar parts. A familiar form does not automatically supply a new meaning. Reference-only and recognition forms are identified with their tables. In verb tables, 1, 2 and 3 denote person; singular and plural address remain distinct.
% mdparagraph:P003:end

\reserve{7\baselineskip}
% mdparagraph:P004:begin
\section*{1. Follow the clauses}
% mdparagraph:P004:end

% mdparagraph:P005:begin
First identify each finite verb, its subject and its complements. Then identify how the clauses are joined. A relative expression can represent a different role inside its clause from the role of its antecedent outside. A content clause reports a proposition. A comparison can repeat an earlier action with dōn. A condition establishes a setting for the following statement or command.
% mdparagraph:P005:end

% teaching: NEW-U17-GRAM-U17-LITERAL-CONTEXTUAL-INTERPRETATION
% mdparagraph:P006:begin
Interpretation comes after this grammatical account. Distinguish an ordinary word meaning, a supplied idiom, and an image whose force depends on context. A complete translation retains the participants and relations as well as the image; an explanation of the image’s point cannot replace the translated sentence.
% mdparagraph:P006:end

% teaching: NEW-U17-EXAMPLE-01-TEACHING
\reserve{5\baselineskip}
% mdparagraph:P007:begin
\textbf{Se cyning sægde ðæt se mann wǣre gōd.}
% mdparagraph:P007:end

% mdparagraph:P008:begin
The king said that the man was good.
% mdparagraph:P008:end

% mdparagraph:P009:begin
Report the man’s condition at the time of the king’s statement; the subjunctive does not itself pronounce the report false.
% mdparagraph:P009:end

% teaching: NEW-U17-EXAMPLE-02-TEACHING
\reserve{5\baselineskip}
% mdparagraph:P010:begin
\textbf{Se cyning sang bęt ðonne se mann dyde.}
% mdparagraph:P010:end

% mdparagraph:P011:begin
The king sang better than the man did.
% mdparagraph:P011:end

% mdparagraph:P012:begin
Dōn repeats singing and has its own subject and tense; bęt modifies the action.
% mdparagraph:P012:end

% teaching: NEW-U17-EXAMPLE-03-TEACHING
\reserve{5\baselineskip}
% mdparagraph:P013:begin
\textbf{Ic geseah ðone mann ðe sang.}
% mdparagraph:P013:end

% mdparagraph:P014:begin
I saw the man who sang.
% mdparagraph:P014:end

% mdparagraph:P015:begin
The man is the object of saw and the subject inside the relative clause; the invariant particle shows no separate case ending.
% mdparagraph:P015:end

% teaching: NEW-U17-EXAMPLE-04-TEACHING
\reserve{5\baselineskip}
% mdparagraph:P016:begin
\textbf{Gif hē cymð, blissiað.}
% mdparagraph:P016:end

% mdparagraph:P017:begin
If he comes, rejoice.
% mdparagraph:P017:end

% mdparagraph:P018:begin
An open condition can use the indicative; the plural command addresses the listeners.
% mdparagraph:P018:end

% mdparagraph:P019:begin
A subjunctive in a report does not by itself say that the report is false. A past form can be simultaneous with a past statement; anteriority needs further evidence. In a condition, establish time and the speaker’s view of its possibility from the whole context. In an equal comparison, swā … swā gives the two sides; in an unequal comparison, identify what the comparative modifies and the standard introduced by ðonne.
% mdparagraph:P019:end

\reserve{8\baselineskip}
% mdparagraph:P020:begin
\subsection*{A01}
% mdparagraph:P020:end

% item: NEW-U17-A01
% mdparagraph:P021:begin
Analyse the report. The man’s goodness is presented as simultaneous with the king’s past statement.
% mdparagraph:P021:end

% mdparagraph:P022:begin
Se cyning sægde ðæt se mann wǣre gōd.
% mdparagraph:P022:end

% output: NEW-U17-A01-O001
% mdparagraph:P023:begin
\itempara{1}{Identify the two subjects and their finite verbs.}
% mdparagraph:P023:end

% output: NEW-U17-A01-O002
% mdparagraph:P024:begin
\itempara{2}{Classify ðæt and explain the reported time relation.}
% mdparagraph:P024:end

% output: NEW-U17-A01-O003
% mdparagraph:P025:begin
\itempara{3}{Does wǣre require the conclusion that the man was not good? Explain.}
% mdparagraph:P025:end

\reserve{8\baselineskip}
% mdparagraph:P026:begin
\subsection*{A02}
% mdparagraph:P026:end

% item: NEW-U17-A02
% mdparagraph:P027:begin
Analyse the comparison; then replace the first subject by the men while keeping the king singular.
% mdparagraph:P027:end

% mdparagraph:P028:begin
Se mann sang bęt ðonne se cyning dyde.
% mdparagraph:P028:end

% output: NEW-U17-A02-O001
% mdparagraph:P029:begin
\itempara{1}{Translate the sentence.}
% mdparagraph:P029:end

% output: NEW-U17-A02-O002
% mdparagraph:P030:begin
\itempara{2}{State the action repeated by dyde and explain why it is singular.}
% mdparagraph:P030:end

% output: NEW-U17-A02-O003
% mdparagraph:P031:begin
\itempara{3}{Write the complete revised Old English sentence.}
% mdparagraph:P031:end

\reserve{8\baselineskip}
% mdparagraph:P032:begin
\subsection*{A03}
% mdparagraph:P032:end

% item: NEW-U17-A03
% mdparagraph:P033:begin
Analyse the relative clause; then replace the man by the men.
% mdparagraph:P033:end

% mdparagraph:P034:begin
Ic geseah ðone mann ðe sang.
% mdparagraph:P034:end

% output: NEW-U17-A03-O001
% mdparagraph:P035:begin
\itempara{1}{Give the outer case and role of ðone mann.}
% mdparagraph:P035:end

% output: NEW-U17-A03-O002
% mdparagraph:P036:begin
\itempara{2}{Give the role represented by ðe inside the relative clause.}
% mdparagraph:P036:end

% output: NEW-U17-A03-O003
% mdparagraph:P037:begin
\itempara{3}{Write the complete revised Old English sentence.}
% mdparagraph:P037:end

\reserve{8\baselineskip}
% mdparagraph:P038:begin
\subsection*{A04}
% mdparagraph:P038:end

% item: NEW-U17-A04
% mdparagraph:P039:begin
A concerns a possible visit tomorrow. In B the speaker treats his coming now as contrary to fact. Analyse both clauses in each sentence.
% mdparagraph:P039:end

% mdparagraph:P040:begin
\textbf{A.} Gif hē cymð, ic fare.
% mdparagraph:P040:end

% mdparagraph:P041:begin
\textbf{B.} Gif hē cōme, ic fōre.
% mdparagraph:P041:end

% output: NEW-U17-A04-O001
% mdparagraph:P042:begin
\itempara{1}{Translate A and name the mood of each verb.}
% mdparagraph:P042:end

% output: NEW-U17-A04-O002
% mdparagraph:P043:begin
\itempara{2}{Translate B and name the tense and mood of each verb.}
% mdparagraph:P043:end

% output: NEW-U17-A04-O003
% mdparagraph:P044:begin
\itempara{3}{Explain which evidence supplies the time and factual status.}
% mdparagraph:P044:end

\reserve{7\baselineskip}
% mdparagraph:P045:begin
\section*{2. Nouns and adjectives}
% mdparagraph:P045:end

% mdparagraph:P046:begin
Choose case from the word’s work in the clause. A subject is nominative, an ordinary direct object accusative, a noun-dependent possessor genitive, and a recipient or governed complement takes its licensed case. A preposition or a particular verb can change the expected choice. Each table below gives all four cases for its stated number and meaning.
% mdparagraph:P046:end

\reserve{7\baselineskip}
% mdparagraph:P047:begin
\subsection*{Strong masculine nouns}
% mdparagraph:P047:end

% mdparagraph:P048:begin
The ordinary masculine pattern has a bare nominative/\allowbreak{}accusative singular, genitive -es and dative -e; the plural has nominative/\allowbreak{}accusative -as, genitive -a and dative -um. The -e and agent-noun patterns and any changed stems are explained individually.
% mdparagraph:P048:end

\reserve{7\baselineskip}
% mdparagraph:P049:begin
\subsection*{hyll}
% mdparagraph:P049:end

% teaching: NEW-U17-VOC-HYLL-U17-SENSE
\par\begin{minipage}{\linewidth}\raggedright
% mdparagraph:P050:begin
\textbf{hyll} — hill.
% mdparagraph:P050:end

% teaching: NEW-U17-NOM-HYLL
% mdparagraph:P051:begin
Masculine; stem \textbf{hyll-}. Add the endings of the stated pattern to this stem, preserving its vowel length.
% mdparagraph:P051:end

% table: NEW-U17-NOM-HYLL-DISPLAY
% mdtable:T01:begin
\par\begingroup
\fontsize{10.5}{13.2}\selectfont
\setlength{\LTleft}{0pt}\setlength{\LTright}{\fill}
\noindent
\begin{tabular}{@{}>{\raggedright\arraybackslash}p{39.46mm}>{\raggedright\arraybackslash}p{39.46mm}>{\raggedright\arraybackslash}p{39.46mm}@{}}
\toprule
\textbf{Case} & \textbf{Singular} & \textbf{Plural} \\
\midrule
Nominative & hyll & hyllas \\
Genitive & hylles & hylla \\
Dative & hylle & hyllum \\
Accusative & hyll & hyllas \\
\bottomrule
\end{tabular}
\par\endgroup

% mdtable:T01:end
\end{minipage}\par
\reserve{7\baselineskip}
% mdparagraph:P052:begin
\subsection*{cēderbēam}
% mdparagraph:P052:end

% teaching: NEW-U17-VOC-CEDERBEAM-U17-SENSE
\par\begin{minipage}{\linewidth}\raggedright
% mdparagraph:P053:begin
\textbf{cēderbēam} — cedar; cedar tree.
% mdparagraph:P053:end

% teaching: NEW-U17-NOM-CEDERBEAM
% mdparagraph:P054:begin
Masculine; stem \textbf{cēderbēam-}. Add the endings of the stated pattern to this stem, preserving its vowel length.
% mdparagraph:P054:end

% table: NEW-U17-NOM-CEDERBEAM-DISPLAY
% mdtable:T02:begin
\par\begingroup
\fontsize{10.5}{13.2}\selectfont
\setlength{\LTleft}{0pt}\setlength{\LTright}{\fill}
\noindent
\begin{tabular}{@{}>{\raggedright\arraybackslash}p{36.99mm}>{\raggedright\arraybackslash}p{40.69mm}>{\raggedright\arraybackslash}p{40.69mm}@{}}
\toprule
\textbf{Case} & \textbf{Singular} & \textbf{Plural} \\
\midrule
Nominative & cēderbēam & cēderbēamas \\
Genitive & cēderbēames & cēderbēama \\
Dative & cēderbēame & cēderbēamum \\
Accusative & cēderbēam & cēderbēamas \\
\bottomrule
\end{tabular}
\par\endgroup

% mdtable:T02:end
\end{minipage}\par
\reserve{7\baselineskip}
% mdparagraph:P055:begin
\subsection*{horn}
% mdparagraph:P055:end

% teaching: NEW-U17-VOC-HORN-U17-SENSE
\par\begin{minipage}{\linewidth}\raggedright
% mdparagraph:P056:begin
\textbf{horn} — horn; figuratively strength or dignity in 37.14.
% mdparagraph:P056:end

% teaching: NEW-U17-NOM-HORN
% mdparagraph:P057:begin
Masculine; stem \textbf{horn-}. This table selects masculine horn, plural hornas. In Cook 37.14 the horn represents strength or dignity; retain the image when translating.
% mdparagraph:P057:end

% table: NEW-U17-NOM-HORN-DISPLAY
% mdtable:T03:begin
\par\begingroup
\fontsize{10.5}{13.2}\selectfont
\setlength{\LTleft}{0pt}\setlength{\LTright}{\fill}
\noindent
\begin{tabular}{@{}>{\raggedright\arraybackslash}p{39.46mm}>{\raggedright\arraybackslash}p{39.46mm}>{\raggedright\arraybackslash}p{39.46mm}@{}}
\toprule
\textbf{Case} & \textbf{Singular} & \textbf{Plural} \\
\midrule
Nominative & horn & hornas \\
Genitive & hornes & horna \\
Dative & horne & hornum \\
Accusative & horn & hornas \\
\bottomrule
\end{tabular}
\par\endgroup

% mdtable:T03:end
\end{minipage}\par
\reserve{7\baselineskip}
% mdparagraph:P058:begin
\subsection*{lof}
% mdparagraph:P058:end

% teaching: NEW-U17-VOC-LOF-U15-SOURCE-SENSE
\par\begin{minipage}{\linewidth}\raggedright
% mdparagraph:P059:begin
\textbf{lof} — praise.
% mdparagraph:P059:end

% teaching: NEW-U17-NOM-LOF
% mdparagraph:P060:begin
Masculine; stem \textbf{lof-}. Add the endings of the stated pattern to this stem, preserving its vowel length.
% mdparagraph:P060:end

% mdparagraph:P061:begin
This table gives the complete singular for the stated abstract or mass meaning. The familiar noun class is reused; plural meanings and forms are not assigned here.
% mdparagraph:P061:end

% table: NEW-U17-NOM-LOF-DISPLAY
% mdtable:T04:begin
\par\begingroup
\fontsize{10.5}{13.2}\selectfont
\setlength{\LTleft}{0pt}\setlength{\LTright}{\fill}
\noindent
\begin{tabular}{@{}>{\raggedright\arraybackslash}p{60.59mm}>{\raggedright\arraybackslash}p{60.59mm}@{}}
\toprule
\textbf{Case} & \textbf{Singular} \\
\midrule
Nominative & lof \\
Genitive & lofes \\
Dative & lofe \\
Accusative & lof \\
\bottomrule
\end{tabular}
\par\endgroup

% mdtable:T04:end
\end{minipage}\par
\reserve{7\baselineskip}
% mdparagraph:P062:begin
\subsection*{olfend}
% mdparagraph:P062:end

% teaching: NEW-U17-VOC-OLFEND-U17-SENSE
\par\begin{minipage}{\linewidth}\raggedright
% mdparagraph:P063:begin
\textbf{olfend} — camel.
% mdparagraph:P063:end

% teaching: NEW-U17-NOM-OLFEND
% mdparagraph:P064:begin
Masculine; stem \textbf{olfend-}. Retain e in the stem olfend-. The final letters -end do not make this an agent noun of the gyltend pattern.
% mdparagraph:P064:end

% table: NEW-U17-NOM-OLFEND-DISPLAY
% mdtable:T05:begin
\par\begingroup
\fontsize{10.5}{13.2}\selectfont
\setlength{\LTleft}{0pt}\setlength{\LTright}{\fill}
\noindent
\begin{tabular}{@{}>{\raggedright\arraybackslash}p{39.46mm}>{\raggedright\arraybackslash}p{39.46mm}>{\raggedright\arraybackslash}p{39.46mm}@{}}
\toprule
\textbf{Case} & \textbf{Singular} & \textbf{Plural} \\
\midrule
Nominative & olfend & olfendas \\
Genitive & olfendes & olfenda \\
Dative & olfende & olfendum \\
Accusative & olfend & olfendas \\
\bottomrule
\end{tabular}
\par\endgroup

% mdtable:T05:end
\end{minipage}\par
\reserve{7\baselineskip}
% mdparagraph:P065:begin
\subsection*{ðegn}
% mdparagraph:P065:end

% teaching: NEW-U17-VOC-THEGN-U17-SENSE
\par\begin{minipage}{\linewidth}\raggedright
% mdparagraph:P066:begin
\textbf{ðegn} — vassal; retainer; thane.
% mdparagraph:P066:end

% teaching: NEW-U17-NOM-THEGN
% mdparagraph:P067:begin
Masculine; stem \textbf{ðegn-}. Add the endings of the stated pattern to this stem, preserving its vowel length.
% mdparagraph:P067:end

% table: NEW-U17-NOM-THEGN-DISPLAY
% mdtable:T06:begin
\par\begingroup
\fontsize{10.5}{13.2}\selectfont
\setlength{\LTleft}{0pt}\setlength{\LTright}{\fill}
\noindent
\begin{tabular}{@{}>{\raggedright\arraybackslash}p{39.46mm}>{\raggedright\arraybackslash}p{39.46mm}>{\raggedright\arraybackslash}p{39.46mm}@{}}
\toprule
\textbf{Case} & \textbf{Singular} & \textbf{Plural} \\
\midrule
Nominative & ðegn & ðegnas \\
Genitive & ðegnes & ðegna \\
Dative & ðegne & ðegnum \\
Accusative & ðegn & ðegnas \\
\bottomrule
\end{tabular}
\par\endgroup

% mdtable:T06:end
\end{minipage}\par
\reserve{7\baselineskip}
% mdparagraph:P068:begin
\subsection*{wīsdōm}
% mdparagraph:P068:end

% teaching: NEW-U17-VOC-WISDOM-U17-SENSE
\par\begin{minipage}{\linewidth}\raggedright
% mdparagraph:P069:begin
\textbf{wīsdōm} — wisdom.
% mdparagraph:P069:end

% teaching: NEW-U17-NOM-WISDOM
% mdparagraph:P070:begin
Masculine; stem \textbf{wīsdōm-}. Retain both lengths in wīsdōm. The final member follows the familiar masculine dōm pattern; only abstract singular wisdom is used here.
% mdparagraph:P070:end

% table: NEW-U17-NOM-WISDOM-DISPLAY
% mdtable:T07:begin
\par\begingroup
\fontsize{10.5}{13.2}\selectfont
\setlength{\LTleft}{0pt}\setlength{\LTright}{\fill}
\noindent
\begin{tabular}{@{}>{\raggedright\arraybackslash}p{60.59mm}>{\raggedright\arraybackslash}p{60.59mm}@{}}
\toprule
\textbf{Case} & \textbf{Singular} \\
\midrule
Nominative & wīsdōm \\
Genitive & wīsdōmes \\
Dative & wīsdōme \\
Accusative & wīsdōm \\
\bottomrule
\end{tabular}
\par\endgroup

% mdtable:T07:end
\end{minipage}\par
\reserve{7\baselineskip}
% mdparagraph:P071:begin
\subsection*{grundweall}
% mdparagraph:P071:end

% teaching: NEW-U17-VOC-GRUNDWEALL-U17-SENSE
\par\begin{minipage}{\linewidth}\raggedright
% mdparagraph:P072:begin
\textbf{grundweall} — foundation.
% mdparagraph:P072:end

% teaching: NEW-U17-NOM-GRUNDWEALL
% mdparagraph:P073:begin
Masculine; stem \textbf{grundweall-}. Add the endings of the stated pattern to this stem, preserving its vowel length.
% mdparagraph:P073:end

% table: NEW-U17-NOM-GRUNDWEALL-DISPLAY
% mdtable:T08:begin
\par\begingroup
\fontsize{10.5}{13.2}\selectfont
\setlength{\LTleft}{0pt}\setlength{\LTright}{\fill}
\noindent
\begin{tabular}{@{}>{\raggedright\arraybackslash}p{34.82mm}>{\raggedright\arraybackslash}p{41.78mm}>{\raggedright\arraybackslash}p{41.78mm}@{}}
\toprule
\textbf{Case} & \textbf{Singular} & \textbf{Plural} \\
\midrule
Nominative & grundweall & grundweallas \\
Genitive & grundwealles & grundwealla \\
Dative & grundwealle & grundweallum \\
Accusative & grundweall & grundweallas \\
\bottomrule
\end{tabular}
\par\endgroup

% mdtable:T08:end
\end{minipage}\par
\reserve{7\baselineskip}
% mdparagraph:P074:begin
\subsection*{cniht}
% mdparagraph:P074:end

% teaching: NEW-U17-VOC-CNIHT-U17-SERVANT
\par\begin{minipage}{\linewidth}\raggedright
% mdparagraph:P075:begin
\textbf{cniht} — servant.
% mdparagraph:P075:end

% teaching: NEW-U17-NOM-CNIHT
% mdparagraph:P076:begin
Masculine; stem \textbf{cniht-}. The meaning required here is servant. It is distinct from the earlier occurrence meaning youth; familiarity with the word in that passage did not supply this new meaning.
% mdparagraph:P076:end

% table: NEW-U17-NOM-CNIHT-DISPLAY
% mdtable:T09:begin
\par\begingroup
\fontsize{10.5}{13.2}\selectfont
\setlength{\LTleft}{0pt}\setlength{\LTright}{\fill}
\noindent
\begin{tabular}{@{}>{\raggedright\arraybackslash}p{39.46mm}>{\raggedright\arraybackslash}p{39.46mm}>{\raggedright\arraybackslash}p{39.46mm}@{}}
\toprule
\textbf{Case} & \textbf{Singular} & \textbf{Plural} \\
\midrule
Nominative & cniht & cnihtas \\
Genitive & cnihtes & cnihta \\
Dative & cnihte & cnihtum \\
Accusative & cniht & cnihtas \\
\bottomrule
\end{tabular}
\par\endgroup

% mdtable:T09:end
\end{minipage}\par
\reserve{7\baselineskip}
% mdparagraph:P077:begin
\subsection*{ymbhwyrft}
% mdparagraph:P077:end

% teaching: NEW-U17-VOC-YMBHWYRFT-U13-SOURCE
\par\begin{minipage}{\linewidth}\raggedright
% mdparagraph:P078:begin
\textbf{ymbhwyrft} — circuit; compass; orbit.
% mdparagraph:P078:end

% teaching: NEW-U17-NOM-YMBHWYRFT
% mdparagraph:P079:begin
Masculine; stem \textbf{ymbhwyrft-}. Add the endings of the stated pattern to this stem, preserving its vowel length.
% mdparagraph:P079:end

% table: NEW-U17-NOM-YMBHWYRFT-DISPLAY
% mdtable:T10:begin
\par\begingroup
\fontsize{10.5}{13.2}\selectfont
\setlength{\LTleft}{0pt}\setlength{\LTright}{\fill}
\noindent
\begin{tabular}{@{}>{\raggedright\arraybackslash}p{36.99mm}>{\raggedright\arraybackslash}p{40.69mm}>{\raggedright\arraybackslash}p{40.69mm}@{}}
\toprule
\textbf{Case} & \textbf{Singular} & \textbf{Plural} \\
\midrule
Nominative & ymbhwyrft & ymbhwyrftas \\
Genitive & ymbhwyrftes & ymbhwyrfta \\
Dative & ymbhwyrfte & ymbhwyrftum \\
Accusative & ymbhwyrft & ymbhwyrftas \\
\bottomrule
\end{tabular}
\par\endgroup

% mdtable:T10:end
\end{minipage}\par
\reserve{7\baselineskip}
% mdparagraph:P080:begin
\subsection*{gylt}
% mdparagraph:P080:end

% teaching: NEW-U17-VOC-GYLT-U17-SENSE
\par\begin{minipage}{\linewidth}\raggedright
% mdparagraph:P081:begin
\textbf{gylt} — debt, including the prayer’s figurative moral debt.
% mdparagraph:P081:end

% teaching: NEW-U17-NOM-GYLT
% mdparagraph:P082:begin
Masculine; stem \textbf{gylt-}. Retain debt in translating the prayer. Its moral force comes from the prayer’s context; an ordinary noun ending does not by itself express that interpretation.
% mdparagraph:P082:end

% table: NEW-U17-NOM-GYLT-DISPLAY
% mdtable:T11:begin
\par\begingroup
\fontsize{10.5}{13.2}\selectfont
\setlength{\LTleft}{0pt}\setlength{\LTright}{\fill}
\noindent
\begin{tabular}{@{}>{\raggedright\arraybackslash}p{39.46mm}>{\raggedright\arraybackslash}p{39.46mm}>{\raggedright\arraybackslash}p{39.46mm}@{}}
\toprule
\textbf{Case} & \textbf{Singular} & \textbf{Plural} \\
\midrule
Nominative & gylt & gyltas \\
Genitive & gyltes & gylta \\
Dative & gylte & gyltum \\
Accusative & gylt & gyltas \\
\bottomrule
\end{tabular}
\par\endgroup

% mdtable:T11:end
\end{minipage}\par
\reserve{7\baselineskip}
% mdparagraph:P083:begin
\subsection*{gyltend}
% mdparagraph:P083:end

% teaching: NEW-U17-VOC-GYLTEND-U17-SENSE
\par\begin{minipage}{\linewidth}\raggedright
% mdparagraph:P084:begin
\textbf{gyltend} — debtor.
% mdparagraph:P084:end

% teaching: NEW-U17-NOM-GYLTEND
% mdparagraph:P085:begin
Masculine; stem \textbf{gyltend-}. This masculine agent noun has selected nominative/\allowbreak{}accusative plural gyltende and genitive plural gyltendra. An endingless nominative/\allowbreak{}accusative plural is also recorded. The selected pattern is not the ordinary -as plural.
% mdparagraph:P085:end

% table: NEW-U17-NOM-GYLTEND-DISPLAY
% mdtable:T12:begin
\par\begingroup
\fontsize{10.5}{13.2}\selectfont
\setlength{\LTleft}{0pt}\setlength{\LTright}{\fill}
\noindent
\begin{tabular}{@{}>{\raggedright\arraybackslash}p{39.46mm}>{\raggedright\arraybackslash}p{39.46mm}>{\raggedright\arraybackslash}p{39.46mm}@{}}
\toprule
\textbf{Case} & \textbf{Singular} & \textbf{Plural} \\
\midrule
Nominative & gyltend & gyltende \\
Genitive & gyltendes & gyltendra \\
Dative & gyltende & gyltendum \\
Accusative & gyltend & gyltende \\
\bottomrule
\end{tabular}
\par\endgroup

% mdtable:T12:end
\end{minipage}\par
\reserve{7\baselineskip}
% mdparagraph:P086:begin
\subsection*{finger}
% mdparagraph:P086:end

% teaching: NEW-U17-VOC-FINGER-U17-SENSE
\par\begin{minipage}{\linewidth}\raggedright
% mdparagraph:P087:begin
\textbf{finger} — finger.
% mdparagraph:P087:end

% teaching: NEW-U17-NOM-FINGER
% mdparagraph:P088:begin
Masculine; stem \textbf{fingr-}. The second-syllable vowel disappears before an ending: finger, but fingres and fingre. Use the oblique stem fingr-.
% mdparagraph:P088:end

% table: NEW-U17-NOM-FINGER-DISPLAY
% mdtable:T13:begin
\par\begingroup
\fontsize{10.5}{13.2}\selectfont
\setlength{\LTleft}{0pt}\setlength{\LTright}{\fill}
\noindent
\begin{tabular}{@{}>{\raggedright\arraybackslash}p{39.46mm}>{\raggedright\arraybackslash}p{39.46mm}>{\raggedright\arraybackslash}p{39.46mm}@{}}
\toprule
\textbf{Case} & \textbf{Singular} & \textbf{Plural} \\
\midrule
Nominative & finger & fingras \\
Genitive & fingres & fingra \\
Dative & fingre & fingrum \\
Accusative & finger & fingras \\
\bottomrule
\end{tabular}
\par\endgroup

% mdtable:T13:end
\end{minipage}\par
\reserve{7\baselineskip}
% mdparagraph:P089:begin
\subsection*{hearpere}
% mdparagraph:P089:end

% teaching: NEW-U17-VOC-HEARPERE-U17-SENSE
\par\begin{minipage}{\linewidth}\raggedright
% mdparagraph:P090:begin
\textbf{hearpere} — harper.
% mdparagraph:P090:end

% teaching: NEW-U17-NOM-HEARPERE
% mdparagraph:P091:begin
Masculine; stem \textbf{hearper-}. The final -e of the headword gives way to endings on hearper-. The genitive plural is hearpera, as in the familiar masculine -ere class.
% mdparagraph:P091:end

% table: NEW-U17-NOM-HEARPERE-DISPLAY
% mdtable:T14:begin
\par\begingroup
\fontsize{10.5}{13.2}\selectfont
\setlength{\LTleft}{0pt}\setlength{\LTright}{\fill}
\noindent
\begin{tabular}{@{}>{\raggedright\arraybackslash}p{39.46mm}>{\raggedright\arraybackslash}p{39.46mm}>{\raggedright\arraybackslash}p{39.46mm}@{}}
\toprule
\textbf{Case} & \textbf{Singular} & \textbf{Plural} \\
\midrule
Nominative & hearpere & hearperas \\
Genitive & hearperes & hearpera \\
Dative & hearpere & hearperum \\
Accusative & hearpere & hearperas \\
\bottomrule
\end{tabular}
\par\endgroup

% mdtable:T14:end
\end{minipage}\par
\reserve{7\baselineskip}
% mdparagraph:P092:begin
\subsection*{weorðscipe}
% mdparagraph:P092:end

% teaching: NEW-U17-VOC-WEORTHSCIPE-U17-SENSE
\par\begin{minipage}{\linewidth}\raggedright
% mdparagraph:P093:begin
\textbf{weorðscipe} — honour; dignity.
% mdparagraph:P093:end

% teaching: NEW-U17-NOM-WEORTHSCIPE
% mdparagraph:P094:begin
Masculine; stem \textbf{weorðscip-}. This masculine -e noun retains -e in the nominative, dative and accusative singular. The genitive has -es on weorðscip-. Only singular honour or dignity is assigned here.
% mdparagraph:P094:end

% mdparagraph:P095:begin
This table gives the complete singular for the stated abstract or mass meaning. The familiar noun class is reused; plural meanings and forms are not assigned here.
% mdparagraph:P095:end

% table: NEW-U17-NOM-WEORTHSCIPE-DISPLAY
% mdtable:T15:begin
\par\begingroup
\fontsize{10.5}{13.2}\selectfont
\setlength{\LTleft}{0pt}\setlength{\LTright}{\fill}
\noindent
\begin{tabular}{@{}>{\raggedright\arraybackslash}p{57.71mm}>{\raggedright\arraybackslash}p{63.48mm}@{}}
\toprule
\textbf{Case} & \textbf{Singular} \\
\midrule
Nominative & weorðscipe \\
Genitive & weorðscipes \\
Dative & weorðscipe \\
Accusative & weorðscipe \\
\bottomrule
\end{tabular}
\par\endgroup

% mdtable:T15:end
\end{minipage}\par
\reserve{7\baselineskip}
% mdparagraph:P096:begin
\subsection*{fox}
% mdparagraph:P096:end

% teaching: NEW-U17-VOC-FOX-U17-SENSE
\par\begin{minipage}{\linewidth}\raggedright
% mdparagraph:P097:begin
\textbf{fox} — fox.
% mdparagraph:P097:end

% teaching: NEW-U17-NOM-FOX
% mdparagraph:P098:begin
Masculine; stem \textbf{fox-}. Add the endings of the stated pattern to this stem, preserving its vowel length.
% mdparagraph:P098:end

% table: NEW-U17-NOM-FOX-DISPLAY
% mdtable:T16:begin
\par\begingroup
\fontsize{10.5}{13.2}\selectfont
\setlength{\LTleft}{0pt}\setlength{\LTright}{\fill}
\noindent
\begin{tabular}{@{}>{\raggedright\arraybackslash}p{39.46mm}>{\raggedright\arraybackslash}p{39.46mm}>{\raggedright\arraybackslash}p{39.46mm}@{}}
\toprule
\textbf{Case} & \textbf{Singular} & \textbf{Plural} \\
\midrule
Nominative & fox & foxas \\
Genitive & foxes & foxa \\
Dative & foxe & foxum \\
Accusative & fox & foxas \\
\bottomrule
\end{tabular}
\par\endgroup

% mdtable:T16:end
\end{minipage}\par
\reserve{7\baselineskip}
% mdparagraph:P099:begin
\subsection*{ðorn}
% mdparagraph:P099:end

% teaching: NEW-U17-VOC-THORN-U17-SENSE
\par\begin{minipage}{\linewidth}\raggedright
% mdparagraph:P100:begin
\textbf{ðorn} — thorn.
% mdparagraph:P100:end

% teaching: NEW-U17-NOM-THORN
% mdparagraph:P101:begin
Masculine; stem \textbf{ðorn-}. Add the endings of the stated pattern to this stem, preserving its vowel length.
% mdparagraph:P101:end

% table: NEW-U17-NOM-THORN-DISPLAY
% mdtable:T17:begin
\par\begingroup
\fontsize{10.5}{13.2}\selectfont
\setlength{\LTleft}{0pt}\setlength{\LTright}{\fill}
\noindent
\begin{tabular}{@{}>{\raggedright\arraybackslash}p{39.46mm}>{\raggedright\arraybackslash}p{39.46mm}>{\raggedright\arraybackslash}p{39.46mm}@{}}
\toprule
\textbf{Case} & \textbf{Singular} & \textbf{Plural} \\
\midrule
Nominative & ðorn & ðornas \\
Genitive & ðornes & ðorna \\
Dative & ðorne & ðornum \\
Accusative & ðorn & ðornas \\
\bottomrule
\end{tabular}
\par\endgroup

% mdtable:T17:end
\end{minipage}\par
\reserve{7\baselineskip}
% mdparagraph:P102:begin
\subsection*{cynehelm}
% mdparagraph:P102:end

% teaching: NEW-U17-VOC-CYNEHELM-U17-SENSE
\par\begin{minipage}{\linewidth}\raggedright
% mdparagraph:P103:begin
\textbf{cynehelm} — crown.
% mdparagraph:P103:end

% teaching: NEW-U17-NOM-CYNEHELM
% mdparagraph:P104:begin
Masculine; stem \textbf{cynehelm-}. Add the endings of the stated pattern to this stem, preserving its vowel length.
% mdparagraph:P104:end

% table: NEW-U17-NOM-CYNEHELM-DISPLAY
% mdtable:T18:begin
\par\begingroup
\fontsize{10.5}{13.2}\selectfont
\setlength{\LTleft}{0pt}\setlength{\LTright}{\fill}
\noindent
\begin{tabular}{@{}>{\raggedright\arraybackslash}p{39.46mm}>{\raggedright\arraybackslash}p{39.46mm}>{\raggedright\arraybackslash}p{39.46mm}@{}}
\toprule
\textbf{Case} & \textbf{Singular} & \textbf{Plural} \\
\midrule
Nominative & cynehelm & cynehelmas \\
Genitive & cynehelmes & cynehelma \\
Dative & cynehelme & cynehelmum \\
Accusative & cynehelm & cynehelmas \\
\bottomrule
\end{tabular}
\par\endgroup

% mdtable:T18:end
\end{minipage}\par
\reserve{7\baselineskip}
% mdparagraph:P105:begin
\subsection*{gerīmcræft}
% mdparagraph:P105:end

% teaching: NEW-U17-VOC-GERIMCRAEFT-U17-SENSE
\par\begin{minipage}{\linewidth}\raggedright
% mdparagraph:P106:begin
\textbf{gerīmcræft} — arithmetic; computation.
% mdparagraph:P106:end

% teaching: NEW-U17-NOM-GERIMCRAEFT
% mdparagraph:P107:begin
Masculine; stem \textbf{gerīmcræft-}. Add the endings of the stated pattern to this stem, preserving its vowel length.
% mdparagraph:P107:end

% mdparagraph:P108:begin
This table gives the complete singular for the stated abstract or mass meaning. The familiar noun class is reused; plural meanings and forms are not assigned here.
% mdparagraph:P108:end

% table: NEW-U17-NOM-GERIMCRAEFT-DISPLAY
% mdtable:T19:begin
\par\begingroup
\fontsize{10.5}{13.2}\selectfont
\setlength{\LTleft}{0pt}\setlength{\LTright}{\fill}
\noindent
\begin{tabular}{@{}>{\raggedright\arraybackslash}p{55.09mm}>{\raggedright\arraybackslash}p{66.1mm}@{}}
\toprule
\textbf{Case} & \textbf{Singular} \\
\midrule
Nominative & gerīmcræft \\
Genitive & gerīmcræftes \\
Dative & gerīmcræfte \\
Accusative & gerīmcræft \\
\bottomrule
\end{tabular}
\par\endgroup

% mdtable:T19:end
\end{minipage}\par
\reserve{7\baselineskip}
% mdparagraph:P109:begin
\subsection*{stān}
% mdparagraph:P109:end

% teaching: NEW-U17-VOC-STAN-U11-SOURCE
\par\begin{minipage}{\linewidth}\raggedright
% mdparagraph:P110:begin
\textbf{stān} — stone.
% mdparagraph:P110:end

% teaching: NEW-U17-NOM-STAN
% mdparagraph:P111:begin
Masculine; stem \textbf{stān-}. The full stone paradigm is available for the gifts, building materials and bread-or-stone contrast.
% mdparagraph:P111:end

% table: NEW-U17-NOM-STAN-DISPLAY
% mdtable:T20:begin
\par\begingroup
\fontsize{10.5}{13.2}\selectfont
\setlength{\LTleft}{0pt}\setlength{\LTright}{\fill}
\noindent
\begin{tabular}{@{}>{\raggedright\arraybackslash}p{39.46mm}>{\raggedright\arraybackslash}p{39.46mm}>{\raggedright\arraybackslash}p{39.46mm}@{}}
\toprule
\textbf{Case} & \textbf{Singular} & \textbf{Plural} \\
\midrule
Nominative & stān & stānas \\
Genitive & stānes & stāna \\
Dative & stāne & stānum \\
Accusative & stān & stānas \\
\bottomrule
\end{tabular}
\par\endgroup

% mdtable:T20:end
\end{minipage}\par
\reserve{7\baselineskip}
% mdparagraph:P112:begin
\subsection*{Strong neuter nouns}
% mdparagraph:P112:end

% mdparagraph:P113:begin
Nominative and accusative are alike in each number. The singular genitive has -es and dative -e; the plural genitive has -a and dative -um. A light stem normally has nominative/\allowbreak{}accusative plural -u, while a heavy stem normally has no ending. Learn the stated compound and exceptional stems with the table.
% mdparagraph:P113:end

\reserve{7\baselineskip}
% mdparagraph:P114:begin
\subsection*{mægen}
% mdparagraph:P114:end

% teaching: NEW-U17-VOC-MAEGEN-U17-POWER
\par\begin{minipage}{\linewidth}\raggedright
% mdparagraph:P115:begin
\textbf{mægen} — power; strength; in 37.2 heavenly powers.
% mdparagraph:P115:end

% teaching: NEW-U17-NOM-MAEGEN
% mdparagraph:P116:begin
Neuter; stem \textbf{mægen-}. Retain the full stem mægen- in mægenes, mægene, mægena and mægenum. The exceptional nominative/\allowbreak{}accusative plural is mægnu. In the praise catalogue the plural denotes heavenly powers; the singular denotes power or strength.
% mdparagraph:P116:end

% table: NEW-U17-NOM-MAEGEN-DISPLAY
% mdtable:T21:begin
\par\begingroup
\fontsize{10.5}{13.2}\selectfont
\setlength{\LTleft}{0pt}\setlength{\LTright}{\fill}
\noindent
\begin{tabular}{@{}>{\raggedright\arraybackslash}p{39.46mm}>{\raggedright\arraybackslash}p{39.46mm}>{\raggedright\arraybackslash}p{39.46mm}@{}}
\toprule
\textbf{Case} & \textbf{Singular} & \textbf{Plural} \\
\midrule
Nominative & mægen & mægnu \\
Genitive & mægenes & mægena \\
Dative & mægene & mægenum \\
Accusative & mægen & mægnu \\
\bottomrule
\end{tabular}
\par\endgroup

% mdtable:T21:end
\end{minipage}\par
\reserve{7\baselineskip}
% mdparagraph:P117:begin
\subsection*{andgiet}
% mdparagraph:P117:end

% teaching: NEW-U17-VOC-ANDGIET-U15-SOURCE-SENSE
% teaching: NEW-U17-VOC-ANDGIET-U15-SOURCE-UNDERSTANDING
\par\begin{minipage}{\linewidth}\raggedright
% mdparagraph:P118:begin
\textbf{andgiet} — meaning; sense; understanding.
% mdparagraph:P118:end

% teaching: NEW-U17-NOM-ANDGIET
% mdparagraph:P119:begin
Neuter; stem \textbf{andgiet-}. Learn the four singular cases for meaning and understanding. These are abstract uses of a familiar neuter class. Plural forms and separate count meanings are outside this unit; the absence of a plural column does not mean that the word has no plural.
% mdparagraph:P119:end

% table: NEW-U17-NOM-ANDGIET-DISPLAY
% mdtable:T22:begin
\par\begingroup
\fontsize{10.5}{13.2}\selectfont
\setlength{\LTleft}{0pt}\setlength{\LTright}{\fill}
\noindent
\begin{tabular}{@{}>{\raggedright\arraybackslash}p{60.59mm}>{\raggedright\arraybackslash}p{60.59mm}@{}}
\toprule
\textbf{Case} & \textbf{Singular} \\
\midrule
Nominative & andgiet \\
Genitive & andgietes \\
Dative & andgiete \\
Accusative & andgiet \\
\bottomrule
\end{tabular}
\par\endgroup

% mdtable:T22:end
\end{minipage}\par
\reserve{7\baselineskip}
% mdparagraph:P120:begin
\subsection*{pund}
% mdparagraph:P120:end

% teaching: NEW-U17-VOC-PUND-U17-SENSE
\par\begin{minipage}{\linewidth}\raggedright
% mdparagraph:P121:begin
\textbf{pund} — pound; source weight measure.
% mdparagraph:P121:end

% teaching: NEW-U17-NOM-PUND
% mdparagraph:P122:begin
Neuter; stem \textbf{pund-}. Add the endings of the stated pattern to this stem, preserving its vowel length.
% mdparagraph:P122:end

% table: NEW-U17-NOM-PUND-DISPLAY
% mdtable:T23:begin
\par\begingroup
\fontsize{10.5}{13.2}\selectfont
\setlength{\LTleft}{0pt}\setlength{\LTright}{\fill}
\noindent
\begin{tabular}{@{}>{\raggedright\arraybackslash}p{39.46mm}>{\raggedright\arraybackslash}p{39.46mm}>{\raggedright\arraybackslash}p{39.46mm}@{}}
\toprule
\textbf{Case} & \textbf{Singular} & \textbf{Plural} \\
\midrule
Nominative & pund & pund \\
Genitive & pundes & punda \\
Dative & punde & pundum \\
Accusative & pund & pund \\
\bottomrule
\end{tabular}
\par\endgroup

% mdtable:T23:end
\end{minipage}\par
\reserve{7\baselineskip}
% mdparagraph:P123:begin
\subsection*{edlēan}
% mdparagraph:P123:end

% teaching: NEW-U17-VOC-EDLEAN-U17-SENSE
\par\begin{minipage}{\linewidth}\raggedright
% mdparagraph:P124:begin
\textbf{edlēan} — reward; recompense.
% mdparagraph:P124:end

% teaching: NEW-U17-NOM-EDLEAN
% mdparagraph:P125:begin
Neuter; stem \textbf{edlēan-}. Add the endings of the stated pattern to this stem, preserving its vowel length.
% mdparagraph:P125:end

% table: NEW-U17-NOM-EDLEAN-DISPLAY
% mdtable:T24:begin
\par\begingroup
\fontsize{10.5}{13.2}\selectfont
\setlength{\LTleft}{0pt}\setlength{\LTright}{\fill}
\noindent
\begin{tabular}{@{}>{\raggedright\arraybackslash}p{39.46mm}>{\raggedright\arraybackslash}p{39.46mm}>{\raggedright\arraybackslash}p{39.46mm}@{}}
\toprule
\textbf{Case} & \textbf{Singular} & \textbf{Plural} \\
\midrule
Nominative & edlēan & edlēan \\
Genitive & edlēanes & edlēana \\
Dative & edlēane & edlēanum \\
Accusative & edlēan & edlēan \\
\bottomrule
\end{tabular}
\par\endgroup

% mdtable:T24:end
\end{minipage}\par
\reserve{7\baselineskip}
% mdparagraph:P126:begin
\subsection*{ðrymsetl}
% mdparagraph:P126:end

% teaching: NEW-U17-VOC-THRYMSETL-U17-SENSE
\par\begin{minipage}{\linewidth}\raggedright
% mdparagraph:P127:begin
\textbf{ðrymsetl} — throne; throne of glory.
% mdparagraph:P127:end

% teaching: NEW-U17-NOM-THRYMSETL
% mdparagraph:P128:begin
Neuter; stem \textbf{ðrymsetl-}. Retain the source spelling ðrymsetl. The final stem setl is heavy, so the nominative/\allowbreak{}accusative plural has no -u.
% mdparagraph:P128:end

% table: NEW-U17-NOM-THRYMSETL-DISPLAY
% mdtable:T25:begin
\par\begingroup
\fontsize{10.5}{13.2}\selectfont
\setlength{\LTleft}{0pt}\setlength{\LTright}{\fill}
\noindent
\begin{tabular}{@{}>{\raggedright\arraybackslash}p{39.46mm}>{\raggedright\arraybackslash}p{39.46mm}>{\raggedright\arraybackslash}p{39.46mm}@{}}
\toprule
\textbf{Case} & \textbf{Singular} & \textbf{Plural} \\
\midrule
Nominative & ðrymsetl & ðrymsetl \\
Genitive & ðrymsetles & ðrymsetla \\
Dative & ðrymsetle & ðrymsetlum \\
Accusative & ðrymsetl & ðrymsetl \\
\bottomrule
\end{tabular}
\par\endgroup

% mdtable:T25:end
\end{minipage}\par
\reserve{7\baselineskip}
% mdparagraph:P129:begin
\subsection*{palmtwīg}
% mdparagraph:P129:end

% teaching: NEW-U17-VOC-PALMTWIG-U17-SENSE
\par\begin{minipage}{\linewidth}\raggedright
% mdparagraph:P130:begin
\textbf{palmtwīg} — palm twig; palm branch.
% mdparagraph:P130:end

% teaching: NEW-U17-NOM-PALMTWIG
% mdparagraph:P131:begin
Neuter; stem \textbf{palmtwīg-}. Long ī makes palmtwīg a heavy neuter. The nominative/\allowbreak{}accusative plural is palmtwīg.
% mdparagraph:P131:end

% table: NEW-U17-NOM-PALMTWIG-DISPLAY
% mdtable:T26:begin
\par\begingroup
\fontsize{10.5}{13.2}\selectfont
\setlength{\LTleft}{0pt}\setlength{\LTright}{\fill}
\noindent
\begin{tabular}{@{}>{\raggedright\arraybackslash}p{39.46mm}>{\raggedright\arraybackslash}p{39.46mm}>{\raggedright\arraybackslash}p{39.46mm}@{}}
\toprule
\textbf{Case} & \textbf{Singular} & \textbf{Plural} \\
\midrule
Nominative & palmtwīg & palmtwīg \\
Genitive & palmtwīges & palmtwīga \\
Dative & palmtwīge & palmtwīgum \\
Accusative & palmtwīg & palmtwīg \\
\bottomrule
\end{tabular}
\par\endgroup

% mdtable:T26:end
\end{minipage}\par
\reserve{7\baselineskip}
% mdparagraph:P132:begin
\subsection*{cnēo}
% mdparagraph:P132:end

% teaching: NEW-U17-VOC-CNEO-U17-KNEE
\par\begin{minipage}{\linewidth}\raggedright
% mdparagraph:P133:begin
\textbf{cnēo} — knee.
% mdparagraph:P133:end

% teaching: NEW-U17-NOM-CNEO
% mdparagraph:P134:begin
Neuter; stem \textbf{cnēow-}. The endingless singular cnēo has cnēow- before endings. Its nominative/\allowbreak{}accusative plural is cnēowu.
% mdparagraph:P134:end

% table: NEW-U17-NOM-CNEO-DISPLAY
% mdtable:T27:begin
\par\begingroup
\fontsize{10.5}{13.2}\selectfont
\setlength{\LTleft}{0pt}\setlength{\LTright}{\fill}
\noindent
\begin{tabular}{@{}>{\raggedright\arraybackslash}p{39.46mm}>{\raggedright\arraybackslash}p{39.46mm}>{\raggedright\arraybackslash}p{39.46mm}@{}}
\toprule
\textbf{Case} & \textbf{Singular} & \textbf{Plural} \\
\midrule
Nominative & cnēo & cnēowu \\
Genitive & cnēowes & cnēowa \\
Dative & cnēowe & cnēowum \\
Accusative & cnēo & cnēowu \\
\bottomrule
\end{tabular}
\par\endgroup

% mdtable:T27:end
\end{minipage}\par
\reserve{7\baselineskip}
% mdparagraph:P135:begin
\subsection*{hol}
% mdparagraph:P135:end

% teaching: NEW-U17-VOC-HOL-U17-SENSE
\par\begin{minipage}{\linewidth}\raggedright
% mdparagraph:P136:begin
\textbf{hol} — hole.
% mdparagraph:P136:end

% teaching: NEW-U17-NOM-HOL
% mdparagraph:P137:begin
Neuter; stem \textbf{hol-}. The short vowel and single final consonant give the light neuter pattern: nominative/\allowbreak{}accusative plural holu.
% mdparagraph:P137:end

% table: NEW-U17-NOM-HOL-DISPLAY
% mdtable:T28:begin
\par\begingroup
\fontsize{10.5}{13.2}\selectfont
\setlength{\LTleft}{0pt}\setlength{\LTright}{\fill}
\noindent
\begin{tabular}{@{}>{\raggedright\arraybackslash}p{39.46mm}>{\raggedright\arraybackslash}p{39.46mm}>{\raggedright\arraybackslash}p{39.46mm}@{}}
\toprule
\textbf{Case} & \textbf{Singular} & \textbf{Plural} \\
\midrule
Nominative & hol & holu \\
Genitive & holes & hola \\
Dative & hole & holum \\
Accusative & hol & holu \\
\bottomrule
\end{tabular}
\par\endgroup

% mdtable:T28:end
\end{minipage}\par
\reserve{7\baselineskip}
% mdparagraph:P138:begin
\subsection*{nest}
% mdparagraph:P138:end

% teaching: NEW-U17-VOC-NEST-U17-SENSE
\par\begin{minipage}{\linewidth}\raggedright
% mdparagraph:P139:begin
\textbf{nest} — nest.
% mdparagraph:P139:end

% teaching: NEW-U17-NOM-NEST
% mdparagraph:P140:begin
Neuter; stem \textbf{nest-}. The final consonant cluster makes the stem heavy. The nominative/\allowbreak{}accusative plural is nest.
% mdparagraph:P140:end

% table: NEW-U17-NOM-NEST-DISPLAY
% mdtable:T29:begin
\par\begingroup
\fontsize{10.5}{13.2}\selectfont
\setlength{\LTleft}{0pt}\setlength{\LTright}{\fill}
\noindent
\begin{tabular}{@{}>{\raggedright\arraybackslash}p{39.46mm}>{\raggedright\arraybackslash}p{39.46mm}>{\raggedright\arraybackslash}p{39.46mm}@{}}
\toprule
\textbf{Case} & \textbf{Singular} & \textbf{Plural} \\
\midrule
Nominative & nest & nest \\
Genitive & nestes & nesta \\
Dative & neste & nestum \\
Accusative & nest & nest \\
\bottomrule
\end{tabular}
\par\endgroup

% mdtable:T29:end
\end{minipage}\par
\reserve{7\baselineskip}
% mdparagraph:P141:begin
\subsection*{hrēod}
% mdparagraph:P141:end

% teaching: NEW-U17-VOC-HREOD-U17-SENSE
\par\begin{minipage}{\linewidth}\raggedright
% mdparagraph:P142:begin
\textbf{hrēod} — reed.
% mdparagraph:P142:end

% teaching: NEW-U17-NOM-HREOD
% mdparagraph:P143:begin
Neuter; stem \textbf{hrēod-}. Long ēo makes hrēod a heavy neuter. Its nominative/\allowbreak{}accusative plural is endingless.
% mdparagraph:P143:end

% table: NEW-U17-NOM-HREOD-DISPLAY
% mdtable:T30:begin
\par\begingroup
\fontsize{10.5}{13.2}\selectfont
\setlength{\LTleft}{0pt}\setlength{\LTright}{\fill}
\noindent
\begin{tabular}{@{}>{\raggedright\arraybackslash}p{39.46mm}>{\raggedright\arraybackslash}p{39.46mm}>{\raggedright\arraybackslash}p{39.46mm}@{}}
\toprule
\textbf{Case} & \textbf{Singular} & \textbf{Plural} \\
\midrule
Nominative & hrēod & hrēod \\
Genitive & hrēodes & hrēoda \\
Dative & hrēode & hrēodum \\
Accusative & hrēod & hrēod \\
\bottomrule
\end{tabular}
\par\endgroup

% mdtable:T30:end
\end{minipage}\par
\reserve{7\baselineskip}
% mdparagraph:P144:begin
\subsection*{Strong feminine nouns}
% mdparagraph:P144:end

% mdparagraph:P145:begin
The selected pattern has nominative singular without an added ending and genitive, dative and accusative singular -e. Plural nominative, genitive and accusative have -a; dative has -um. The endingless accusative of cwēn and the doubled consonants of the -nes and related stems require particular attention.
% mdparagraph:P145:end

\reserve{7\baselineskip}
% mdparagraph:P146:begin
\subsection*{hēanes}
% mdparagraph:P146:end

% teaching: NEW-U17-VOC-HEANES-HEIGHTS-U13-SOURCE
\par\begin{minipage}{\linewidth}\raggedright
% mdparagraph:P147:begin
\textbf{hēanes} — heights; highest places.
% mdparagraph:P147:end

% teaching: NEW-U17-VOC-HEANES-U13-SOURCE
% mdparagraph:P148:begin
\textbf{hēanes} — height; highest point. Supplied reference only.
% mdparagraph:P148:end

% teaching: NEW-U17-NOM-HEANES
% mdparagraph:P149:begin
Feminine; stem \textbf{hēaness-}. Use the plural for heights or the highest places. The singular column gives reference forms only; singular height is not required from memory here. The stem is hēaness-, with doubled s before every ending and with e retained.
% mdparagraph:P149:end

% table: NEW-U17-NOM-HEANES-DISPLAY
% mdtable:T31:begin
\par\begingroup
\fontsize{10.5}{13.2}\selectfont
\setlength{\LTleft}{0pt}\setlength{\LTright}{\fill}
\noindent
\begin{tabular}{@{}>{\raggedright\arraybackslash}p{39.46mm}>{\raggedright\arraybackslash}p{39.46mm}>{\raggedright\arraybackslash}p{39.46mm}@{}}
\toprule
\textbf{Case} & \textbf{Singular} & \textbf{Plural} \\
\midrule
Nominative & hēanes & hēanessa \\
Genitive & hēanesse & hēanessa \\
Dative & hēanesse & hēanessum \\
Accusative & hēanesse & hēanessa \\
\bottomrule
\end{tabular}
\par\endgroup

% mdtable:T31:end
\end{minipage}\par
\reserve{7\baselineskip}
% mdparagraph:P150:begin
\subsection*{ȳst}
% mdparagraph:P150:end

% teaching: NEW-U17-VOC-YST-U17-SENSE
\par\begin{minipage}{\linewidth}\raggedright
% mdparagraph:P151:begin
\textbf{ȳst} — storm.
% mdparagraph:P151:end

% teaching: NEW-U17-NOM-YST
% mdparagraph:P152:begin
Feminine; stem \textbf{ȳst-}. Add the endings of the stated pattern to this stem, preserving its vowel length.
% mdparagraph:P152:end

% table: NEW-U17-NOM-YST-DISPLAY
% mdtable:T32:begin
\par\begingroup
\fontsize{10.5}{13.2}\selectfont
\setlength{\LTleft}{0pt}\setlength{\LTright}{\fill}
\noindent
\begin{tabular}{@{}>{\raggedright\arraybackslash}p{39.46mm}>{\raggedright\arraybackslash}p{39.46mm}>{\raggedright\arraybackslash}p{39.46mm}@{}}
\toprule
\textbf{Case} & \textbf{Singular} & \textbf{Plural} \\
\midrule
Nominative & ȳst & ȳsta \\
Genitive & ȳste & ȳsta \\
Dative & ȳste & ȳstum \\
Accusative & ȳste & ȳsta \\
\bottomrule
\end{tabular}
\par\endgroup

% mdtable:T32:end
\end{minipage}\par
\reserve{7\baselineskip}
% mdparagraph:P153:begin
\subsection*{cwēn}
% mdparagraph:P153:end

% teaching: NEW-U17-VOC-CWEN-U17-QUEEN
% teaching: NEW-U17-VOC-CWEN-U17-PRINCESS
\par\begin{minipage}{\linewidth}\raggedright
% mdparagraph:P154:begin
\textbf{cwēn} — queen; princess.
% mdparagraph:P154:end

% teaching: NEW-U17-NOM-CWEN
% mdparagraph:P155:begin
Feminine; stem \textbf{cwēn-}. The selected accusative singular is endingless cwēn; the genitive and dative have cwēne. The ordinary-class accusative cwēne is also a defensible variant, but the contrast exercise explicitly asks for cwēn. Queen and princess are separately assigned contextual meanings.
% mdparagraph:P155:end

% table: NEW-U17-NOM-CWEN-DISPLAY
% mdtable:T33:begin
\par\begingroup
\fontsize{10.5}{13.2}\selectfont
\setlength{\LTleft}{0pt}\setlength{\LTright}{\fill}
\noindent
\begin{tabular}{@{}>{\raggedright\arraybackslash}p{39.46mm}>{\raggedright\arraybackslash}p{39.46mm}>{\raggedright\arraybackslash}p{39.46mm}@{}}
\toprule
\textbf{Case} & \textbf{Singular} & \textbf{Plural} \\
\midrule
Nominative & cwēn & cwēna \\
Genitive & cwēne & cwēna \\
Dative & cwēne & cwēnum \\
Accusative & cwēn & cwēna \\
\bottomrule
\end{tabular}
\par\endgroup

% mdtable:T33:end
\end{minipage}\par
\reserve{7\baselineskip}
% mdparagraph:P156:begin
\subsection*{ðēnung}
% mdparagraph:P156:end

% teaching: NEW-U17-VOC-THENUNG-U16-READING
\par\begin{minipage}{\linewidth}\raggedright
% mdparagraph:P157:begin
\textbf{ðēnung} — service; attendance; service performed.
% mdparagraph:P157:end

% teaching: NEW-U17-NOM-THENUNG
% mdparagraph:P158:begin
Feminine; stem \textbf{ðēnung-}. The selected singular oblique ending is -e. Nouns in -ung can also have -a in the dative, and sometimes in the genitive and accusative; retain such a form when it is given in a source.
% mdparagraph:P158:end

% table: NEW-U17-NOM-THENUNG-DISPLAY
% mdtable:T34:begin
\par\begingroup
\fontsize{10.5}{13.2}\selectfont
\setlength{\LTleft}{0pt}\setlength{\LTright}{\fill}
\noindent
\begin{tabular}{@{}>{\raggedright\arraybackslash}p{39.46mm}>{\raggedright\arraybackslash}p{39.46mm}>{\raggedright\arraybackslash}p{39.46mm}@{}}
\toprule
\textbf{Case} & \textbf{Singular} & \textbf{Plural} \\
\midrule
Nominative & ðēnung & ðēnunga \\
Genitive & ðēnunge & ðēnunga \\
Dative & ðēnunge & ðēnungum \\
Accusative & ðēnunge & ðēnunga \\
\bottomrule
\end{tabular}
\par\endgroup

% mdtable:T34:end
\end{minipage}\par
\reserve{7\baselineskip}
% mdparagraph:P159:begin
\subsection*{getimbrung}
% mdparagraph:P159:end

% teaching: NEW-U17-VOC-GETIMBRUNG-U17-SENSE
\par\begin{minipage}{\linewidth}\raggedright
% mdparagraph:P160:begin
\textbf{getimbrung} — building; structure built.
% mdparagraph:P160:end

% teaching: NEW-U17-NOM-GETIMBRUNG
% mdparagraph:P161:begin
Feminine; stem \textbf{getimbrung-}. The noun is feminine. A pronoun resuming the building must agree with that grammatical gender.
% mdparagraph:P161:end

% table: NEW-U17-NOM-GETIMBRUNG-DISPLAY
% mdtable:T35:begin
\par\begingroup
\fontsize{10.5}{13.2}\selectfont
\setlength{\LTleft}{0pt}\setlength{\LTright}{\fill}
\noindent
\begin{tabular}{@{}>{\raggedright\arraybackslash}p{35.87mm}>{\raggedright\arraybackslash}p{39.46mm}>{\raggedright\arraybackslash}p{43.05mm}@{}}
\toprule
\textbf{Case} & \textbf{Singular} & \textbf{Plural} \\
\midrule
Nominative & getimbrung & getimbrunga \\
Genitive & getimbrunge & getimbrunga \\
Dative & getimbrunge & getimbrungum \\
Accusative & getimbrunge & getimbrunga \\
\bottomrule
\end{tabular}
\par\endgroup

% mdtable:T35:end
\end{minipage}\par
\reserve{7\baselineskip}
% mdparagraph:P162:begin
\subsection*{ēaðmōdnes}
% mdparagraph:P162:end

% teaching: NEW-U17-VOC-EADMODNES-U17-SENSE
\par\begin{minipage}{\linewidth}\raggedright
% mdparagraph:P163:begin
\textbf{ēaðmōdnes} — humility.
% mdparagraph:P163:end

% teaching: NEW-U17-NOM-EADMODNES
% mdparagraph:P164:begin
Feminine; stem \textbf{ēaðmōdness-}. Use ēaðmōdness- before the endings, with doubled s and retained e. Only abstract singular humility is assigned.
% mdparagraph:P164:end

% table: NEW-U17-NOM-EADMODNES-DISPLAY
% mdtable:T36:begin
\par\begingroup
\fontsize{10.5}{13.2}\selectfont
\setlength{\LTleft}{0pt}\setlength{\LTright}{\fill}
\noindent
\begin{tabular}{@{}>{\raggedright\arraybackslash}p{57.71mm}>{\raggedright\arraybackslash}p{63.48mm}@{}}
\toprule
\textbf{Case} & \textbf{Singular} \\
\midrule
Nominative & ēaðmōdnes \\
Genitive & ēaðmōdnesse \\
Dative & ēaðmōdnesse \\
Accusative & ēaðmōdnesse \\
\bottomrule
\end{tabular}
\par\endgroup

% mdtable:T36:end
\end{minipage}\par
\reserve{7\baselineskip}
% mdparagraph:P165:begin
\subsection*{ðīnen}
% mdparagraph:P165:end

% teaching: NEW-U17-VOC-THINEN-U17-SENSE
\par\begin{minipage}{\linewidth}\raggedright
% mdparagraph:P166:begin
\textbf{ðīnen} — handmaiden.
% mdparagraph:P166:end

% teaching: NEW-U17-NOM-THINEN
% mdparagraph:P167:begin
Feminine; stem \textbf{ðīnenn-}. The oblique stem is ðīnenn-, with doubled n and retained e: ðīnenne. The nominative singular conceals the doubling.
% mdparagraph:P167:end

% table: NEW-U17-NOM-THINEN-DISPLAY
% mdtable:T37:begin
\par\begingroup
\fontsize{10.5}{13.2}\selectfont
\setlength{\LTleft}{0pt}\setlength{\LTright}{\fill}
\noindent
\begin{tabular}{@{}>{\raggedright\arraybackslash}p{39.46mm}>{\raggedright\arraybackslash}p{39.46mm}>{\raggedright\arraybackslash}p{39.46mm}@{}}
\toprule
\textbf{Case} & \textbf{Singular} & \textbf{Plural} \\
\midrule
Nominative & ðīnen & ðīnenna \\
Genitive & ðīnenne & ðīnenna \\
Dative & ðīnenne & ðīnennum \\
Accusative & ðīnenne & ðīnenna \\
\bottomrule
\end{tabular}
\par\endgroup

% mdtable:T37:end
\end{minipage}\par
\reserve{7\baselineskip}
% mdparagraph:P168:begin
\subsection*{cnēoris}
% mdparagraph:P168:end

% teaching: NEW-U17-VOC-CNEORIS-U17-SENSE
\par\begin{minipage}{\linewidth}\raggedright
% mdparagraph:P169:begin
\textbf{cnēoris} — generation.
% mdparagraph:P169:end

% teaching: NEW-U17-NOM-CNEORIS
% mdparagraph:P170:begin
Feminine; stem \textbf{cnēoriss-}. Use cnēoriss- before endings: cnēorisse, cnēorissa and cnēorissum. The assigned meaning is generation.
% mdparagraph:P170:end

% table: NEW-U17-NOM-CNEORIS-DISPLAY
% mdtable:T38:begin
\par\begingroup
\fontsize{10.5}{13.2}\selectfont
\setlength{\LTleft}{0pt}\setlength{\LTright}{\fill}
\noindent
\begin{tabular}{@{}>{\raggedright\arraybackslash}p{39.46mm}>{\raggedright\arraybackslash}p{39.46mm}>{\raggedright\arraybackslash}p{39.46mm}@{}}
\toprule
\textbf{Case} & \textbf{Singular} & \textbf{Plural} \\
\midrule
Nominative & cnēoris & cnēorissa \\
Genitive & cnēorisse & cnēorissa \\
Dative & cnēorisse & cnēorissum \\
Accusative & cnēorisse & cnēorissa \\
\bottomrule
\end{tabular}
\par\endgroup

% mdtable:T38:end
\end{minipage}\par
\reserve{7\baselineskip}
% mdparagraph:P171:begin
\subsection*{mǣgð}
% mdparagraph:P171:end

% teaching: NEW-U17-VOC-MAEGTH-TRIBE-U17-SENSE
\par\begin{minipage}{\linewidth}\raggedright
% mdparagraph:P172:begin
\textbf{mǣgð} — tribe; nation.
% mdparagraph:P172:end

% teaching: NEW-U17-NOM-MAEGTH-TRIBE
% mdparagraph:P173:begin
Feminine; stem \textbf{mǣgð-}. Long ǣ distinguishes mǣgð, tribe or nation, from short-vowel mægð, maiden. Keep both the length and this noun’s feminine pattern.
% mdparagraph:P173:end

% table: NEW-U17-NOM-MAEGTH-TRIBE-DISPLAY
% mdtable:T39:begin
\par\begingroup
\fontsize{10.5}{13.2}\selectfont
\setlength{\LTleft}{0pt}\setlength{\LTright}{\fill}
\noindent
\begin{tabular}{@{}>{\raggedright\arraybackslash}p{39.46mm}>{\raggedright\arraybackslash}p{39.46mm}>{\raggedright\arraybackslash}p{39.46mm}@{}}
\toprule
\textbf{Case} & \textbf{Singular} & \textbf{Plural} \\
\midrule
Nominative & mǣgð & mǣgða \\
Genitive & mǣgðe & mǣgða \\
Dative & mǣgðe & mǣgðum \\
Accusative & mǣgðe & mǣgða \\
\bottomrule
\end{tabular}
\par\endgroup

% mdtable:T39:end
\end{minipage}\par
\reserve{7\baselineskip}
% mdparagraph:P174:begin
\subsection*{blētsung}
% mdparagraph:P174:end

% teaching: NEW-U17-VOC-BLETSUNG-U17-SENSE
\par\begin{minipage}{\linewidth}\raggedright
% mdparagraph:P175:begin
\textbf{blētsung} — blessing.
% mdparagraph:P175:end

% teaching: NEW-U17-NOM-BLETSUNG
% mdparagraph:P176:begin
Feminine; stem \textbf{blētsung-}. Add the endings of the stated pattern to this stem, preserving its vowel length.
% mdparagraph:P176:end

% mdparagraph:P177:begin
This table gives the complete singular for the stated abstract or mass meaning. The familiar noun class is reused; plural meanings and forms are not assigned here.
% mdparagraph:P177:end

% table: NEW-U17-NOM-BLETSUNG-DISPLAY
% mdtable:T40:begin
\par\begingroup
\fontsize{10.5}{13.2}\selectfont
\setlength{\LTleft}{0pt}\setlength{\LTright}{\fill}
\noindent
\begin{tabular}{@{}>{\raggedright\arraybackslash}p{60.59mm}>{\raggedright\arraybackslash}p{60.59mm}@{}}
\toprule
\textbf{Case} & \textbf{Singular} \\
\midrule
Nominative & blētsung \\
Genitive & blētsunge \\
Dative & blētsunge \\
Accusative & blētsunge \\
\bottomrule
\end{tabular}
\par\endgroup

% mdtable:T40:end
\end{minipage}\par
\reserve{7\baselineskip}
% mdparagraph:P178:begin
\subsection*{beorhtnes}
% mdparagraph:P178:end

% teaching: NEW-U17-VOC-BEORHTNES-U17-SENSE
\par\begin{minipage}{\linewidth}\raggedright
% mdparagraph:P179:begin
\textbf{beorhtnes} — brightness.
% mdparagraph:P179:end

% teaching: NEW-U17-NOM-BEORHTNES
% mdparagraph:P180:begin
Feminine; stem \textbf{beorhtness-}. Add the endings of the stated pattern to this stem, preserving its vowel length.
% mdparagraph:P180:end

% mdparagraph:P181:begin
This table gives the complete singular for the stated abstract or mass meaning. The familiar noun class is reused; plural meanings and forms are not assigned here.
% mdparagraph:P181:end

% table: NEW-U17-NOM-BEORHTNES-DISPLAY
% mdtable:T41:begin
\par\begingroup
\fontsize{10.5}{13.2}\selectfont
\setlength{\LTleft}{0pt}\setlength{\LTright}{\fill}
\noindent
\begin{tabular}{@{}>{\raggedright\arraybackslash}p{57.71mm}>{\raggedright\arraybackslash}p{63.48mm}@{}}
\toprule
\textbf{Case} & \textbf{Singular} \\
\midrule
Nominative & beorhtnes \\
Genitive & beorhtnesse \\
Dative & beorhtnesse \\
Accusative & beorhtnesse \\
\bottomrule
\end{tabular}
\par\endgroup

% mdtable:T41:end
\end{minipage}\par
\reserve{7\baselineskip}
% mdparagraph:P182:begin
\subsection*{ðancung}
% mdparagraph:P182:end

% teaching: NEW-U17-VOC-THANCUNG-U17-SENSE
\par\begin{minipage}{\linewidth}\raggedright
% mdparagraph:P183:begin
\textbf{ðancung} — thanksgiving.
% mdparagraph:P183:end

% teaching: NEW-U17-NOM-THANCUNG
% mdparagraph:P184:begin
Feminine; stem \textbf{ðancung-}. Add the endings of the stated pattern to this stem, preserving its vowel length.
% mdparagraph:P184:end

% mdparagraph:P185:begin
This table gives the complete singular for the stated abstract or mass meaning. The familiar noun class is reused; plural meanings and forms are not assigned here.
% mdparagraph:P185:end

% table: NEW-U17-NOM-THANCUNG-DISPLAY
% mdtable:T42:begin
\par\begingroup
\fontsize{10.5}{13.2}\selectfont
\setlength{\LTleft}{0pt}\setlength{\LTright}{\fill}
\noindent
\begin{tabular}{@{}>{\raggedright\arraybackslash}p{60.59mm}>{\raggedright\arraybackslash}p{60.59mm}@{}}
\toprule
\textbf{Case} & \textbf{Singular} \\
\midrule
Nominative & ðancung \\
Genitive & ðancunge \\
Dative & ðancunge \\
Accusative & ðancunge \\
\bottomrule
\end{tabular}
\par\endgroup

% mdtable:T42:end
\end{minipage}\par
\reserve{7\baselineskip}
% mdparagraph:P186:begin
\subsection*{Weak nouns}
% mdparagraph:P186:end

% mdparagraph:P187:begin
Weak masculines normally have nominative singular -a, other singular cases -an, plural nominative/\allowbreak{}accusative -an, genitive -ena and dative -um. Weak feminines have nominative singular -e and otherwise the same endings. The weak neuter ēare keeps -e also in the accusative singular and has a contracted genitive plural.
% mdparagraph:P187:end

\reserve{7\baselineskip}
% mdparagraph:P188:begin
\subsection*{draca}
% mdparagraph:P188:end

% teaching: NEW-U17-VOC-DRACA-U17-SENSE
\par\begin{minipage}{\linewidth}\raggedright
% mdparagraph:P189:begin
\textbf{draca} — dragon.
% mdparagraph:P189:end

% teaching: NEW-U17-NOM-DRACA
% mdparagraph:P190:begin
Masculine; stem \textbf{drac-}. Add the endings of the stated pattern to this stem, preserving its vowel length.
% mdparagraph:P190:end

% table: NEW-U17-NOM-DRACA-DISPLAY
% mdtable:T43:begin
\par\begingroup
\fontsize{10.5}{13.2}\selectfont
\setlength{\LTleft}{0pt}\setlength{\LTright}{\fill}
\noindent
\begin{tabular}{@{}>{\raggedright\arraybackslash}p{39.46mm}>{\raggedright\arraybackslash}p{39.46mm}>{\raggedright\arraybackslash}p{39.46mm}@{}}
\toprule
\textbf{Case} & \textbf{Singular} & \textbf{Plural} \\
\midrule
Nominative & draca & dracan \\
Genitive & dracan & dracena \\
Dative & dracan & dracum \\
Accusative & dracan & dracan \\
\bottomrule
\end{tabular}
\par\endgroup

% mdtable:T43:end
\end{minipage}\par
\reserve{7\baselineskip}
% mdparagraph:P191:begin
\subsection*{nædre}
% mdparagraph:P191:end

% teaching: NEW-U17-VOC-NAEDRE-U17-SENSE
\par\begin{minipage}{\linewidth}\raggedright
% mdparagraph:P192:begin
\textbf{nædre} — serpent.
% mdparagraph:P192:end

% teaching: NEW-U17-NOM-NAEDRE
% mdparagraph:P193:begin
Feminine; stem \textbf{nædr-}. The selected spelling is nædre, with short æ and one d. The source also gives næddre. Add the weak feminine endings to nædr-.
% mdparagraph:P193:end

% table: NEW-U17-NOM-NAEDRE-DISPLAY
% mdtable:T44:begin
\par\begingroup
\fontsize{10.5}{13.2}\selectfont
\setlength{\LTleft}{0pt}\setlength{\LTright}{\fill}
\noindent
\begin{tabular}{@{}>{\raggedright\arraybackslash}p{39.46mm}>{\raggedright\arraybackslash}p{39.46mm}>{\raggedright\arraybackslash}p{39.46mm}@{}}
\toprule
\textbf{Case} & \textbf{Singular} & \textbf{Plural} \\
\midrule
Nominative & nædre & nædran \\
Genitive & nædran & nædrena \\
Dative & nædran & nædrum \\
Accusative & nædran & nædran \\
\bottomrule
\end{tabular}
\par\endgroup

% mdtable:T44:end
\end{minipage}\par
\reserve{7\baselineskip}
% mdparagraph:P194:begin
\subsection*{dēma}
% mdparagraph:P194:end

% teaching: NEW-U17-VOC-DEMA-U17-SENSE
\par\begin{minipage}{\linewidth}\raggedright
% mdparagraph:P195:begin
\textbf{dēma} — judge.
% mdparagraph:P195:end

% teaching: NEW-U17-NOM-DEMA
% mdparagraph:P196:begin
Masculine; stem \textbf{dēm-}. Add the endings of the stated pattern to this stem, preserving its vowel length.
% mdparagraph:P196:end

% table: NEW-U17-NOM-DEMA-DISPLAY
% mdtable:T45:begin
\par\begingroup
\fontsize{10.5}{13.2}\selectfont
\setlength{\LTleft}{0pt}\setlength{\LTright}{\fill}
\noindent
\begin{tabular}{@{}>{\raggedright\arraybackslash}p{39.46mm}>{\raggedright\arraybackslash}p{39.46mm}>{\raggedright\arraybackslash}p{39.46mm}@{}}
\toprule
\textbf{Case} & \textbf{Singular} & \textbf{Plural} \\
\midrule
Nominative & dēma & dēman \\
Genitive & dēman & dēmena \\
Dative & dēman & dēmum \\
Accusative & dēman & dēman \\
\bottomrule
\end{tabular}
\par\endgroup

% mdtable:T45:end
\end{minipage}\par
\reserve{7\baselineskip}
% mdparagraph:P197:begin
\subsection*{hālga}
% mdparagraph:P197:end

% teaching: NEW-U17-VOC-HALGA-U17-SENSE
\par\begin{minipage}{\linewidth}\raggedright
% mdparagraph:P198:begin
\textbf{hālga} — saint; holy one.
% mdparagraph:P198:end

% teaching: NEW-U17-NOM-HALGA
% mdparagraph:P199:begin
Masculine; stem \textbf{hālg-}. This is a weak masculine noun, saint or holy one. Its genitive plural is hālgena. Keep it distinct from the adjective hālig and its genitive plural hālgra.
% mdparagraph:P199:end

% table: NEW-U17-NOM-HALGA-DISPLAY
% mdtable:T46:begin
\par\begingroup
\fontsize{10.5}{13.2}\selectfont
\setlength{\LTleft}{0pt}\setlength{\LTright}{\fill}
\noindent
\begin{tabular}{@{}>{\raggedright\arraybackslash}p{39.46mm}>{\raggedright\arraybackslash}p{39.46mm}>{\raggedright\arraybackslash}p{39.46mm}@{}}
\toprule
\textbf{Case} & \textbf{Singular} & \textbf{Plural} \\
\midrule
Nominative & hālga & hālgan \\
Genitive & hālgan & hālgena \\
Dative & hālgan & hālgum \\
Accusative & hālgan & hālgan \\
\bottomrule
\end{tabular}
\par\endgroup

% mdtable:T46:end
\end{minipage}\par
\reserve{7\baselineskip}
% mdparagraph:P200:begin
\subsection*{hlīsa}
% mdparagraph:P200:end

% teaching: NEW-U17-VOC-HLISA-U17-FAME
% teaching: NEW-U17-VOC-HLISA-U17-REPORT
\par\begin{minipage}{\linewidth}\raggedright
% mdparagraph:P201:begin
\textbf{hlīsa} — fame; report; account heard.
% mdparagraph:P201:end

% teaching: NEW-U17-NOM-HLISA
% mdparagraph:P202:begin
Masculine; stem \textbf{hlīs-}. Fame is ordinarily singular in this passage. A report or account can be treated as an individual item. The complete table serves those stated meanings; context determines which is intended.
% mdparagraph:P202:end

% table: NEW-U17-NOM-HLISA-DISPLAY
% mdtable:T47:begin
\par\begingroup
\fontsize{10.5}{13.2}\selectfont
\setlength{\LTleft}{0pt}\setlength{\LTright}{\fill}
\noindent
\begin{tabular}{@{}>{\raggedright\arraybackslash}p{39.46mm}>{\raggedright\arraybackslash}p{39.46mm}>{\raggedright\arraybackslash}p{39.46mm}@{}}
\toprule
\textbf{Case} & \textbf{Singular} & \textbf{Plural} \\
\midrule
Nominative & hlīsa & hlīsan \\
Genitive & hlīsan & hlīsena \\
Dative & hlīsan & hlīsum \\
Accusative & hlīsan & hlīsan \\
\bottomrule
\end{tabular}
\par\endgroup

% mdtable:T47:end
\end{minipage}\par
\reserve{7\baselineskip}
% mdparagraph:P203:begin
\subsection*{ðēowa}
% mdparagraph:P203:end

% teaching: NEW-U17-VOC-THEOWA-U17-SENSE
\par\begin{minipage}{\linewidth}\raggedright
% mdparagraph:P204:begin
\textbf{ðēowa} — servant.
% mdparagraph:P204:end

% teaching: NEW-U17-NOM-THEOWA
% mdparagraph:P205:begin
Masculine; stem \textbf{ðēow-}. This is weak ðēowa, with nominative/\allowbreak{}accusative plural ðēowan. Do not substitute the plural of the separate strong form ðēow.
% mdparagraph:P205:end

% table: NEW-U17-NOM-THEOWA-DISPLAY
% mdtable:T48:begin
\par\begingroup
\fontsize{10.5}{13.2}\selectfont
\setlength{\LTleft}{0pt}\setlength{\LTright}{\fill}
\noindent
\begin{tabular}{@{}>{\raggedright\arraybackslash}p{39.46mm}>{\raggedright\arraybackslash}p{39.46mm}>{\raggedright\arraybackslash}p{39.46mm}@{}}
\toprule
\textbf{Case} & \textbf{Singular} & \textbf{Plural} \\
\midrule
Nominative & ðēowa & ðēowan \\
Genitive & ðēowan & ðēowena \\
Dative & ðēowan & ðēowum \\
Accusative & ðēowan & ðēowan \\
\bottomrule
\end{tabular}
\par\endgroup

% mdtable:T48:end
\end{minipage}\par
\reserve{7\baselineskip}
% mdparagraph:P206:begin
\subsection*{ēare}
% mdparagraph:P206:end

% teaching: NEW-U17-VOC-EARE-U17-SENSE
\par\begin{minipage}{\linewidth}\raggedright
% mdparagraph:P207:begin
\textbf{ēare} — ear.
% mdparagraph:P207:end

% teaching: NEW-U17-NOM-EARE
% mdparagraph:P208:begin
Neuter; stem \textbf{ēar-}. This weak neuter has ēare in both nominative and accusative singular. The other singular cases end in -an. Its genitive plural contracts to ēarna; the dative plural is ēarum.
% mdparagraph:P208:end

% table: NEW-U17-NOM-EARE-DISPLAY
% mdtable:T49:begin
\par\begingroup
\fontsize{10.5}{13.2}\selectfont
\setlength{\LTleft}{0pt}\setlength{\LTright}{\fill}
\noindent
\begin{tabular}{@{}>{\raggedright\arraybackslash}p{39.46mm}>{\raggedright\arraybackslash}p{39.46mm}>{\raggedright\arraybackslash}p{39.46mm}@{}}
\toprule
\textbf{Case} & \textbf{Singular} & \textbf{Plural} \\
\midrule
Nominative & ēare & ēaran \\
Genitive & ēaran & ēarna \\
Dative & ēaran & ēarum \\
Accusative & ēare & ēaran \\
\bottomrule
\end{tabular}
\par\endgroup

% mdtable:T49:end
\end{minipage}\par
\reserve{7\baselineskip}
% mdparagraph:P209:begin
\subsection*{tabele}
% mdparagraph:P209:end

% teaching: NEW-U17-VOC-TABELE-U17-SENSE
\par\begin{minipage}{\linewidth}\raggedright
% mdparagraph:P210:begin
\textbf{tabele} — tablet; table of the law.
% mdparagraph:P210:end

% teaching: NEW-U17-NOM-TABELE
% mdparagraph:P211:begin
Feminine; stem \textbf{tabel-}. The weak feminine stem is tabel-. Retain e in the dative plural tabelum.
% mdparagraph:P211:end

% table: NEW-U17-NOM-TABELE-DISPLAY
% mdtable:T50:begin
\par\begingroup
\fontsize{10.5}{13.2}\selectfont
\setlength{\LTleft}{0pt}\setlength{\LTright}{\fill}
\noindent
\begin{tabular}{@{}>{\raggedright\arraybackslash}p{39.46mm}>{\raggedright\arraybackslash}p{39.46mm}>{\raggedright\arraybackslash}p{39.46mm}@{}}
\toprule
\textbf{Case} & \textbf{Singular} & \textbf{Plural} \\
\midrule
Nominative & tabele & tabelan \\
Genitive & tabelan & tabelena \\
Dative & tabelan & tabelum \\
Accusative & tabelan & tabelan \\
\bottomrule
\end{tabular}
\par\endgroup

% mdtable:T50:end
\end{minipage}\par
\reserve{7\baselineskip}
% mdparagraph:P212:begin
\subsection*{hearpe}
% mdparagraph:P212:end

% teaching: NEW-U17-VOC-HEARPE-U17-SENSE
\par\begin{minipage}{\linewidth}\raggedright
% mdparagraph:P213:begin
\textbf{hearpe} — harp.
% mdparagraph:P213:end

% teaching: NEW-U17-NOM-HEARPE
% mdparagraph:P214:begin
Feminine; stem \textbf{hearp-}. Add the endings of the stated pattern to this stem, preserving its vowel length.
% mdparagraph:P214:end

% table: NEW-U17-NOM-HEARPE-DISPLAY
% mdtable:T51:begin
\par\begingroup
\fontsize{10.5}{13.2}\selectfont
\setlength{\LTleft}{0pt}\setlength{\LTright}{\fill}
\noindent
\begin{tabular}{@{}>{\raggedright\arraybackslash}p{39.46mm}>{\raggedright\arraybackslash}p{39.46mm}>{\raggedright\arraybackslash}p{39.46mm}@{}}
\toprule
\textbf{Case} & \textbf{Singular} & \textbf{Plural} \\
\midrule
Nominative & hearpe & hearpan \\
Genitive & hearpan & hearpena \\
Dative & hearpan & hearpum \\
Accusative & hearpan & hearpan \\
\bottomrule
\end{tabular}
\par\endgroup

% mdtable:T51:end
\end{minipage}\par
\reserve{7\baselineskip}
% mdparagraph:P215:begin
\subsection*{blōstma}
% mdparagraph:P215:end

% teaching: NEW-U17-VOC-BLOSTMA-U17-SENSE
\par\begin{minipage}{\linewidth}\raggedright
% mdparagraph:P216:begin
\textbf{blōstma} — blossom.
% mdparagraph:P216:end

% teaching: NEW-U17-NOM-BLOSTMA
% mdparagraph:P217:begin
Masculine; stem \textbf{blōstm-}. Add the endings of the stated pattern to this stem, preserving its vowel length.
% mdparagraph:P217:end

% table: NEW-U17-NOM-BLOSTMA-DISPLAY
% mdtable:T52:begin
\par\begingroup
\fontsize{10.5}{13.2}\selectfont
\setlength{\LTleft}{0pt}\setlength{\LTright}{\fill}
\noindent
\begin{tabular}{@{}>{\raggedright\arraybackslash}p{39.46mm}>{\raggedright\arraybackslash}p{39.46mm}>{\raggedright\arraybackslash}p{39.46mm}@{}}
\toprule
\textbf{Case} & \textbf{Singular} & \textbf{Plural} \\
\midrule
Nominative & blōstma & blōstman \\
Genitive & blōstman & blōstmena \\
Dative & blōstman & blōstmum \\
Accusative & blōstman & blōstman \\
\bottomrule
\end{tabular}
\par\endgroup

% mdtable:T52:end
\end{minipage}\par
\reserve{7\baselineskip}
% mdparagraph:P218:begin
\subsection*{biema}
% mdparagraph:P218:end

% teaching: NEW-U17-VOC-BIEMA-U17-SENSE
\par\begin{minipage}{\linewidth}\raggedright
% mdparagraph:P219:begin
\textbf{biema} — trumpet; clarion.
% mdparagraph:P219:end

% teaching: NEW-U17-NOM-BIEMA
% mdparagraph:P220:begin
Masculine; stem \textbf{biem-}. Retain the source spelling biema, without adding a length mark. It follows the weak masculine pattern.
% mdparagraph:P220:end

% table: NEW-U17-NOM-BIEMA-DISPLAY
% mdtable:T53:begin
\par\begingroup
\fontsize{10.5}{13.2}\selectfont
\setlength{\LTleft}{0pt}\setlength{\LTright}{\fill}
\noindent
\begin{tabular}{@{}>{\raggedright\arraybackslash}p{39.46mm}>{\raggedright\arraybackslash}p{39.46mm}>{\raggedright\arraybackslash}p{39.46mm}@{}}
\toprule
\textbf{Case} & \textbf{Singular} & \textbf{Plural} \\
\midrule
Nominative & biema & bieman \\
Genitive & bieman & biemena \\
Dative & bieman & biemum \\
Accusative & bieman & bieman \\
\bottomrule
\end{tabular}
\par\endgroup

% mdtable:T53:end
\end{minipage}\par
\reserve{7\baselineskip}
% mdparagraph:P221:begin
\subsection*{hlǣfdige}
% mdparagraph:P221:end

% teaching: NEW-U17-VOC-HLAEFDIGE-U17-SENSE
\par\begin{minipage}{\linewidth}\raggedright
% mdparagraph:P222:begin
\textbf{hlǣfdige} — lady.
% mdparagraph:P222:end

% teaching: NEW-U17-NOM-HLAEFDIGE
% mdparagraph:P223:begin
Feminine; stem \textbf{hlǣfdig-}. Add the endings of the stated pattern to this stem, preserving its vowel length.
% mdparagraph:P223:end

% table: NEW-U17-NOM-HLAEFDIGE-DISPLAY
% mdtable:T54:begin
\par\begingroup
\fontsize{10.5}{13.2}\selectfont
\setlength{\LTleft}{0pt}\setlength{\LTright}{\fill}
\noindent
\begin{tabular}{@{}>{\raggedright\arraybackslash}p{39.46mm}>{\raggedright\arraybackslash}p{39.46mm}>{\raggedright\arraybackslash}p{39.46mm}@{}}
\toprule
\textbf{Case} & \textbf{Singular} & \textbf{Plural} \\
\midrule
Nominative & hlǣfdige & hlǣfdigan \\
Genitive & hlǣfdigan & hlǣfdigena \\
Dative & hlǣfdigan & hlǣfdigum \\
Accusative & hlǣfdigan & hlǣfdigan \\
\bottomrule
\end{tabular}
\par\endgroup

% mdtable:T54:end
\end{minipage}\par
\reserve{7\baselineskip}
% mdparagraph:P224:begin
\subsection*{Strong and weak adjectives}
% mdparagraph:P224:end

% mdparagraph:P225:begin
An adjective agrees with its noun in gender, number and case. Use the familiar strong and weak choice: an article or comparable definite determiner normally calls for weak inflection, while a predicative adjective is normally strong. Direct address also favours weak forms; the praise catalogue uses this pattern even while its global instruction removes the article. The selected personal uses are explained after the tables.
% mdparagraph:P225:end

% mdparagraph:P226:begin
Each panel has masculine, neuter and feminine columns; the singular panel precedes the plural panel. Strong instrumental singular forms in the masculine and neuter columns are for recognition; their display adds no new instrumental construction. A dash is an inapplicable display position. Weak genitive plural -ra is the selected form. These tables print one consistent set of source-licensed forms; a documented alternative is not automatically ungrammatical.
% mdparagraph:P226:end

\reserve{7\baselineskip}
% mdparagraph:P227:begin
\subsection*{wæstmbǣre}
% mdparagraph:P227:end

% teaching: NEW-U17-VOC-WAESTMBAERE-U17-SENSE
\par\begin{minipage}{\linewidth}\raggedright
% mdparagraph:P228:begin
\textbf{wæstmbǣre} — fruitful.
% mdparagraph:P228:end

% teaching: NEW-U17-ADJ-WAESTMBAERE
% mdparagraph:P229:begin
The headword ends in -e; replace that -e with an overt ending on \textbf{wæstmbǣr-}. The displayed strong forms select feminine nominative singular and neuter nominative/\allowbreak{}accusative plural -u, and feminine plural -a. The source also permits -e in those positions.
% mdparagraph:P229:end

% mdparagraph:P230:begin
\textbf{Strong.}
% mdparagraph:P230:end

% table: NEW-U17-ADJ-WAESTMBAERE-STRONG-DISPLAY
% mdtable:T55:begin
% course-table-style: 2026-10
% course-panel:T55-P01:begin
\par\begingroup
\fontsize{10.5}{13.2}\selectfont
\setlength{\LTleft}{0pt}\setlength{\LTright}{\fill}
\noindent
\begin{tabular}{@{}>{\raggedright\arraybackslash}p{22.0mm}>{\raggedright\arraybackslash}p{31.188mm}>{\raggedright\arraybackslash}p{31.188mm}>{\raggedright\arraybackslash}p{31.188mm}@{}}
\toprule
\textbf{Case} & \textbf{Masculine singular} & \textbf{Neuter singular} & \textbf{Feminine singular} \\
\midrule
Nominative & wæstmbǣre & wæstmbǣre & wæstmbǣru \\
Genitive & wæstmbǣres & wæstmbǣres & wæstmbǣrre \\
Dative & wæstmbǣrum & wæstmbǣrum & wæstmbǣrre \\
Accusative & wæstmbǣrne & wæstmbǣre & wæstmbǣre \\
Instrumental & wæstmbǣre & wæstmbǣre & — \\
\bottomrule
\end{tabular}
\par\endgroup
% course-panel:T55-P01:end
\end{minipage}\par
\par\addvspace{5pt}
% course-panel:T55-P02:begin
\par\begingroup
\fontsize{10.5}{13.2}\selectfont
\setlength{\LTleft}{0pt}\setlength{\LTright}{\fill}
\noindent
\begin{tabular}{@{}>{\raggedright\arraybackslash}p{22.0mm}>{\raggedright\arraybackslash}p{31.188mm}>{\raggedright\arraybackslash}p{31.188mm}>{\raggedright\arraybackslash}p{31.188mm}@{}}
\toprule
\textbf{Case} & \textbf{Masculine plural} & \textbf{Neuter plural} & \textbf{Feminine plural} \\
\midrule
Nominative & wæstmbǣre & wæstmbǣru & wæstmbǣra \\
Genitive & wæstmbǣrra & wæstmbǣrra & wæstmbǣrra \\
Dative & wæstmbǣrum & wæstmbǣrum & wæstmbǣrum \\
Accusative & wæstmbǣre & wæstmbǣru & wæstmbǣra \\
Instrumental & — & — & — \\
\bottomrule
\end{tabular}
\par\endgroup
% course-panel:T55-P02:end
% mdtable:T55:end

\reserve{18\baselineskip}
\reserve{16\baselineskip}
% mdparagraph:P231:begin
\textbf{Weak.}
% mdparagraph:P231:end

% table: NEW-U17-ADJ-WAESTMBAERE-WEAK-DISPLAY
% mdtable:T56:begin
% course-table-style: 2026-10
% course-panel:T56-P01:begin
\par\begingroup
\fontsize{10.5}{13.2}\selectfont
\setlength{\LTleft}{0pt}\setlength{\LTright}{\fill}
\noindent
\begin{tabular}{@{}>{\raggedright\arraybackslash}p{22.0mm}>{\raggedright\arraybackslash}p{31.188mm}>{\raggedright\arraybackslash}p{31.188mm}>{\raggedright\arraybackslash}p{31.188mm}@{}}
\toprule
\textbf{Case} & \textbf{Masculine singular} & \textbf{Neuter singular} & \textbf{Feminine singular} \\
\midrule
Nominative & wæstmbǣra & wæstmbǣre & wæstmbǣre \\
Genitive & wæstmbǣran & wæstmbǣran & wæstmbǣran \\
Dative & wæstmbǣran & wæstmbǣran & wæstmbǣran \\
Accusative & wæstmbǣran & wæstmbǣre & wæstmbǣran \\
\bottomrule
\end{tabular}
\par\endgroup
% course-panel:T56-P01:end
\par\addvspace{5pt}
% course-panel:T56-P02:begin
\par\begingroup
\fontsize{10.5}{13.2}\selectfont
\setlength{\LTleft}{0pt}\setlength{\LTright}{\fill}
\noindent
\begin{tabular}{@{}>{\raggedright\arraybackslash}p{22.0mm}>{\raggedright\arraybackslash}p{31.188mm}>{\raggedright\arraybackslash}p{31.188mm}>{\raggedright\arraybackslash}p{31.188mm}@{}}
\toprule
\textbf{Case} & \textbf{Masculine plural} & \textbf{Neuter plural} & \textbf{Feminine plural} \\
\midrule
Nominative & wæstmbǣran & wæstmbǣran & wæstmbǣran \\
Genitive & wæstmbǣrra & wæstmbǣrra & wæstmbǣrra \\
Dative & wæstmbǣrum & wæstmbǣrum & wæstmbǣrum \\
Accusative & wæstmbǣran & wæstmbǣran & wæstmbǣran \\
\bottomrule
\end{tabular}
\par\endgroup
% course-panel:T56-P02:end
% mdtable:T56:end

\reserve{7\baselineskip}
% mdparagraph:P232:begin
\subsection*{geong}
% mdparagraph:P232:end

% teaching: NEW-U17-VOC-GEONG-YOUNG-SOURCE
\par\begin{minipage}{\linewidth}\raggedright
% mdparagraph:P233:begin
\textbf{geong} — young; substantivally young people.
% mdparagraph:P233:end

% teaching: NEW-U17-ADJ-GEONG
% mdparagraph:P234:begin
The stem \textbf{geong-} is heavy. Strong feminine nominative singular and neuter nominative/\allowbreak{}accusative plural therefore have no added -u.
% mdparagraph:P234:end

% mdparagraph:P235:begin
\textbf{Strong.}
% mdparagraph:P235:end

% table: NEW-U17-ADJ-GEONG-STRONG-DISPLAY
% mdtable:T57:begin
% course-table-style: 2026-10
% course-panel:T57-P01:begin
\par\begingroup
\fontsize{10.5}{13.2}\selectfont
\setlength{\LTleft}{0pt}\setlength{\LTright}{\fill}
\noindent
\begin{tabular}{@{}>{\raggedright\arraybackslash}p{22.0mm}>{\raggedright\arraybackslash}p{31.188mm}>{\raggedright\arraybackslash}p{31.188mm}>{\raggedright\arraybackslash}p{31.188mm}@{}}
\toprule
\textbf{Case} & \textbf{Masculine singular} & \textbf{Neuter singular} & \textbf{Feminine singular} \\
\midrule
Nominative & geong & geong & geong \\
Genitive & geonges & geonges & geongre \\
Dative & geongum & geongum & geongre \\
Accusative & geongne & geong & geonge \\
Instrumental & geonge & geonge & — \\
\bottomrule
\end{tabular}
\par\endgroup
% course-panel:T57-P01:end
\end{minipage}\par
\par\addvspace{5pt}
% course-panel:T57-P02:begin
\par\begingroup
\fontsize{10.5}{13.2}\selectfont
\setlength{\LTleft}{0pt}\setlength{\LTright}{\fill}
\noindent
\begin{tabular}{@{}>{\raggedright\arraybackslash}p{22.0mm}>{\raggedright\arraybackslash}p{31.188mm}>{\raggedright\arraybackslash}p{31.188mm}>{\raggedright\arraybackslash}p{31.188mm}@{}}
\toprule
\textbf{Case} & \textbf{Masculine plural} & \textbf{Neuter plural} & \textbf{Feminine plural} \\
\midrule
Nominative & geonge & geong & geonga \\
Genitive & geongra & geongra & geongra \\
Dative & geongum & geongum & geongum \\
Accusative & geonge & geong & geonga \\
Instrumental & — & — & — \\
\bottomrule
\end{tabular}
\par\endgroup
% course-panel:T57-P02:end
% mdtable:T57:end

\reserve{18\baselineskip}
\reserve{16\baselineskip}
% mdparagraph:P236:begin
\textbf{Weak.}
% mdparagraph:P236:end

% table: NEW-U17-ADJ-GEONG-WEAK-DISPLAY
% mdtable:T58:begin
% course-table-style: 2026-10
% course-panel:T58-P01:begin
\par\begingroup
\fontsize{10.5}{13.2}\selectfont
\setlength{\LTleft}{0pt}\setlength{\LTright}{\fill}
\noindent
\begin{tabular}{@{}>{\raggedright\arraybackslash}p{22.0mm}>{\raggedright\arraybackslash}p{31.188mm}>{\raggedright\arraybackslash}p{31.188mm}>{\raggedright\arraybackslash}p{31.188mm}@{}}
\toprule
\textbf{Case} & \textbf{Masculine singular} & \textbf{Neuter singular} & \textbf{Feminine singular} \\
\midrule
Nominative & geonga & geonge & geonge \\
Genitive & geongan & geongan & geongan \\
Dative & geongan & geongan & geongan \\
Accusative & geongan & geonge & geongan \\
\bottomrule
\end{tabular}
\par\endgroup
% course-panel:T58-P01:end
\par\addvspace{5pt}
% course-panel:T58-P02:begin
\par\begingroup
\fontsize{10.5}{13.2}\selectfont
\setlength{\LTleft}{0pt}\setlength{\LTright}{\fill}
\noindent
\begin{tabular}{@{}>{\raggedright\arraybackslash}p{22.0mm}>{\raggedright\arraybackslash}p{31.188mm}>{\raggedright\arraybackslash}p{31.188mm}>{\raggedright\arraybackslash}p{31.188mm}@{}}
\toprule
\textbf{Case} & \textbf{Masculine plural} & \textbf{Neuter plural} & \textbf{Feminine plural} \\
\midrule
Nominative & geongan & geongan & geongan \\
Genitive & geongra & geongra & geongra \\
Dative & geongum & geongum & geongum \\
Accusative & geongan & geongan & geongan \\
\bottomrule
\end{tabular}
\par\endgroup
% course-panel:T58-P02:end
% mdtable:T58:end

\reserve{7\baselineskip}
% mdparagraph:P237:begin
\subsection*{manigfeald}
% mdparagraph:P237:end

% teaching: NEW-U17-VOC-MANIGFEALD-U15-SOURCE-SENSE
\par\begin{minipage}{\linewidth}\raggedright
% mdparagraph:P238:begin
\textbf{manigfeald} — manifold; numerous and varied.
% mdparagraph:P238:end

% teaching: NEW-U17-ADJ-MANIGFEALD
% mdparagraph:P239:begin
The stem \textbf{manigfeald-} is heavy. Strong feminine nominative singular and neuter nominative/\allowbreak{}accusative plural therefore have no added -u.
% mdparagraph:P239:end

% mdparagraph:P240:begin
\textbf{Strong.}
% mdparagraph:P240:end

% table: NEW-U17-ADJ-MANIGFEALD-STRONG-DISPLAY
% mdtable:T59:begin
% course-table-style: 2026-10
% course-panel:T59-P01:begin
\par\begingroup
\fontsize{10.5}{13.2}\selectfont
\setlength{\LTleft}{0pt}\setlength{\LTright}{\fill}
\noindent
\begin{tabular}{@{}>{\raggedright\arraybackslash}p{22.0mm}>{\raggedright\arraybackslash}p{31.188mm}>{\raggedright\arraybackslash}p{31.188mm}>{\raggedright\arraybackslash}p{31.188mm}@{}}
\toprule
\textbf{Case} & \textbf{Masculine singular} & \textbf{Neuter singular} & \textbf{Feminine singular} \\
\midrule
Nominative & manigfeald & manigfeald & manigfeald \\
Genitive & manigfealdes & manigfealdes & manigfealdre \\
Dative & manigfealdum & manigfealdum & manigfealdre \\
Accusative & manigfealdne & manigfeald & manigfealde \\
Instrumental & manigfealde & manigfealde & — \\
\bottomrule
\end{tabular}
\par\endgroup
% course-panel:T59-P01:end
\end{minipage}\par
\par\addvspace{5pt}
% course-panel:T59-P02:begin
\par\begingroup
\fontsize{10.5}{13.2}\selectfont
\setlength{\LTleft}{0pt}\setlength{\LTright}{\fill}
\noindent
\begin{tabular}{@{}>{\raggedright\arraybackslash}p{22.0mm}>{\raggedright\arraybackslash}p{31.188mm}>{\raggedright\arraybackslash}p{31.188mm}>{\raggedright\arraybackslash}p{31.188mm}@{}}
\toprule
\textbf{Case} & \textbf{Masculine plural} & \textbf{Neuter plural} & \textbf{Feminine plural} \\
\midrule
Nominative & manigfealde & manigfeald & manigfealda \\
Genitive & manigfealdra & manigfealdra & manigfealdra \\
Dative & manigfealdum & manigfealdum & manigfealdum \\
Accusative & manigfealde & manigfeald & manigfealda \\
Instrumental & — & — & — \\
\bottomrule
\end{tabular}
\par\endgroup
% course-panel:T59-P02:end
% mdtable:T59:end

\reserve{18\baselineskip}
\reserve{16\baselineskip}
% mdparagraph:P241:begin
\textbf{Weak.}
% mdparagraph:P241:end

% table: NEW-U17-ADJ-MANIGFEALD-WEAK-DISPLAY
% mdtable:T60:begin
% course-table-style: 2026-10
% course-panel:T60-P01:begin
\par\begingroup
\fontsize{10.5}{13.2}\selectfont
\setlength{\LTleft}{0pt}\setlength{\LTright}{\fill}
\noindent
\begin{tabular}{@{}>{\raggedright\arraybackslash}p{22.0mm}>{\raggedright\arraybackslash}p{31.188mm}>{\raggedright\arraybackslash}p{31.188mm}>{\raggedright\arraybackslash}p{31.188mm}@{}}
\toprule
\textbf{Case} & \textbf{Masculine singular} & \textbf{Neuter singular} & \textbf{Feminine singular} \\
\midrule
Nominative & manigfealda & manigfealde & manigfealde \\
Genitive & manigfealdan & manigfealdan & manigfealdan \\
Dative & manigfealdan & manigfealdan & manigfealdan \\
Accusative & manigfealdan & manigfealde & manigfealdan \\
\bottomrule
\end{tabular}
\par\endgroup
% course-panel:T60-P01:end
\par\addvspace{5pt}
% course-panel:T60-P02:begin
\par\begingroup
\fontsize{10.5}{13.2}\selectfont
\setlength{\LTleft}{0pt}\setlength{\LTright}{\fill}
\noindent
\begin{tabular}{@{}>{\raggedright\arraybackslash}p{22.0mm}>{\raggedright\arraybackslash}p{31.188mm}>{\raggedright\arraybackslash}p{31.188mm}>{\raggedright\arraybackslash}p{31.188mm}@{}}
\toprule
\textbf{Case} & \textbf{Masculine plural} & \textbf{Neuter plural} & \textbf{Feminine plural} \\
\midrule
Nominative & manigfealdan & manigfealdan & manigfealdan \\
Genitive & manigfealdra & manigfealdra & manigfealdra \\
Dative & manigfealdum & manigfealdum & manigfealdum \\
Accusative & manigfealdan & manigfealdan & manigfealdan \\
\bottomrule
\end{tabular}
\par\endgroup
% course-panel:T60-P02:end
% mdtable:T60:end

\reserve{7\baselineskip}
% mdparagraph:P242:begin
\subsection*{rīce}
% mdparagraph:P242:end

% teaching: NEW-U17-VOC-RICE-ADJECTIVE-U17-SENSE
\par\begin{minipage}{\linewidth}\raggedright
% mdparagraph:P243:begin
\textbf{rīce} — rich; substantivally rich people.
% mdparagraph:P243:end

% teaching: NEW-U17-ADJ-RICE-ADJECTIVE
% mdparagraph:P244:begin
The headword ends in -e; replace that -e with an overt ending on \textbf{rīc-}. The displayed strong forms select feminine nominative singular and neuter nominative/\allowbreak{}accusative plural -u, and feminine plural -a. The source also permits -e in those positions.
% mdparagraph:P244:end

% mdparagraph:P245:begin
\textbf{Strong.}
% mdparagraph:P245:end

% table: NEW-U17-ADJ-RICE-ADJECTIVE-STRONG-DISPLAY
% mdtable:T61:begin
% course-table-style: 2026-10
% course-panel:T61-P01:begin
\par\begingroup
\fontsize{10.5}{13.2}\selectfont
\setlength{\LTleft}{0pt}\setlength{\LTright}{\fill}
\noindent
\begin{tabular}{@{}>{\raggedright\arraybackslash}p{22.0mm}>{\raggedright\arraybackslash}p{31.188mm}>{\raggedright\arraybackslash}p{31.188mm}>{\raggedright\arraybackslash}p{31.188mm}@{}}
\toprule
\textbf{Case} & \textbf{Masculine singular} & \textbf{Neuter singular} & \textbf{Feminine singular} \\
\midrule
Nominative & rīce & rīce & rīcu \\
Genitive & rīces & rīces & rīcre \\
Dative & rīcum & rīcum & rīcre \\
Accusative & rīcne & rīce & rīce \\
Instrumental & rīce & rīce & — \\
\bottomrule
\end{tabular}
\par\endgroup
% course-panel:T61-P01:end
\end{minipage}\par
\par\addvspace{5pt}
% course-panel:T61-P02:begin
\par\begingroup
\fontsize{10.5}{13.2}\selectfont
\setlength{\LTleft}{0pt}\setlength{\LTright}{\fill}
\noindent
\begin{tabular}{@{}>{\raggedright\arraybackslash}p{22.0mm}>{\raggedright\arraybackslash}p{31.188mm}>{\raggedright\arraybackslash}p{31.188mm}>{\raggedright\arraybackslash}p{31.188mm}@{}}
\toprule
\textbf{Case} & \textbf{Masculine plural} & \textbf{Neuter plural} & \textbf{Feminine plural} \\
\midrule
Nominative & rīce & rīcu & rīca \\
Genitive & rīcra & rīcra & rīcra \\
Dative & rīcum & rīcum & rīcum \\
Accusative & rīce & rīcu & rīca \\
Instrumental & — & — & — \\
\bottomrule
\end{tabular}
\par\endgroup
% course-panel:T61-P02:end
% mdtable:T61:end

\reserve{18\baselineskip}
\reserve{16\baselineskip}
% mdparagraph:P246:begin
\textbf{Weak.}
% mdparagraph:P246:end

% table: NEW-U17-ADJ-RICE-ADJECTIVE-WEAK-DISPLAY
% mdtable:T62:begin
% course-table-style: 2026-10
% course-panel:T62-P01:begin
\par\begingroup
\fontsize{10.5}{13.2}\selectfont
\setlength{\LTleft}{0pt}\setlength{\LTright}{\fill}
\noindent
\begin{tabular}{@{}>{\raggedright\arraybackslash}p{22.0mm}>{\raggedright\arraybackslash}p{31.188mm}>{\raggedright\arraybackslash}p{31.188mm}>{\raggedright\arraybackslash}p{31.188mm}@{}}
\toprule
\textbf{Case} & \textbf{Masculine singular} & \textbf{Neuter singular} & \textbf{Feminine singular} \\
\midrule
Nominative & rīca & rīce & rīce \\
Genitive & rīcan & rīcan & rīcan \\
Dative & rīcan & rīcan & rīcan \\
Accusative & rīcan & rīce & rīcan \\
\bottomrule
\end{tabular}
\par\endgroup
% course-panel:T62-P01:end
\par\addvspace{5pt}
% course-panel:T62-P02:begin
\par\begingroup
\fontsize{10.5}{13.2}\selectfont
\setlength{\LTleft}{0pt}\setlength{\LTright}{\fill}
\noindent
\begin{tabular}{@{}>{\raggedright\arraybackslash}p{22.0mm}>{\raggedright\arraybackslash}p{31.188mm}>{\raggedright\arraybackslash}p{31.188mm}>{\raggedright\arraybackslash}p{31.188mm}@{}}
\toprule
\textbf{Case} & \textbf{Masculine plural} & \textbf{Neuter plural} & \textbf{Feminine plural} \\
\midrule
Nominative & rīcan & rīcan & rīcan \\
Genitive & rīcra & rīcra & rīcra \\
Dative & rīcum & rīcum & rīcum \\
Accusative & rīcan & rīcan & rīcan \\
\bottomrule
\end{tabular}
\par\endgroup
% course-panel:T62-P02:end
% mdtable:T62:end

\reserve{7\baselineskip}
% mdparagraph:P247:begin
\subsection*{dæghwāmlīc}
% mdparagraph:P247:end

% teaching: NEW-U17-VOC-DAEGHWAMLIC-U17-SENSE
\par\begin{minipage}{\linewidth}\raggedright
% mdparagraph:P248:begin
\textbf{dæghwāmlīc} — daily.
% mdparagraph:P248:end

% teaching: NEW-U17-ADJ-DAEGHWAMLIC
% mdparagraph:P249:begin
The stem \textbf{dæghwāmlīc-} is heavy. Strong feminine nominative singular and neuter nominative/\allowbreak{}accusative plural therefore have no added -u.
% mdparagraph:P249:end

% mdparagraph:P250:begin
\textbf{Strong.}
% mdparagraph:P250:end

% table: NEW-U17-ADJ-DAEGHWAMLIC-STRONG-DISPLAY
% mdtable:T63:begin
% course-table-style: 2026-10
% course-panel:T63-P01:begin
\par\begingroup
\fontsize{10.5}{13.2}\selectfont
\setlength{\LTleft}{0pt}\setlength{\LTright}{\fill}
\noindent
\begin{tabular}{@{}>{\raggedright\arraybackslash}p{22.0mm}>{\raggedright\arraybackslash}p{31.188mm}>{\raggedright\arraybackslash}p{31.188mm}>{\raggedright\arraybackslash}p{31.188mm}@{}}
\toprule
\textbf{Case} & \textbf{Masculine singular} & \textbf{Neuter singular} & \textbf{Feminine singular} \\
\midrule
Nominative & dæghwāmlīc & dæghwāmlīc & dæghwāmlīc \\
Genitive & dæghwāmlīces & dæghwāmlīces & dæghwāmlīcre \\
Dative & dæghwāmlīcum & dæghwāmlīcum & dæghwāmlīcre \\
Accusative & dæghwāmlīcne & dæghwāmlīc & dæghwāmlīce \\
Instrumental & dæghwāmlīce & dæghwāmlīce & — \\
\bottomrule
\end{tabular}
\par\endgroup
% course-panel:T63-P01:end
\end{minipage}\par
\par\addvspace{5pt}
% course-panel:T63-P02:begin
\par\begingroup
\fontsize{10.5}{13.2}\selectfont
\setlength{\LTleft}{0pt}\setlength{\LTright}{\fill}
\noindent
\begin{tabular}{@{}>{\raggedright\arraybackslash}p{22.0mm}>{\raggedright\arraybackslash}p{31.188mm}>{\raggedright\arraybackslash}p{31.188mm}>{\raggedright\arraybackslash}p{31.188mm}@{}}
\toprule
\textbf{Case} & \textbf{Masculine plural} & \textbf{Neuter plural} & \textbf{Feminine plural} \\
\midrule
Nominative & dæghwāmlīce & dæghwāmlīc & dæghwāmlīca \\
Genitive & dæghwāmlīcra & dæghwāmlīcra & dæghwāmlīcra \\
Dative & dæghwāmlīcum & dæghwāmlīcum & dæghwāmlīcum \\
Accusative & dæghwāmlīce & dæghwāmlīc & dæghwāmlīca \\
Instrumental & — & — & — \\
\bottomrule
\end{tabular}
\par\endgroup
% course-panel:T63-P02:end
% mdtable:T63:end

\reserve{18\baselineskip}
\reserve{16\baselineskip}
% mdparagraph:P251:begin
\textbf{Weak.}
% mdparagraph:P251:end

% table: NEW-U17-ADJ-DAEGHWAMLIC-WEAK-DISPLAY
% mdtable:T64:begin
% course-table-style: 2026-10
% course-panel:T64-P01:begin
\par\begingroup
\fontsize{10.5}{13.2}\selectfont
\setlength{\LTleft}{0pt}\setlength{\LTright}{\fill}
\noindent
\begin{tabular}{@{}>{\raggedright\arraybackslash}p{22.0mm}>{\raggedright\arraybackslash}p{31.188mm}>{\raggedright\arraybackslash}p{31.188mm}>{\raggedright\arraybackslash}p{31.188mm}@{}}
\toprule
\textbf{Case} & \textbf{Masculine singular} & \textbf{Neuter singular} & \textbf{Feminine singular} \\
\midrule
Nominative & dæghwāmlīca & dæghwāmlīce & dæghwāmlīce \\
Genitive & dæghwāmlīcan & dæghwāmlīcan & dæghwāmlīcan \\
Dative & dæghwāmlīcan & dæghwāmlīcan & dæghwāmlīcan \\
Accusative & dæghwāmlīcan & dæghwāmlīce & dæghwāmlīcan \\
\bottomrule
\end{tabular}
\par\endgroup
% course-panel:T64-P01:end
\par\addvspace{5pt}
% course-panel:T64-P02:begin
\par\begingroup
\fontsize{10.5}{13.2}\selectfont
\setlength{\LTleft}{0pt}\setlength{\LTright}{\fill}
\noindent
\begin{tabular}{@{}>{\raggedright\arraybackslash}p{22.0mm}>{\raggedright\arraybackslash}p{31.188mm}>{\raggedright\arraybackslash}p{31.188mm}>{\raggedright\arraybackslash}p{31.188mm}@{}}
\toprule
\textbf{Case} & \textbf{Masculine plural} & \textbf{Neuter plural} & \textbf{Feminine plural} \\
\midrule
Nominative & dæghwāmlīcan & dæghwāmlīcan & dæghwāmlīcan \\
Genitive & dæghwāmlīcra & dæghwāmlīcra & dæghwāmlīcra \\
Dative & dæghwāmlīcum & dæghwāmlīcum & dæghwāmlīcum \\
Accusative & dæghwāmlīcan & dæghwāmlīcan & dæghwāmlīcan \\
\bottomrule
\end{tabular}
\par\endgroup
% course-panel:T64-P02:end
% mdtable:T64:end

\reserve{7\baselineskip}
% mdparagraph:P252:begin
\subsection*{hlūd}
% mdparagraph:P252:end

% teaching: NEW-U17-VOC-HLUD-U17-SENSE
\par\begin{minipage}{\linewidth}\raggedright
% mdparagraph:P253:begin
\textbf{hlūd} — loud.
% mdparagraph:P253:end

% teaching: NEW-U17-ADJ-HLUD
% mdparagraph:P254:begin
The stem \textbf{hlūd-} is heavy. Strong feminine nominative singular and neuter nominative/\allowbreak{}accusative plural therefore have no added -u.
% mdparagraph:P254:end

% mdparagraph:P255:begin
\textbf{Strong.}
% mdparagraph:P255:end

% table: NEW-U17-ADJ-HLUD-STRONG-DISPLAY
% mdtable:T65:begin
% course-table-style: 2026-10
% course-panel:T65-P01:begin
\par\begingroup
\fontsize{10.5}{13.2}\selectfont
\setlength{\LTleft}{0pt}\setlength{\LTright}{\fill}
\noindent
\begin{tabular}{@{}>{\raggedright\arraybackslash}p{22.0mm}>{\raggedright\arraybackslash}p{31.188mm}>{\raggedright\arraybackslash}p{31.188mm}>{\raggedright\arraybackslash}p{31.188mm}@{}}
\toprule
\textbf{Case} & \textbf{Masculine singular} & \textbf{Neuter singular} & \textbf{Feminine singular} \\
\midrule
Nominative & hlūd & hlūd & hlūd \\
Genitive & hlūdes & hlūdes & hlūdre \\
Dative & hlūdum & hlūdum & hlūdre \\
Accusative & hlūdne & hlūd & hlūde \\
Instrumental & hlūde & hlūde & — \\
\bottomrule
\end{tabular}
\par\endgroup
% course-panel:T65-P01:end
\end{minipage}\par
\par\addvspace{5pt}
% course-panel:T65-P02:begin
\par\begingroup
\fontsize{10.5}{13.2}\selectfont
\setlength{\LTleft}{0pt}\setlength{\LTright}{\fill}
\noindent
\begin{tabular}{@{}>{\raggedright\arraybackslash}p{22.0mm}>{\raggedright\arraybackslash}p{31.188mm}>{\raggedright\arraybackslash}p{31.188mm}>{\raggedright\arraybackslash}p{31.188mm}@{}}
\toprule
\textbf{Case} & \textbf{Masculine plural} & \textbf{Neuter plural} & \textbf{Feminine plural} \\
\midrule
Nominative & hlūde & hlūd & hlūda \\
Genitive & hlūdra & hlūdra & hlūdra \\
Dative & hlūdum & hlūdum & hlūdum \\
Accusative & hlūde & hlūd & hlūda \\
Instrumental & — & — & — \\
\bottomrule
\end{tabular}
\par\endgroup
% course-panel:T65-P02:end
% mdtable:T65:end

\reserve{18\baselineskip}
\reserve{16\baselineskip}
% mdparagraph:P256:begin
\textbf{Weak.}
% mdparagraph:P256:end

% table: NEW-U17-ADJ-HLUD-WEAK-DISPLAY
% mdtable:T66:begin
% course-table-style: 2026-10
% course-panel:T66-P01:begin
\par\begingroup
\fontsize{10.5}{13.2}\selectfont
\setlength{\LTleft}{0pt}\setlength{\LTright}{\fill}
\noindent
\begin{tabular}{@{}>{\raggedright\arraybackslash}p{22.0mm}>{\raggedright\arraybackslash}p{31.188mm}>{\raggedright\arraybackslash}p{31.188mm}>{\raggedright\arraybackslash}p{31.188mm}@{}}
\toprule
\textbf{Case} & \textbf{Masculine singular} & \textbf{Neuter singular} & \textbf{Feminine singular} \\
\midrule
Nominative & hlūda & hlūde & hlūde \\
Genitive & hlūdan & hlūdan & hlūdan \\
Dative & hlūdan & hlūdan & hlūdan \\
Accusative & hlūdan & hlūde & hlūdan \\
\bottomrule
\end{tabular}
\par\endgroup
% course-panel:T66-P01:end
\par\addvspace{5pt}
% course-panel:T66-P02:begin
\par\begingroup
\fontsize{10.5}{13.2}\selectfont
\setlength{\LTleft}{0pt}\setlength{\LTright}{\fill}
\noindent
\begin{tabular}{@{}>{\raggedright\arraybackslash}p{22.0mm}>{\raggedright\arraybackslash}p{31.188mm}>{\raggedright\arraybackslash}p{31.188mm}>{\raggedright\arraybackslash}p{31.188mm}@{}}
\toprule
\textbf{Case} & \textbf{Masculine plural} & \textbf{Neuter plural} & \textbf{Feminine plural} \\
\midrule
Nominative & hlūdan & hlūdan & hlūdan \\
Genitive & hlūdra & hlūdra & hlūdra \\
Dative & hlūdum & hlūdum & hlūdum \\
Accusative & hlūdan & hlūdan & hlūdan \\
\bottomrule
\end{tabular}
\par\endgroup
% course-panel:T66-P02:end
% mdtable:T66:end

\reserve{7\baselineskip}
% mdparagraph:P257:begin
\subsection*{æðele}
% mdparagraph:P257:end

% teaching: NEW-U17-VOC-AETHELE-U17-NOBLE
\par\begin{minipage}{\linewidth}\raggedright
% mdparagraph:P258:begin
\textbf{æðele} — noble; illustrious.
% mdparagraph:P258:end

% teaching: NEW-U17-ADJ-AETHELE
% mdparagraph:P259:begin
The headword ends in -e; replace that -e with an overt ending on \textbf{æðel-}. The displayed strong forms select feminine nominative singular and neuter nominative/\allowbreak{}accusative plural -u, and feminine plural -a. The source also permits -e in those positions.
% mdparagraph:P259:end

% mdparagraph:P260:begin
\textbf{Strong.}
% mdparagraph:P260:end

% table: NEW-U17-ADJ-AETHELE-STRONG-DISPLAY
% mdtable:T67:begin
% course-table-style: 2026-10
% course-panel:T67-P01:begin
\par\begingroup
\fontsize{10.5}{13.2}\selectfont
\setlength{\LTleft}{0pt}\setlength{\LTright}{\fill}
\noindent
\begin{tabular}{@{}>{\raggedright\arraybackslash}p{22.0mm}>{\raggedright\arraybackslash}p{31.188mm}>{\raggedright\arraybackslash}p{31.188mm}>{\raggedright\arraybackslash}p{31.188mm}@{}}
\toprule
\textbf{Case} & \textbf{Masculine singular} & \textbf{Neuter singular} & \textbf{Feminine singular} \\
\midrule
Nominative & æðele & æðele & æðelu \\
Genitive & æðeles & æðeles & æðelre \\
Dative & æðelum & æðelum & æðelre \\
Accusative & æðelne & æðele & æðele \\
Instrumental & æðele & æðele & — \\
\bottomrule
\end{tabular}
\par\endgroup
% course-panel:T67-P01:end
\end{minipage}\par
\par\addvspace{5pt}
% course-panel:T67-P02:begin
\par\begingroup
\fontsize{10.5}{13.2}\selectfont
\setlength{\LTleft}{0pt}\setlength{\LTright}{\fill}
\noindent
\begin{tabular}{@{}>{\raggedright\arraybackslash}p{22.0mm}>{\raggedright\arraybackslash}p{31.188mm}>{\raggedright\arraybackslash}p{31.188mm}>{\raggedright\arraybackslash}p{31.188mm}@{}}
\toprule
\textbf{Case} & \textbf{Masculine plural} & \textbf{Neuter plural} & \textbf{Feminine plural} \\
\midrule
Nominative & æðele & æðelu & æðela \\
Genitive & æðelra & æðelra & æðelra \\
Dative & æðelum & æðelum & æðelum \\
Accusative & æðele & æðelu & æðela \\
Instrumental & — & — & — \\
\bottomrule
\end{tabular}
\par\endgroup
% course-panel:T67-P02:end
% mdtable:T67:end

\reserve{18\baselineskip}
\reserve{16\baselineskip}
% mdparagraph:P261:begin
\textbf{Weak.}
% mdparagraph:P261:end

% table: NEW-U17-ADJ-AETHELE-WEAK-DISPLAY
% mdtable:T68:begin
% course-table-style: 2026-10
% course-panel:T68-P01:begin
\par\begingroup
\fontsize{10.5}{13.2}\selectfont
\setlength{\LTleft}{0pt}\setlength{\LTright}{\fill}
\noindent
\begin{tabular}{@{}>{\raggedright\arraybackslash}p{22.0mm}>{\raggedright\arraybackslash}p{31.188mm}>{\raggedright\arraybackslash}p{31.188mm}>{\raggedright\arraybackslash}p{31.188mm}@{}}
\toprule
\textbf{Case} & \textbf{Masculine singular} & \textbf{Neuter singular} & \textbf{Feminine singular} \\
\midrule
Nominative & æðela & æðele & æðele \\
Genitive & æðelan & æðelan & æðelan \\
Dative & æðelan & æðelan & æðelan \\
Accusative & æðelan & æðele & æðelan \\
\bottomrule
\end{tabular}
\par\endgroup
% course-panel:T68-P01:end
\par\addvspace{5pt}
% course-panel:T68-P02:begin
\par\begingroup
\fontsize{10.5}{13.2}\selectfont
\setlength{\LTleft}{0pt}\setlength{\LTright}{\fill}
\noindent
\begin{tabular}{@{}>{\raggedright\arraybackslash}p{22.0mm}>{\raggedright\arraybackslash}p{31.188mm}>{\raggedright\arraybackslash}p{31.188mm}>{\raggedright\arraybackslash}p{31.188mm}@{}}
\toprule
\textbf{Case} & \textbf{Masculine plural} & \textbf{Neuter plural} & \textbf{Feminine plural} \\
\midrule
Nominative & æðelan & æðelan & æðelan \\
Genitive & æðelra & æðelra & æðelra \\
Dative & æðelum & æðelum & æðelum \\
Accusative & æðelan & æðelan & æðelan \\
\bottomrule
\end{tabular}
\par\endgroup
% course-panel:T68-P02:end
% mdtable:T68:end

\reserve{7\baselineskip}
% mdparagraph:P262:begin
\subsection*{æðelboren}
% mdparagraph:P262:end

% teaching: NEW-U17-VOC-AETHELBOREN-U17-SENSE
\par\begin{minipage}{\linewidth}\raggedright
% mdparagraph:P263:begin
\textbf{æðelboren} — high-born.
% mdparagraph:P263:end

% teaching: NEW-U17-ADJ-AETHELBOREN
% mdparagraph:P264:begin
Use the retained stem æðelboren- throughout this selected table. Do not remove a vowel merely by analogy with a different word.
% mdparagraph:P264:end

% mdparagraph:P265:begin
\textbf{Strong.}
% mdparagraph:P265:end

% table: NEW-U17-ADJ-AETHELBOREN-STRONG-DISPLAY
% mdtable:T69:begin
% course-table-style: 2026-10
% course-panel:T69-P01:begin
\par\begingroup
\fontsize{10.5}{13.2}\selectfont
\setlength{\LTleft}{0pt}\setlength{\LTright}{\fill}
\noindent
\begin{tabular}{@{}>{\raggedright\arraybackslash}p{22.0mm}>{\raggedright\arraybackslash}p{31.188mm}>{\raggedright\arraybackslash}p{31.188mm}>{\raggedright\arraybackslash}p{31.188mm}@{}}
\toprule
\textbf{Case} & \textbf{Masculine singular} & \textbf{Neuter singular} & \textbf{Feminine singular} \\
\midrule
Nominative & æðelboren & æðelboren & æðelborenu \\
Genitive & æðelborenes & æðelborenes & æðelborenre \\
Dative & æðelborenum & æðelborenum & æðelborenre \\
Accusative & æðelborenne & æðelboren & æðelborene \\
Instrumental & æðelborene & æðelborene & — \\
\bottomrule
\end{tabular}
\par\endgroup
% course-panel:T69-P01:end
\end{minipage}\par
\par\addvspace{5pt}
% course-panel:T69-P02:begin
\par\begingroup
\fontsize{10.5}{13.2}\selectfont
\setlength{\LTleft}{0pt}\setlength{\LTright}{\fill}
\noindent
\begin{tabular}{@{}>{\raggedright\arraybackslash}p{22.0mm}>{\raggedright\arraybackslash}p{31.188mm}>{\raggedright\arraybackslash}p{31.188mm}>{\raggedright\arraybackslash}p{31.188mm}@{}}
\toprule
\textbf{Case} & \textbf{Masculine plural} & \textbf{Neuter plural} & \textbf{Feminine plural} \\
\midrule
Nominative & æðelborene & æðelborenu & æðelborena \\
Genitive & æðelborenra & æðelborenra & æðelborenra \\
Dative & æðelborenum & æðelborenum & æðelborenum \\
Accusative & æðelborene & æðelborenu & æðelborena \\
Instrumental & — & — & — \\
\bottomrule
\end{tabular}
\par\endgroup
% course-panel:T69-P02:end
% mdtable:T69:end

\reserve{18\baselineskip}
\reserve{16\baselineskip}
% mdparagraph:P266:begin
\textbf{Weak.}
% mdparagraph:P266:end

% table: NEW-U17-ADJ-AETHELBOREN-WEAK-DISPLAY
% mdtable:T70:begin
% course-table-style: 2026-10
% course-panel:T70-P01:begin
\par\begingroup
\fontsize{10.5}{13.2}\selectfont
\setlength{\LTleft}{0pt}\setlength{\LTright}{\fill}
\noindent
\begin{tabular}{@{}>{\raggedright\arraybackslash}p{22.0mm}>{\raggedright\arraybackslash}p{31.188mm}>{\raggedright\arraybackslash}p{31.188mm}>{\raggedright\arraybackslash}p{31.188mm}@{}}
\toprule
\textbf{Case} & \textbf{Masculine singular} & \textbf{Neuter singular} & \textbf{Feminine singular} \\
\midrule
Nominative & æðelborena & æðelborene & æðelborene \\
Genitive & æðelborenan & æðelborenan & æðelborenan \\
Dative & æðelborenan & æðelborenan & æðelborenan \\
Accusative & æðelborenan & æðelborene & æðelborenan \\
\bottomrule
\end{tabular}
\par\endgroup
% course-panel:T70-P01:end
\par\addvspace{5pt}
% course-panel:T70-P02:begin
\par\begingroup
\fontsize{10.5}{13.2}\selectfont
\setlength{\LTleft}{0pt}\setlength{\LTright}{\fill}
\noindent
\begin{tabular}{@{}>{\raggedright\arraybackslash}p{22.0mm}>{\raggedright\arraybackslash}p{31.188mm}>{\raggedright\arraybackslash}p{31.188mm}>{\raggedright\arraybackslash}p{31.188mm}@{}}
\toprule
\textbf{Case} & \textbf{Masculine plural} & \textbf{Neuter plural} & \textbf{Feminine plural} \\
\midrule
Nominative & æðelborenan & æðelborenan & æðelborenan \\
Genitive & æðelborenra & æðelborenra & æðelborenra \\
Dative & æðelborenum & æðelborenum & æðelborenum \\
Accusative & æðelborenan & æðelborenan & æðelborenan \\
\bottomrule
\end{tabular}
\par\endgroup
% course-panel:T70-P02:end
% mdtable:T70:end

\reserve{7\baselineskip}
% mdparagraph:P267:begin
\subsection*{Adjectives naming people}
% mdparagraph:P267:end

% teaching: NEW-U17-GRAM-U17-PERSONAL-ADJECTIVE-GROUPS
% teaching: NEW-U17-VOC-MIHTIG-U17-POWERFUL-PEOPLE
% teaching: NEW-U17-VOC-EADMOD-U17-HUMBLE-PEOPLE
% mdparagraph:P268:begin
The selected adjectives geong, eald, mihtig, ēaðmōd and rīce can name young, old, powerful, humble and rich people. Use weak masculine plural forms in the assigned address and object expressions. Geongan and ealdan address people without an expressed men or article under Cook 37’s instruction. Ðā mihtgan, ðā ēaðmōdan and ðā rīcan are accusative plural objects. These people meanings are distinct from merely describing a noun; this unit assigns these particular plural uses, not an unrestricted noun declension.
% mdparagraph:P268:end

\reserve{7\baselineskip}
% mdparagraph:P269:begin
\section*{3. Verbs and nonfinite forms}
% mdparagraph:P269:end

% mdparagraph:P270:begin
Each heading states exactly which verbal system is being learned. A complete indicative table gives all three persons in both numbers for the present and preterite. A command pair gives singular and plural; a subjunctive table gives all persons in its named tense. A two-form nonfinite display gives the stated pair, rather than an unprinted finite conjugation.
% mdparagraph:P270:end

% teaching: NEW-U17-DICTIONARY-REFERENCES
% mdparagraph:P271:begin
\textbf{Dictionary references:} hęrian, praise; ondrǣdan, fear; hyngran, be hungry. These citation forms remain visible. A citation identifies the word; it does not by itself teach a new infinitival construction. The independent forms are the ones printed in the bounded systems below.
% mdparagraph:P271:end

% mdparagraph:P272:begin
Strong verbs change their stem vowels in the preterite; the second singular preterite normally shares the plural stem. Weak verbs use a dental preterite, whose spelling can merge with the last consonant of the stem. Present subjunctives generally have singular -e and plural -en; weak -ian verbs retain i in that pattern. Learn an exceptional stem together with its endings.
% mdparagraph:P272:end

\reserve{7\baselineskip}
% mdparagraph:P273:begin
\subsection*{hęrian: indicative, commands and present subjunctive}
% mdparagraph:P273:end

% teaching: NEW-U17-VOC-HERIAN-U17-SELECTED
\par\begin{minipage}{\linewidth}\raggedright
% mdparagraph:P274:begin
\textbf{hęrian} — praise.
% mdparagraph:P274:end

% teaching: u17-herian-indicative12-imperative2-present-subjunctive6
% mdparagraph:P275:begin
This weak verb follows the nęrian type. The present has hęrie, hęrest and hęreð, with -i- in plural hęriað. The preterite stem is hęred-. The command is singular hęre, plural hęriað; the present subjunctive is singular hęrie, plural hęrien. Thus hęriað can be a statement or a command, whereas hęrien expresses a plural wish in the praise exercise.
% mdparagraph:P275:end

% mdparagraph:P276:begin
\textbf{Indicative.}
% mdparagraph:P276:end

% table: u17-herian-indicative12-imperative2-present-subjunctive6-INDICATIVE-DISPLAY
% mdtable:T71:begin
% course-table-style: 2026-10
% course-panel:T71-P01:begin
\par\begingroup
\fontsize{10.5}{13.2}\selectfont
\setlength{\LTleft}{0pt}\setlength{\LTright}{\fill}
\noindent
\begin{tabular}{@{}>{\raggedright\arraybackslash}p{39.459mm}>{\raggedright\arraybackslash}p{39.459mm}>{\raggedright\arraybackslash}p{39.459mm}@{}}
\toprule
\textbf{Person and number} & \textbf{Present} & \textbf{Preterite} \\
\midrule
1 sg. & hęrie & hęrede \\
2 sg. & hęrest & hęredest \\
3 sg. & hęreð & hęrede \\
1–3 pl. & hęriað & hęredon \\
\bottomrule
\end{tabular}
\par\endgroup
% course-panel:T71-P01:end
% mdtable:T71:end
\end{minipage}\par
\reserve{9\baselineskip}
\reserve{7\baselineskip}
% mdparagraph:P277:begin
\textbf{Imperative.}
% mdparagraph:P277:end

% table: u17-herian-indicative12-imperative2-present-subjunctive6-IMPERATIVE-DISPLAY
% mdtable:T72:begin
% course-table-style: 2026-10
% course-panel:T72-P01:begin
\par\begingroup
\fontsize{10.5}{13.2}\selectfont
\setlength{\LTleft}{0pt}\setlength{\LTright}{\fill}
\noindent
\begin{tabular}{@{}>{\raggedright\arraybackslash}p{63.342mm}>{\raggedright\arraybackslash}p{57.846mm}@{}}
\toprule
\textbf{Person and number} & \textbf{Command} \\
\midrule
2 sg. & hęre \\
2 pl. & hęriað \\
\bottomrule
\end{tabular}
\par\endgroup
% course-panel:T72-P01:end
% mdtable:T72:end
\reserve{13\baselineskip}
\reserve{11\baselineskip}
% mdparagraph:P278:begin
\textbf{Subjunctive.}
% mdparagraph:P278:end

% table: u17-herian-indicative12-imperative2-present-subjunctive6-SUBJUNCTIVE-DISPLAY
% mdtable:T73:begin
% course-table-style: 2026-10
% course-panel:T73-P01:begin
\par\begingroup
\fontsize{10.5}{13.2}\selectfont
\setlength{\LTleft}{0pt}\setlength{\LTright}{\fill}
\noindent
\begin{tabular}{@{}>{\raggedright\arraybackslash}p{63.342mm}>{\raggedright\arraybackslash}p{57.846mm}@{}}
\toprule
\textbf{Person and number} & \textbf{Present} \\
\midrule
1 sg. & hęrie \\
2 sg. & hęrie \\
3 sg. & hęrie \\
1–3 pl. & hęrien \\
\bottomrule
\end{tabular}
\par\endgroup
% course-panel:T73-P01:end
% mdtable:T73:end
\reserve{7\baselineskip}
% mdparagraph:P279:begin
\subsection*{gecweðan: present and preterite indicative}
% mdparagraph:P279:end

% teaching: NEW-U17-VOC-GECWETHAN-U17-SELECTED
\par\begin{minipage}{\linewidth}\raggedright
% mdparagraph:P280:begin
\textbf{gecweðan} — speak; say.
% mdparagraph:P280:end

% teaching: u17-gecwethan-indicative12
% mdparagraph:P281:begin
The present second and third singular have i. In the preterite, first and third singular gecwæð retain æ and ð; second singular gecwǣde and plural gecwǣdon have ǣ and d. Retain the lexical prefix ge-.
% mdparagraph:P281:end

% table: u17-gecwethan-indicative12-INDICATIVE-DISPLAY
% mdtable:T74:begin
% course-table-style: 2026-10
% course-panel:T74-P01:begin
\par\begingroup
\fontsize{10.5}{13.2}\selectfont
\setlength{\LTleft}{0pt}\setlength{\LTright}{\fill}
\noindent
\begin{tabular}{@{}>{\raggedright\arraybackslash}p{39.459mm}>{\raggedright\arraybackslash}p{39.459mm}>{\raggedright\arraybackslash}p{39.459mm}@{}}
\toprule
\textbf{Person and number} & \textbf{Present} & \textbf{Preterite} \\
\midrule
1 sg. & gecweðe & gecwæð \\
2 sg. & gecwist & gecwǣde \\
3 sg. & gecwið & gecwæð \\
1–3 pl. & gecweðað & gecwǣdon \\
\bottomrule
\end{tabular}
\par\endgroup
% course-panel:T74-P01:end
% mdtable:T74:end
\end{minipage}\par
\reserve{7\baselineskip}
% mdparagraph:P282:begin
\subsection*{ūpāhębban: indicative and participles}
% mdparagraph:P282:end

% teaching: NEW-U17-VOC-UP-AHEBBAN-U17-SELECTED
\par\begin{minipage}{\linewidth}\raggedright
% mdparagraph:P283:begin
\textbf{ūp-ā-hębban} — raise; exalt.
% mdparagraph:P283:end

% teaching: u17-up-ahebban-indicative12-participle-pair2
% mdparagraph:P284:begin
The present stem has doubled bb, but the second and third singular are ūpāhęfst and ūpāhęfð. The preterite has hōf and the past participle hafen, with ūpā- retained. Raising can be physical or, in the explicitly taught exaltation sense, an increase in standing.
% mdparagraph:P284:end

% mdparagraph:P285:begin
\textbf{Indicative.}
% mdparagraph:P285:end

% table: u17-up-ahebban-indicative12-participle-pair2-INDICATIVE-DISPLAY
% mdtable:T75:begin
% course-table-style: 2026-10
% course-panel:T75-P01:begin
\par\begingroup
\fontsize{10.5}{13.2}\selectfont
\setlength{\LTleft}{0pt}\setlength{\LTright}{\fill}
\noindent
\begin{tabular}{@{}>{\raggedright\arraybackslash}p{39.459mm}>{\raggedright\arraybackslash}p{39.459mm}>{\raggedright\arraybackslash}p{39.459mm}@{}}
\toprule
\textbf{Person and number} & \textbf{Present} & \textbf{Preterite} \\
\midrule
1 sg. & ūpāhębbe & ūpāhōf \\
2 sg. & ūpāhęfst & ūpāhōfe \\
3 sg. & ūpāhęfð & ūpāhōf \\
1–3 pl. & ūpāhębbað & ūpāhōfon \\
\bottomrule
\end{tabular}
\par\endgroup
% course-panel:T75-P01:end
% mdtable:T75:end
\end{minipage}\par
\reserve{9\baselineskip}
\reserve{7\baselineskip}
% mdparagraph:P286:begin
\textbf{Participle formation pair.}
% mdparagraph:P286:end

% table: u17-up-ahebban-indicative12-participle-pair2-PARTICIPLES-DISPLAY
% mdtable:T76:begin
\par\begingroup
\fontsize{10.5}{13.2}\selectfont
\setlength{\LTleft}{0pt}\setlength{\LTright}{\fill}
\noindent
\begin{tabular}{@{}>{\raggedright\arraybackslash}p{53.15mm}>{\raggedright\arraybackslash}p{68.04mm}@{}}
\toprule
\textbf{Form} & \textbf{Literal surface} \\
\midrule
present & ūpāhębbende \\
past & ūpāhafen \\
\bottomrule
\end{tabular}
\par\endgroup

% mdtable:T76:end
\reserve{7\baselineskip}
% mdparagraph:P287:begin
\subsection*{oferwyrcean: present and preterite indicative}
% mdparagraph:P287:end

% teaching: NEW-U17-VOC-OFERWYRCEAN-U17-SELECTED
\par\begin{minipage}{\linewidth}\raggedright
% mdparagraph:P288:begin
\textbf{oferwyrcean} — overlay; cover with a worked covering.
% mdparagraph:P288:end

% teaching: u17-oferwyrcean-indicative12
% mdparagraph:P289:begin
The selected present uses -ean and plural -eað. The weak preterite has worht-. Retain ofer- throughout. The thing overlaid is accusative; the covering material follows mid with the dative.
% mdparagraph:P289:end

% table: u17-oferwyrcean-indicative12-INDICATIVE-DISPLAY
% mdtable:T77:begin
% course-table-style: 2026-10
% course-panel:T77-P01:begin
\par\begingroup
\fontsize{10.5}{13.2}\selectfont
\setlength{\LTleft}{0pt}\setlength{\LTright}{\fill}
\noindent
\begin{tabular}{@{}>{\raggedright\arraybackslash}p{39.459mm}>{\raggedright\arraybackslash}p{39.459mm}>{\raggedright\arraybackslash}p{39.459mm}@{}}
\toprule
\textbf{Person and number} & \textbf{Present} & \textbf{Preterite} \\
\midrule
1 sg. & oferwyrce & oferworhte \\
2 sg. & oferwyrcst & oferworhtest \\
3 sg. & oferwyrcð & oferworhte \\
1–3 pl. & oferwyrceað & oferworhton \\
\bottomrule
\end{tabular}
\par\endgroup
% course-panel:T77-P01:end
% mdtable:T77:end
\end{minipage}\par
\reserve{7\baselineskip}
% mdparagraph:P290:begin
\subsection*{gesęttan: present and preterite indicative}
% mdparagraph:P290:end

% teaching: NEW-U17-VOC-GESETTAN-U17-SELECTED
\par\begin{minipage}{\linewidth}\raggedright
% mdparagraph:P291:begin
\textbf{gesęttan} — set; place; establish.
% mdparagraph:P291:end

% teaching: u17-gesettan-indicative12
% mdparagraph:P292:begin
Retain tt before a vowel. The selected present second singular has -tst and third singular -tt. The weak dental preterite merges with the final t; it does not add a separate -tde sequence.
% mdparagraph:P292:end

% table: u17-gesettan-indicative12-INDICATIVE-DISPLAY
% mdtable:T78:begin
% course-table-style: 2026-10
% course-panel:T78-P01:begin
\par\begingroup
\fontsize{10.5}{13.2}\selectfont
\setlength{\LTleft}{0pt}\setlength{\LTright}{\fill}
\noindent
\begin{tabular}{@{}>{\raggedright\arraybackslash}p{39.459mm}>{\raggedright\arraybackslash}p{39.459mm}>{\raggedright\arraybackslash}p{39.459mm}@{}}
\toprule
\textbf{Person and number} & \textbf{Present} & \textbf{Preterite} \\
\midrule
1 sg. & gesętte & gesętte \\
2 sg. & gesętst & gesęttest \\
3 sg. & gesętt & gesętte \\
1–3 pl. & gesęttað & gesętton \\
\bottomrule
\end{tabular}
\par\endgroup
% course-panel:T78-P01:end
% mdtable:T78:end
\end{minipage}\par
\reserve{7\baselineskip}
% mdparagraph:P293:begin
\subsection*{giernan: present and preterite indicative}
% mdparagraph:P293:end

% teaching: NEW-U17-VOC-GIERNAN-U17-SELECTED
\par\begin{minipage}{\linewidth}\raggedright
% mdparagraph:P294:begin
\textbf{giernan} — desire; request; ask for; seek in marriage.
% mdparagraph:P294:end

% teaching: u17-giernan-indicative12
% mdparagraph:P295:begin
This long-stem weak verb has present -st and -ð and a preterite in -de. The thing desired or person sought in marriage is genitive. Æt with the dative names the person from whom something is requested.
% mdparagraph:P295:end

% table: u17-giernan-indicative12-INDICATIVE-DISPLAY
% mdtable:T79:begin
% course-table-style: 2026-10
% course-panel:T79-P01:begin
\par\begingroup
\fontsize{10.5}{13.2}\selectfont
\setlength{\LTleft}{0pt}\setlength{\LTright}{\fill}
\noindent
\begin{tabular}{@{}>{\raggedright\arraybackslash}p{39.459mm}>{\raggedright\arraybackslash}p{39.459mm}>{\raggedright\arraybackslash}p{39.459mm}@{}}
\toprule
\textbf{Person and number} & \textbf{Present} & \textbf{Preterite} \\
\midrule
1 sg. & gierne & giernde \\
2 sg. & giernst & gierndest \\
3 sg. & giernð & giernde \\
1–3 pl. & giernað & gierndon \\
\bottomrule
\end{tabular}
\par\endgroup
% course-panel:T79-P01:end
% mdtable:T79:end
\end{minipage}\par
\reserve{7\baselineskip}
% mdparagraph:P296:begin
\subsection*{mǣrsian: present and preterite indicative}
% mdparagraph:P296:end

% teaching: NEW-U17-VOC-MAERSIAN-U17-SELECTED
\par\begin{minipage}{\linewidth}\raggedright
% mdparagraph:P297:begin
\textbf{mǣrsian} — make great; magnify; glorify.
% mdparagraph:P297:end

% teaching: u17-maersian-indicative12
% mdparagraph:P298:begin
The weak -ian pattern has present singular -ie, -ast, -að and plural -iað; its preterite uses -ode, -odest, -odon. The participles have -iende and -od. In the passive, mǣrsod means celebrated or exalted.
% mdparagraph:P298:end

% table: u17-maersian-indicative12-INDICATIVE-DISPLAY
% mdtable:T80:begin
% course-table-style: 2026-10
% course-panel:T80-P01:begin
\par\begingroup
\fontsize{10.5}{13.2}\selectfont
\setlength{\LTleft}{0pt}\setlength{\LTright}{\fill}
\noindent
\begin{tabular}{@{}>{\raggedright\arraybackslash}p{39.459mm}>{\raggedright\arraybackslash}p{39.459mm}>{\raggedright\arraybackslash}p{39.459mm}@{}}
\toprule
\textbf{Person and number} & \textbf{Present} & \textbf{Preterite} \\
\midrule
1 sg. & mǣrsie & mǣrsode \\
2 sg. & mǣrsast & mǣrsodest \\
3 sg. & mǣrsað & mǣrsode \\
1–3 pl. & mǣrsiað & mǣrsodon \\
\bottomrule
\end{tabular}
\par\endgroup
% course-panel:T80-P01:end
% mdtable:T80:end
\end{minipage}\par
\reserve{7\baselineskip}
% page-review repair: U17LF01: participle heading, rules and complete two-form table
\par\begin{minipage}{\linewidth}\raggedright
% mdparagraph:P299:begin
\subsection*{mǣrsian: present and past participles}
% mdparagraph:P299:end

% teaching: u17-maersian-participle-pair2
% mdparagraph:P300:begin
The present participle has -iende and the past participle -od. The latter is used in the passive description celebrated or exalted.
% mdparagraph:P300:end

% mdparagraph:P301:begin
This display assigns the complete present/\allowbreak{}past formation pair. It does not add other finite moods or an unrestricted attributive declension; the past form is used in the stated passive frame.
% mdparagraph:P301:end

% table: u17-maersian-participle-pair2-PARTICIPLES-DISPLAY
% mdtable:T81:begin
\par\begingroup
\fontsize{10.5}{13.2}\selectfont
\setlength{\LTleft}{0pt}\setlength{\LTright}{\fill}
\noindent
\begin{tabular}{@{}>{\raggedright\arraybackslash}p{55.85mm}>{\raggedright\arraybackslash}p{65.34mm}@{}}
\toprule
\textbf{Form} & \textbf{Literal surface} \\
\midrule
present & mǣrsiende \\
past & mǣrsod \\
\bottomrule
\end{tabular}
\par\endgroup

% mdtable:T81:end
\end{minipage}\par

\reserve{7\baselineskip}
% mdparagraph:P302:begin
\subsection*{gewilnian: present and preterite indicative}
% mdparagraph:P302:end

% teaching: NEW-U17-VOC-GEWILNIAN-U17-SELECTED
\par\begin{minipage}{\linewidth}\raggedright
% mdparagraph:P303:begin
\textbf{gewilnian} — desire; wish.
% mdparagraph:P303:end

% teaching: u17-gewilnian-indicative12
% mdparagraph:P304:begin
Use the weak -ian endings: present -ie, -ast, -að, -iað; preterite -ode, -odest, -odon. A following ðæt clause can express what is desired.
% mdparagraph:P304:end

% table: u17-gewilnian-indicative12-INDICATIVE-DISPLAY
% mdtable:T82:begin
% course-table-style: 2026-10
% course-panel:T82-P01:begin
\par\begingroup
\fontsize{10.5}{13.2}\selectfont
\setlength{\LTleft}{0pt}\setlength{\LTright}{\fill}
\noindent
\begin{tabular}{@{}>{\raggedright\arraybackslash}p{39.459mm}>{\raggedright\arraybackslash}p{39.459mm}>{\raggedright\arraybackslash}p{39.459mm}@{}}
\toprule
\textbf{Person and number} & \textbf{Present} & \textbf{Preterite} \\
\midrule
1 sg. & gewilnie & gewilnode \\
2 sg. & gewilnast & gewilnodest \\
3 sg. & gewilnað & gewilnode \\
1–3 pl. & gewilniað & gewilnodon \\
\bottomrule
\end{tabular}
\par\endgroup
% course-panel:T82-P01:end
% mdtable:T82:end
\end{minipage}\par
\reserve{7\baselineskip}
% mdparagraph:P305:begin
\subsection*{ðurhwunian: present and preterite indicative}
% mdparagraph:P305:end

% teaching: NEW-U17-VOC-THURHWUNIAN-U17-SELECTED
\par\begin{minipage}{\linewidth}\raggedright
% mdparagraph:P306:begin
\textbf{ðurhwunian} — remain; endure; last.
% mdparagraph:P306:end

% teaching: u17-thurhwunian-indicative12
% mdparagraph:P307:begin
Use the weak -ian endings: present -ie, -ast, -að, -iað; preterite -ode, -odest, -odon. The meaning is remaining or lasting.
% mdparagraph:P307:end

% table: u17-thurhwunian-indicative12-INDICATIVE-DISPLAY
% mdtable:T83:begin
% course-table-style: 2026-10
% course-panel:T83-P01:begin
\par\begingroup
\fontsize{10.5}{13.2}\selectfont
\setlength{\LTleft}{0pt}\setlength{\LTright}{\fill}
\noindent
\begin{tabular}{@{}>{\raggedright\arraybackslash}p{39.459mm}>{\raggedright\arraybackslash}p{39.459mm}>{\raggedright\arraybackslash}p{39.459mm}@{}}
\toprule
\textbf{Person and number} & \textbf{Present} & \textbf{Preterite} \\
\midrule
1 sg. & ðurhwunie & ðurhwunode \\
2 sg. & ðurhwunast & ðurhwunodest \\
3 sg. & ðurhwunað & ðurhwunode \\
1–3 pl. & ðurhwuniað & ðurhwunodon \\
\bottomrule
\end{tabular}
\par\endgroup
% course-panel:T83-P01:end
% mdtable:T83:end
\end{minipage}\par
\reserve{7\baselineskip}
% mdparagraph:P308:begin
\subsection*{wiðstandan: present and preterite indicative}
% mdparagraph:P308:end

% teaching: NEW-U17-VOC-WITHSTANDAN-U17-SELECTED
\par\begin{minipage}{\linewidth}\raggedright
% mdparagraph:P309:begin
\textbf{wiðstandan} — withstand; resist.
% mdparagraph:P309:end

% teaching: u17-withstandan-indicative12
% mdparagraph:P310:begin
The present second and third singular have stęntst and stęnt. The preterite has stōd and stōdon, without the present stem’s n. The thing resisted is dative.
% mdparagraph:P310:end

% table: u17-withstandan-indicative12-INDICATIVE-DISPLAY
% mdtable:T84:begin
% course-table-style: 2026-10
% course-panel:T84-P01:begin
\par\begingroup
\fontsize{10.5}{13.2}\selectfont
\setlength{\LTleft}{0pt}\setlength{\LTright}{\fill}
\noindent
\begin{tabular}{@{}>{\raggedright\arraybackslash}p{39.459mm}>{\raggedright\arraybackslash}p{39.459mm}>{\raggedright\arraybackslash}p{39.459mm}@{}}
\toprule
\textbf{Person and number} & \textbf{Present} & \textbf{Preterite} \\
\midrule
1 sg. & wiðstande & wiðstōd \\
2 sg. & wiðstęntst & wiðstōde \\
3 sg. & wiðstęnt & wiðstōd \\
1–3 pl. & wiðstandað & wiðstōdon \\
\bottomrule
\end{tabular}
\par\endgroup
% course-panel:T84-P01:end
% mdtable:T84:end
\end{minipage}\par
\reserve{7\baselineskip}
% mdparagraph:P311:begin
\subsection*{gearcian: present and preterite indicative}
% mdparagraph:P311:end

% teaching: NEW-U17-VOC-GEARCIAN-U17-SELECTED
\par\begin{minipage}{\linewidth}\raggedright
% mdparagraph:P312:begin
\textbf{gearcian} — prepare.
% mdparagraph:P312:end

% teaching: u17-gearcian-indicative12
% mdparagraph:P313:begin
The verb is unprefixed gearcian. Use the weak -ian endings: present -ie, -ast, -að, -iað; preterite -ode, -odest, -odon. The things prepared are accusative and their beneficiaries dative.
% mdparagraph:P313:end

% table: u17-gearcian-indicative12-INDICATIVE-DISPLAY
% mdtable:T85:begin
% course-table-style: 2026-10
% course-panel:T85-P01:begin
\par\begingroup
\fontsize{10.5}{13.2}\selectfont
\setlength{\LTleft}{0pt}\setlength{\LTright}{\fill}
\noindent
\begin{tabular}{@{}>{\raggedright\arraybackslash}p{39.459mm}>{\raggedright\arraybackslash}p{39.459mm}>{\raggedright\arraybackslash}p{39.459mm}@{}}
\toprule
\textbf{Person and number} & \textbf{Present} & \textbf{Preterite} \\
\midrule
1 sg. & gearcie & gearcode \\
2 sg. & gearcast & gearcodest \\
3 sg. & gearcað & gearcode \\
1–3 pl. & gearciað & gearcodon \\
\bottomrule
\end{tabular}
\par\endgroup
% course-panel:T85-P01:end
% mdtable:T85:end
\end{minipage}\par
\reserve{7\baselineskip}
% mdparagraph:P314:begin
\subsection*{āfyllan: indicative and participles}
% mdparagraph:P314:end

% teaching: NEW-U17-VOC-AFYLLAN-U17-SELECTED
\par\begin{minipage}{\linewidth}\raggedright
% mdparagraph:P315:begin
\textbf{āfyllan} — fill.
% mdparagraph:P315:end

% teaching: u17-afyllan-indicative12-participle-pair2
% mdparagraph:P316:begin
Double ll before a vowel, but use single l in āfylst, āfylð and āfylde. The selected past participle is āfylled; contracted āfyld is also recorded. The filled entity is accusative in an active clause and nominative in a passive; mid with the dative names the contents.
% mdparagraph:P316:end

% mdparagraph:P317:begin
\textbf{Indicative.}
% mdparagraph:P317:end

% table: u17-afyllan-indicative12-participle-pair2-INDICATIVE-DISPLAY
% mdtable:T86:begin
% course-table-style: 2026-10
% course-panel:T86-P01:begin
\par\begingroup
\fontsize{10.5}{13.2}\selectfont
\setlength{\LTleft}{0pt}\setlength{\LTright}{\fill}
\noindent
\begin{tabular}{@{}>{\raggedright\arraybackslash}p{39.459mm}>{\raggedright\arraybackslash}p{39.459mm}>{\raggedright\arraybackslash}p{39.459mm}@{}}
\toprule
\textbf{Person and number} & \textbf{Present} & \textbf{Preterite} \\
\midrule
1 sg. & āfylle & āfylde \\
2 sg. & āfylst & āfyldest \\
3 sg. & āfylð & āfylde \\
1–3 pl. & āfyllað & āfyldon \\
\bottomrule
\end{tabular}
\par\endgroup
% course-panel:T86-P01:end
% mdtable:T86:end
\end{minipage}\par
\reserve{9\baselineskip}
\reserve{7\baselineskip}
% mdparagraph:P318:begin
\textbf{Participle formation pair.}
% mdparagraph:P318:end

% table: u17-afyllan-indicative12-participle-pair2-PARTICIPLES-DISPLAY
% mdtable:T87:begin
\par\begingroup
\fontsize{10.5}{13.2}\selectfont
\setlength{\LTleft}{0pt}\setlength{\LTright}{\fill}
\noindent
\begin{tabular}{@{}>{\raggedright\arraybackslash}p{55.85mm}>{\raggedright\arraybackslash}p{65.34mm}@{}}
\toprule
\textbf{Form} & \textbf{Literal surface} \\
\midrule
present & āfyllende \\
past & āfylled \\
\bottomrule
\end{tabular}
\par\endgroup

% mdtable:T87:end
\reserve{7\baselineskip}
% mdparagraph:P319:begin
\subsection*{āweorpan: present and preterite indicative}
% mdparagraph:P319:end

% teaching: NEW-U17-VOC-AWEORPAN-U17-SELECTED
\par\begin{minipage}{\linewidth}\raggedright
% mdparagraph:P320:begin
\textbf{āweorpan} — cast down; cast away.
% mdparagraph:P320:end

% teaching: u17-aweorpan-indicative12
% mdparagraph:P321:begin
The strong stems alternate eo, ea, u, o. Present second and third singular have ie; the preterite second singular uses the same u-stem as the plural. The displaced person is accusative and of with the dative gives the source.
% mdparagraph:P321:end

% table: u17-aweorpan-indicative12-INDICATIVE-DISPLAY
% mdtable:T88:begin
% course-table-style: 2026-10
% course-panel:T88-P01:begin
\par\begingroup
\fontsize{10.5}{13.2}\selectfont
\setlength{\LTleft}{0pt}\setlength{\LTright}{\fill}
\noindent
\begin{tabular}{@{}>{\raggedright\arraybackslash}p{39.459mm}>{\raggedright\arraybackslash}p{39.459mm}>{\raggedright\arraybackslash}p{39.459mm}@{}}
\toprule
\textbf{Person and number} & \textbf{Present} & \textbf{Preterite} \\
\midrule
1 sg. & āweorpe & āwearp \\
2 sg. & āwierpst & āwurpe \\
3 sg. & āwierpð & āwearp \\
1–3 pl. & āweorpað & āwurpon \\
\bottomrule
\end{tabular}
\par\endgroup
% course-panel:T88-P01:end
% mdtable:T88:end
\end{minipage}\par
\reserve{7\baselineskip}
% mdparagraph:P322:begin
\subsection*{āwrītan: present and preterite indicative}
% mdparagraph:P322:end

% teaching: NEW-U17-VOC-AWRITAN-U17-SELECTED
\par\begin{minipage}{\linewidth}\raggedright
% mdparagraph:P323:begin
\textbf{āwrītan} — write.
% mdparagraph:P323:end

% teaching: u17-awritan-indicative12
% mdparagraph:P324:begin
The strong vowel sequence is ī, ā, i, i. Present second singular has -tst and the selected third singular -tt; a source variant has single -t. There is no new vowel mutation in these two present forms.
% mdparagraph:P324:end

% table: u17-awritan-indicative12-INDICATIVE-DISPLAY
% mdtable:T89:begin
% course-table-style: 2026-10
% course-panel:T89-P01:begin
\par\begingroup
\fontsize{10.5}{13.2}\selectfont
\setlength{\LTleft}{0pt}\setlength{\LTright}{\fill}
\noindent
\begin{tabular}{@{}>{\raggedright\arraybackslash}p{39.459mm}>{\raggedright\arraybackslash}p{39.459mm}>{\raggedright\arraybackslash}p{39.459mm}@{}}
\toprule
\textbf{Person and number} & \textbf{Present} & \textbf{Preterite} \\
\midrule
1 sg. & āwrīte & āwrāt \\
2 sg. & āwrītst & āwrite \\
3 sg. & āwrītt & āwrāt \\
1–3 pl. & āwrītað & āwriton \\
\bottomrule
\end{tabular}
\par\endgroup
% course-panel:T89-P01:end
% mdtable:T89:end
\end{minipage}\par
\reserve{7\baselineskip}
% mdparagraph:P325:begin
\subsection*{ālūtan: present and preterite indicative}
% mdparagraph:P325:end

% teaching: NEW-U17-VOC-ALUTAN-U17-SELECTED
\par\begin{minipage}{\linewidth}\raggedright
% mdparagraph:P326:begin
\textbf{ālūtan} — bow; incline oneself.
% mdparagraph:P326:end

% teaching: u17-alutan-indicative12
% mdparagraph:P327:begin
The strong stems alternate ū, ēa, u, o. Present second and third singular change ū to ȳ and have -tst and selected -tt; single -t is also recorded. The actor bows; dative Gode names the one bowed to.
% mdparagraph:P327:end

% table: u17-alutan-indicative12-INDICATIVE-DISPLAY
% mdtable:T90:begin
% course-table-style: 2026-10
% course-panel:T90-P01:begin
\par\begingroup
\fontsize{10.5}{13.2}\selectfont
\setlength{\LTleft}{0pt}\setlength{\LTright}{\fill}
\noindent
\begin{tabular}{@{}>{\raggedright\arraybackslash}p{39.459mm}>{\raggedright\arraybackslash}p{39.459mm}>{\raggedright\arraybackslash}p{39.459mm}@{}}
\toprule
\textbf{Person and number} & \textbf{Present} & \textbf{Preterite} \\
\midrule
1 sg. & ālūte & ālēat \\
2 sg. & ālȳtst & ālute \\
3 sg. & ālȳtt & ālēat \\
1–3 pl. & ālūtað & āluton \\
\bottomrule
\end{tabular}
\par\endgroup
% course-panel:T90-P01:end
% mdtable:T90:end
\end{minipage}\par
\reserve{7\baselineskip}
% mdparagraph:P328:begin
\subsection*{windan: present and preterite indicative}
% mdparagraph:P328:end

% teaching: NEW-U17-VOC-WINDAN-U17-SELECTED
\par\begin{minipage}{\linewidth}\raggedright
% mdparagraph:P329:begin
\textbf{windan} — wind; plait.
% mdparagraph:P329:end

% teaching: u17-windan-indicative12
% mdparagraph:P330:begin
The strong stems alternate i, a, u, u. Stem d becomes t before the second- and third-singular present endings. The thing plaited is accusative; of with the dative names its material.
% mdparagraph:P330:end

% table: u17-windan-indicative12-INDICATIVE-DISPLAY
% mdtable:T91:begin
% course-table-style: 2026-10
% course-panel:T91-P01:begin
\par\begingroup
\fontsize{10.5}{13.2}\selectfont
\setlength{\LTleft}{0pt}\setlength{\LTright}{\fill}
\noindent
\begin{tabular}{@{}>{\raggedright\arraybackslash}p{39.459mm}>{\raggedright\arraybackslash}p{39.459mm}>{\raggedright\arraybackslash}p{39.459mm}@{}}
\toprule
\textbf{Person and number} & \textbf{Present} & \textbf{Preterite} \\
\midrule
1 sg. & winde & wand \\
2 sg. & wintst & wunde \\
3 sg. & wint & wand \\
1–3 pl. & windað & wundon \\
\bottomrule
\end{tabular}
\par\endgroup
% course-panel:T91-P01:end
% mdtable:T91:end
\end{minipage}\par
\reserve{7\baselineskip}
% mdparagraph:P331:begin
\subsection*{spǣtan: present and preterite indicative}
% mdparagraph:P331:end

% teaching: NEW-U17-VOC-SPAETAN-U17-SELECTED
\par\begin{minipage}{\linewidth}\raggedright
% mdparagraph:P332:begin
\textbf{spǣtan} — spit.
% mdparagraph:P332:end

% teaching: u17-spaetan-indicative12
% mdparagraph:P333:begin
Retain long ǣ. The t-stem has selected present third singular -tt and weak preterite -tte, -ttest, -tton. In the exercise, on takes the accusative person spat upon.
% mdparagraph:P333:end

% table: u17-spaetan-indicative12-INDICATIVE-DISPLAY
% mdtable:T92:begin
% course-table-style: 2026-10
% course-panel:T92-P01:begin
\par\begingroup
\fontsize{10.5}{13.2}\selectfont
\setlength{\LTleft}{0pt}\setlength{\LTright}{\fill}
\noindent
\begin{tabular}{@{}>{\raggedright\arraybackslash}p{39.459mm}>{\raggedright\arraybackslash}p{39.459mm}>{\raggedright\arraybackslash}p{39.459mm}@{}}
\toprule
\textbf{Person and number} & \textbf{Present} & \textbf{Preterite} \\
\midrule
1 sg. & spǣte & spǣtte \\
2 sg. & spǣtst & spǣttest \\
3 sg. & spǣtt & spǣtte \\
1–3 pl. & spǣtað & spǣtton \\
\bottomrule
\end{tabular}
\par\endgroup
% course-panel:T92-P01:end
% mdtable:T92:end
\end{minipage}\par
\reserve{7\baselineskip}
% mdparagraph:P334:begin
\subsection*{tǣlan: present and preterite indicative}
% mdparagraph:P334:end

% teaching: NEW-U17-VOC-TAELAN-U17-SELECTED
\par\begin{minipage}{\linewidth}\raggedright
% mdparagraph:P335:begin
\textbf{tǣlan} — blame; censure.
% mdparagraph:P335:end

% teaching: u17-taelan-indicative12
% mdparagraph:P336:begin
This long-stem weak verb has a preterite in -de. The person blamed is accusative. In the reading, tǣlst is second singular present indicative.
% mdparagraph:P336:end

% table: u17-taelan-indicative12-INDICATIVE-DISPLAY
% mdtable:T93:begin
% course-table-style: 2026-10
% course-panel:T93-P01:begin
\par\begingroup
\fontsize{10.5}{13.2}\selectfont
\setlength{\LTleft}{0pt}\setlength{\LTright}{\fill}
\noindent
\begin{tabular}{@{}>{\raggedright\arraybackslash}p{39.459mm}>{\raggedright\arraybackslash}p{39.459mm}>{\raggedright\arraybackslash}p{39.459mm}@{}}
\toprule
\textbf{Person and number} & \textbf{Present} & \textbf{Preterite} \\
\midrule
1 sg. & tǣle & tǣlde \\
2 sg. & tǣlst & tǣldest \\
3 sg. & tǣlð & tǣlde \\
1–3 pl. & tǣlað & tǣldon \\
\bottomrule
\end{tabular}
\par\endgroup
% course-panel:T93-P01:end
% mdtable:T93:end
\end{minipage}\par
\reserve{7\baselineskip}
% mdparagraph:P337:begin
\subsection*{bīegan: present and preterite indicative}
% mdparagraph:P337:end

% teaching: NEW-U17-VOC-BIEGAN-U17-SELECTED
\par\begin{minipage}{\linewidth}\raggedright
% mdparagraph:P338:begin
\textbf{bīegan} — bend.
% mdparagraph:P338:end

% teaching: u17-biegan-indicative12
% mdparagraph:P339:begin
Use the selected long-diphthong stem bīeg- with weak endings. This transitive verb takes the bent knee as an accusative object. Keep it distinct from the intransitive bowing verb.
% mdparagraph:P339:end

% table: u17-biegan-indicative12-INDICATIVE-DISPLAY
% mdtable:T94:begin
% course-table-style: 2026-10
% course-panel:T94-P01:begin
\par\begingroup
\fontsize{10.5}{13.2}\selectfont
\setlength{\LTleft}{0pt}\setlength{\LTright}{\fill}
\noindent
\begin{tabular}{@{}>{\raggedright\arraybackslash}p{39.459mm}>{\raggedright\arraybackslash}p{39.459mm}>{\raggedright\arraybackslash}p{39.459mm}@{}}
\toprule
\textbf{Person and number} & \textbf{Present} & \textbf{Preterite} \\
\midrule
1 sg. & bīege & bīegde \\
2 sg. & bīegst & bīegdest \\
3 sg. & bīegð & bīegde \\
1–3 pl. & bīegað & bīegdon \\
\bottomrule
\end{tabular}
\par\endgroup
% course-panel:T94-P01:end
% mdtable:T94:end
\end{minipage}\par
\reserve{7\baselineskip}
% mdparagraph:P340:begin
\subsection*{gemunan: present and preterite indicative}
% mdparagraph:P340:end

% teaching: NEW-U17-VOC-GEMUNAN-U17-SELECTED
\par\begin{minipage}{\linewidth}\raggedright
% mdparagraph:P341:begin
\textbf{gemunan} — remember.
% mdparagraph:P341:end

% teaching: u17-gemunan-indicative12
% mdparagraph:P342:begin
This preterite-present verb has present first/\allowbreak{}third singular geman, second singular gemanst and plural gemunon. Its dental preterite uses gemund-. The thing remembered is accusative.
% mdparagraph:P342:end

% table: u17-gemunan-indicative12-INDICATIVE-DISPLAY
% mdtable:T95:begin
% course-table-style: 2026-10
% course-panel:T95-P01:begin
\par\begingroup
\fontsize{10.5}{13.2}\selectfont
\setlength{\LTleft}{0pt}\setlength{\LTright}{\fill}
\noindent
\begin{tabular}{@{}>{\raggedright\arraybackslash}p{39.459mm}>{\raggedright\arraybackslash}p{39.459mm}>{\raggedright\arraybackslash}p{39.459mm}@{}}
\toprule
\textbf{Person and number} & \textbf{Present} & \textbf{Preterite} \\
\midrule
1 sg. & geman & gemunde \\
2 sg. & gemanst & gemundest \\
3 sg. & geman & gemunde \\
1–3 pl. & gemunon & gemundon \\
\bottomrule
\end{tabular}
\par\endgroup
% course-panel:T95-P01:end
% mdtable:T95:end
\end{minipage}\par
\reserve{7\baselineskip}
% mdparagraph:P343:begin
\subsection*{ðęncean: present and preterite indicative}
% mdparagraph:P343:end

% teaching: NEW-U17-VOC-THENCAN-U17-SELECTED
\par\begin{minipage}{\linewidth}\raggedright
% mdparagraph:P344:begin
\textbf{ðęncan} — think.
% mdparagraph:P344:end

% teaching: u17-thencan-indicative12
% mdparagraph:P345:begin
The selected headword variant is ðęncean. Its weak preterite ðōht- changes the vowel and replaces c before the dental ending. In the Sheba exchange, a free-choice expression gives what the queen had thought.
% mdparagraph:P345:end

% table: u17-thencan-indicative12-INDICATIVE-DISPLAY
% mdtable:T96:begin
% course-table-style: 2026-10
% course-panel:T96-P01:begin
\par\begingroup
\fontsize{10.5}{13.2}\selectfont
\setlength{\LTleft}{0pt}\setlength{\LTright}{\fill}
\noindent
\begin{tabular}{@{}>{\raggedright\arraybackslash}p{39.459mm}>{\raggedright\arraybackslash}p{39.459mm}>{\raggedright\arraybackslash}p{39.459mm}@{}}
\toprule
\textbf{Person and number} & \textbf{Present} & \textbf{Preterite} \\
\midrule
1 sg. & ðęnce & ðōhte \\
2 sg. & ðęncst & ðōhtest \\
3 sg. & ðęncð & ðōhte \\
1–3 pl. & ðęnceað & ðōhton \\
\bottomrule
\end{tabular}
\par\endgroup
% course-panel:T96-P01:end
% mdtable:T96:end
\end{minipage}\par
\reserve{7\baselineskip}
% mdparagraph:P346:begin
\subsection*{blētsian: present and past participles}
% mdparagraph:P346:end

% teaching: NEW-U17-VOC-BLETSIAN-U17-BLESS
\par\begin{minipage}{\linewidth}\raggedright
% mdparagraph:P347:begin
\textbf{blētsian} — bless.
% mdparagraph:P347:end

% teaching: u17-bletsian-participle-pair2
% mdparagraph:P348:begin
The present participle has -iende and the past participle -od in this weak -ian pattern.
% mdparagraph:P348:end

% mdparagraph:P349:begin
This display assigns the complete present/\allowbreak{}past formation pair. It does not add other finite moods or an unrestricted attributive declension; the past form is used in the stated passive frame.
% mdparagraph:P349:end

% table: u17-bletsian-participle-pair2-PARTICIPLES-DISPLAY
% mdtable:T97:begin
\par\begingroup
\fontsize{10.5}{13.2}\selectfont
\setlength{\LTleft}{0pt}\setlength{\LTright}{\fill}
\noindent
\begin{tabular}{@{}>{\raggedright\arraybackslash}p{54.47mm}>{\raggedright\arraybackslash}p{66.72mm}@{}}
\toprule
\textbf{Form} & \textbf{Literal surface} \\
\midrule
present & blētsiende \\
past & blētsod \\
\bottomrule
\end{tabular}
\par\endgroup

% mdtable:T97:end
\end{minipage}\par
\reserve{7\baselineskip}
% mdparagraph:P350:begin
\subsection*{ofwundrian: present and past participles}
% mdparagraph:P350:end

% teaching: NEW-U17-VOC-OFWUNDRIAN-U17-ASTONISHED
\par\begin{minipage}{\linewidth}\raggedright
% mdparagraph:P351:begin
\textbf{ofwundrian} — be astonished.
% mdparagraph:P351:end

% teaching: u17-ofwundrian-participle-pair2
% mdparagraph:P352:begin
The present participle has -iende and the past participle -od in this weak -ian pattern.
% mdparagraph:P352:end

% mdparagraph:P353:begin
This display assigns the complete present/\allowbreak{}past formation pair. It does not add other finite moods or an unrestricted attributive declension; the past form is used in the stated passive frame.
% mdparagraph:P353:end

% table: u17-ofwundrian-participle-pair2-PARTICIPLES-DISPLAY
% mdtable:T98:begin
\par\begingroup
\fontsize{10.5}{13.2}\selectfont
\setlength{\LTleft}{0pt}\setlength{\LTright}{\fill}
\noindent
\begin{tabular}{@{}>{\raggedright\arraybackslash}p{51.9mm}>{\raggedright\arraybackslash}p{69.29mm}@{}}
\toprule
\textbf{Form} & \textbf{Literal surface} \\
\midrule
present & ofwundriende \\
past & ofwundrod \\
\bottomrule
\end{tabular}
\par\endgroup

% mdtable:T98:end
\end{minipage}\par
\reserve{7\baselineskip}
% mdparagraph:P354:begin
\subsection*{wyrcean: present and past participles}
% mdparagraph:P354:end

% teaching: u17-wyrcean-participle-pair2
% mdparagraph:P355:begin
The present participle uses the present stem; the past participle is worht, with the same changed stem as the weak preterite.
% mdparagraph:P355:end

\reserve{10\baselineskip}
% mdparagraph:P356:begin
This display assigns the complete present/\allowbreak{}past formation pair. It does not add other finite moods or an unrestricted attributive declension; the past form is used in the stated passive frame.
% mdparagraph:P356:end

% table: u17-wyrcean-participle-pair2-PARTICIPLES-DISPLAY
% mdtable:T99:begin
\par\begingroup
\fontsize{10.5}{13.2}\selectfont
\setlength{\LTleft}{0pt}\setlength{\LTright}{\fill}
\noindent
\begin{tabular}{@{}>{\raggedright\arraybackslash}p{57.3mm}>{\raggedright\arraybackslash}p{63.89mm}@{}}
\toprule
\textbf{Form} & \textbf{Literal surface} \\
\midrule
present & wyrcende \\
past & worht \\
\bottomrule
\end{tabular}
\par\endgroup

% mdtable:T99:end
\reserve{7\baselineskip}
% mdparagraph:P357:begin
\subsection*{gesweotolian: present and past participles}
% mdparagraph:P357:end

% teaching: NEW-U17-VOC-GESWEOTOLIAN-U17-MANIFEST
\par\begin{minipage}{\linewidth}\raggedright
% mdparagraph:P358:begin
\textbf{gesweotolian} — make manifest; reveal.
% mdparagraph:P358:end

% teaching: u17-gesweotolian-participle-pair2
% mdparagraph:P359:begin
The present participle has -iende and the past participle -od in this weak -ian pattern.
% mdparagraph:P359:end

% mdparagraph:P360:begin
This display assigns the complete present/\allowbreak{}past formation pair. It does not add other finite moods or an unrestricted attributive declension; the past form is used in the stated passive frame.
% mdparagraph:P360:end

% table: u17-gesweotolian-participle-pair2-PARTICIPLES-DISPLAY
% mdtable:T100:begin
\par\begingroup
\fontsize{10.5}{13.2}\selectfont
\setlength{\LTleft}{0pt}\setlength{\LTright}{\fill}
\noindent
\begin{tabular}{@{}>{\raggedright\arraybackslash}p{49.57mm}>{\raggedright\arraybackslash}p{71.62mm}@{}}
\toprule
\textbf{Form} & \textbf{Literal surface} \\
\midrule
present & gesweotoliende \\
past & gesweotolod \\
\bottomrule
\end{tabular}
\par\endgroup

% mdtable:T100:end
\end{minipage}\par
\reserve{7\baselineskip}
% mdparagraph:P361:begin
\subsection*{ācęnnan: present and past participles}
% mdparagraph:P361:end

% teaching: NEW-U17-VOC-ACENNAN-U17-BEGET
\par\begin{minipage}{\linewidth}\raggedright
% mdparagraph:P362:begin
\textbf{ācęnnan} — beget; bring forth.
% mdparagraph:P362:end

% teaching: u17-acennan-participle-pair2
% mdparagraph:P363:begin
The present participle keeps the present stem with nn; the past participle has -ed. The begetting expression below is supplied in full.
% mdparagraph:P363:end

% mdparagraph:P364:begin
This display assigns the complete present/\allowbreak{}past formation pair. It does not add other finite moods or an unrestricted attributive declension; the past form is used in the stated passive frame.
% mdparagraph:P364:end

% table: u17-acennan-participle-pair2-PARTICIPLES-DISPLAY
% mdtable:T101:begin
\par\begingroup
\fontsize{10.5}{13.2}\selectfont
\setlength{\LTleft}{0pt}\setlength{\LTright}{\fill}
\noindent
\begin{tabular}{@{}>{\raggedright\arraybackslash}p{55.85mm}>{\raggedright\arraybackslash}p{65.34mm}@{}}
\toprule
\textbf{Form} & \textbf{Literal surface} \\
\midrule
present & ācęnnende \\
past & ācęnned \\
\bottomrule
\end{tabular}
\par\endgroup

% mdtable:T101:end
\end{minipage}\par
\reserve{7\baselineskip}
% mdparagraph:P365:begin
\subsection*{geopenian: present and past participles}
% mdparagraph:P365:end

% teaching: u17-geopenian-participle-pair2
% mdparagraph:P366:begin
The present participle has -iende and the past participle -od in this weak -ian pattern.
% mdparagraph:P366:end

\reserve{10\baselineskip}
% mdparagraph:P367:begin
This display assigns the complete present/\allowbreak{}past formation pair. It does not add other finite moods or an unrestricted attributive declension; the past form is used in the stated passive frame.
% mdparagraph:P367:end

% table: u17-geopenian-participle-pair2-PARTICIPLES-DISPLAY
% mdtable:T102:begin
\par\begingroup
\fontsize{10.5}{13.2}\selectfont
\setlength{\LTleft}{0pt}\setlength{\LTright}{\fill}
\noindent
\begin{tabular}{@{}>{\raggedright\arraybackslash}p{53.15mm}>{\raggedright\arraybackslash}p{68.04mm}@{}}
\toprule
\textbf{Form} & \textbf{Literal surface} \\
\midrule
present & geopeniende \\
past & geopenod \\
\bottomrule
\end{tabular}
\par\endgroup

% mdtable:T102:end
\reserve{7\baselineskip}
% mdparagraph:P368:begin
\subsection*{cnucian: singular and plural commands}
% mdparagraph:P368:end

% teaching: NEW-U17-VOC-CNUCIAN-U17-KNOCK
\par\begin{minipage}{\linewidth}\raggedright
% mdparagraph:P369:begin
\textbf{cnucian} — knock.
% mdparagraph:P369:end

% teaching: u17-cnucian-imperative2
% mdparagraph:P370:begin
The weak -ian command has singular -a and plural -iað. Only this complete command pair is added here.
% mdparagraph:P370:end

% table: u17-cnucian-imperative2-IMPERATIVE-DISPLAY
% mdtable:T103:begin
% course-table-style: 2026-10
% course-panel:T103-P01:begin
\par\begingroup
\fontsize{10.5}{13.2}\selectfont
\setlength{\LTleft}{0pt}\setlength{\LTright}{\fill}
\noindent
\begin{tabular}{@{}>{\raggedright\arraybackslash}p{63.342mm}>{\raggedright\arraybackslash}p{57.846mm}@{}}
\toprule
\textbf{Person and number} & \textbf{Command} \\
\midrule
2 sg. & cnuca \\
2 pl. & cnuciað \\
\bottomrule
\end{tabular}
\par\endgroup
% course-panel:T103-P01:end
% mdtable:T103:end
\end{minipage}\par
\reserve{18\baselineskip}
% mdparagraph:P371:begin
\subsection*{āhębban: singular and plural commands}
% mdparagraph:P371:end

% teaching: u17-ahebban-imperative2
% mdparagraph:P372:begin
The strong j-present has singular āhęfe and plural āhębbað. Only this complete command pair is added here.
% mdparagraph:P372:end

% table: u17-ahebban-imperative2-IMPERATIVE-DISPLAY
% mdtable:T104:begin
% course-table-style: 2026-10
% course-panel:T104-P01:begin
\par\begingroup
\fontsize{10.5}{13.2}\selectfont
\setlength{\LTleft}{0pt}\setlength{\LTright}{\fill}
\noindent
\begin{tabular}{@{}>{\raggedright\arraybackslash}p{63.342mm}>{\raggedright\arraybackslash}p{57.846mm}@{}}
\toprule
\textbf{Person and number} & \textbf{Command} \\
\midrule
2 sg. & āhęfe \\
2 pl. & āhębbað \\
\bottomrule
\end{tabular}
\par\endgroup
% course-panel:T104-P01:end
% mdtable:T104:end
\reserve{18\baselineskip}
% mdparagraph:P373:begin
\subsection*{plegian: singular and plural commands}
% mdparagraph:P373:end

% teaching: u17-plegian-imperative2
% mdparagraph:P374:begin
The weak -ian command has singular -a and plural -iað. The assigned ordinary sense is play, including playing with a ball. Only this complete command pair is added here.
% mdparagraph:P374:end

% table: u17-plegian-imperative2-IMPERATIVE-DISPLAY
% mdtable:T105:begin
% course-table-style: 2026-10
% course-panel:T105-P01:begin
\par\begingroup
\fontsize{10.5}{13.2}\selectfont
\setlength{\LTleft}{0pt}\setlength{\LTright}{\fill}
\noindent
\begin{tabular}{@{}>{\raggedright\arraybackslash}p{63.342mm}>{\raggedright\arraybackslash}p{57.846mm}@{}}
\toprule
\textbf{Person and number} & \textbf{Command} \\
\midrule
2 sg. & plega \\
2 pl. & plegiað \\
\bottomrule
\end{tabular}
\par\endgroup
% course-panel:T105-P01:end
% mdtable:T105:end
\reserve{18\baselineskip}
% mdparagraph:P375:begin
\subsection*{cēosan: singular and plural commands}
% mdparagraph:P375:end

% teaching: u17-ceosan-imperative2
% mdparagraph:P376:begin
Use the unmutated present stem in the singular command and add -að for the plural. Only this complete command pair is added here.
% mdparagraph:P376:end

% table: u17-ceosan-imperative2-IMPERATIVE-DISPLAY
% mdtable:T106:begin
% course-table-style: 2026-10
% course-panel:T106-P01:begin
\par\begingroup
\fontsize{10.5}{13.2}\selectfont
\setlength{\LTleft}{0pt}\setlength{\LTright}{\fill}
\noindent
\begin{tabular}{@{}>{\raggedright\arraybackslash}p{63.342mm}>{\raggedright\arraybackslash}p{57.846mm}@{}}
\toprule
\textbf{Person and number} & \textbf{Command} \\
\midrule
2 sg. & cēos \\
2 pl. & cēosað \\
\bottomrule
\end{tabular}
\par\endgroup
% course-panel:T106-P01:end
% mdtable:T106:end
\reserve{7\baselineskip}
% mdparagraph:P377:begin
\subsection*{tōcuman: present subjunctive}
% mdparagraph:P377:end

% teaching: NEW-U17-VOC-TOCUMAN-U17-COME
\par\begin{minipage}{\linewidth}\raggedright
% mdparagraph:P378:begin
\textbf{tōcuman} — come; arrive.
% mdparagraph:P378:end

% teaching: u17-tocuman-present-subjunctive6
% mdparagraph:P379:begin
All three singular persons use -e and all three plural persons -en on the present stem, without an indicative second/\allowbreak{}third-singular vowel change. Only this tense of the subjunctive is added here.
% mdparagraph:P379:end

% table: u17-tocuman-present-subjunctive6-SUBJUNCTIVE-DISPLAY
% mdtable:T107:begin
% course-table-style: 2026-10
% course-panel:T107-P01:begin
\par\begingroup
\fontsize{10.5}{13.2}\selectfont
\setlength{\LTleft}{0pt}\setlength{\LTright}{\fill}
\noindent
\begin{tabular}{@{}>{\raggedright\arraybackslash}p{63.342mm}>{\raggedright\arraybackslash}p{57.846mm}@{}}
\toprule
\textbf{Person and number} & \textbf{Present} \\
\midrule
1 sg. & tōcume \\
2 sg. & tōcume \\
3 sg. & tōcume \\
1–3 pl. & tōcumen \\
\bottomrule
\end{tabular}
\par\endgroup
% course-panel:T107-P01:end
% mdtable:T107:end
\end{minipage}\par
\reserve{7\baselineskip}
% mdparagraph:P380:begin
\subsection*{geweorðan: present subjunctive}
% mdparagraph:P380:end

% teaching: NEW-U17-VOC-GEWEORTHAN-U17-BEFALL
\par\begin{minipage}{\linewidth}\raggedright
% mdparagraph:P381:begin
\textbf{geweorðan} — come to pass; befall.
% mdparagraph:P381:end

% teaching: u17-geweorthan-present-subjunctive6
% mdparagraph:P382:begin
All three singular persons use -e and all three plural persons -en on the present stem, without an indicative second/\allowbreak{}third-singular vowel change. Only this tense of the subjunctive is added here.
% mdparagraph:P382:end

% table: u17-geweorthan-present-subjunctive6-SUBJUNCTIVE-DISPLAY
% mdtable:T108:begin
% course-table-style: 2026-10
% course-panel:T108-P01:begin
\par\begingroup
\fontsize{10.5}{13.2}\selectfont
\setlength{\LTleft}{0pt}\setlength{\LTright}{\fill}
\noindent
\begin{tabular}{@{}>{\raggedright\arraybackslash}p{63.342mm}>{\raggedright\arraybackslash}p{57.846mm}@{}}
\toprule
\textbf{Person and number} & \textbf{Present} \\
\midrule
1 sg. & geweorðe \\
2 sg. & geweorðe \\
3 sg. & geweorðe \\
1–3 pl. & geweorðen \\
\bottomrule
\end{tabular}
\par\endgroup
% course-panel:T108-P01:end
% mdtable:T108:end
\end{minipage}\par
\reserve{18\baselineskip}
% mdparagraph:P383:begin
\subsection*{leornian: present subjunctive}
% mdparagraph:P383:end

% teaching: u17-leornian-present-subjunctive6
% mdparagraph:P384:begin
All three singular persons use -e and all three plural persons -en on the present stem; there is no indicative second/\allowbreak{}third-singular vowel change. With leornian, retain i in leornie and leornien.
% mdparagraph:P384:end

% table: u17-leornian-present-subjunctive6-SUBJUNCTIVE-DISPLAY
% mdtable:T109:begin
% course-table-style: 2026-10
% course-panel:T109-P01:begin
\par\begingroup
\fontsize{10.5}{13.2}\selectfont
\setlength{\LTleft}{0pt}\setlength{\LTright}{\fill}
\noindent
\begin{tabular}{@{}>{\raggedright\arraybackslash}p{63.342mm}>{\raggedright\arraybackslash}p{57.846mm}@{}}
\toprule
\textbf{Person and number} & \textbf{Present} \\
\midrule
1 sg. & leornie \\
2 sg. & leornie \\
3 sg. & leornie \\
1–3 pl. & leornien \\
\bottomrule
\end{tabular}
\par\endgroup
% course-panel:T109-P01:end
% mdtable:T109:end
\reserve{7\baselineskip}
% mdparagraph:P385:begin
\subsection*{āhyldan: present subjunctive}
% mdparagraph:P385:end

% teaching: NEW-U17-VOC-AHYLDAN-U17-LAY
\par\begin{minipage}{\linewidth}\raggedright
% mdparagraph:P386:begin
\textbf{āhyldan} — lay; incline.
% mdparagraph:P386:end

% teaching: u17-ahyldan-present-subjunctive6
% mdparagraph:P387:begin
All three singular persons use -e and all three plural persons -en on the present stem, without an indicative second/\allowbreak{}third-singular vowel change. Only this tense of the subjunctive is added here.
% mdparagraph:P387:end

% table: u17-ahyldan-present-subjunctive6-SUBJUNCTIVE-DISPLAY
% mdtable:T110:begin
% course-table-style: 2026-10
% course-panel:T110-P01:begin
\par\begingroup
\fontsize{10.5}{13.2}\selectfont
\setlength{\LTleft}{0pt}\setlength{\LTright}{\fill}
\noindent
\begin{tabular}{@{}>{\raggedright\arraybackslash}p{63.342mm}>{\raggedright\arraybackslash}p{57.846mm}@{}}
\toprule
\textbf{Person and number} & \textbf{Present} \\
\midrule
1 sg. & āhylde \\
2 sg. & āhylde \\
3 sg. & āhylde \\
1–3 pl. & āhylden \\
\bottomrule
\end{tabular}
\par\endgroup
% course-panel:T110-P01:end
% mdtable:T110:end
\end{minipage}\par
\reserve{7\baselineskip}
% page-review repair: U17LF02: participial agreement introduction with first complete singular panel
\par\begin{minipage}{\linewidth}\raggedright
% mdparagraph:P388:begin
\subsection*{ondrǣdan: present participle: strong and weak agreement}
% mdparagraph:P388:end

% teaching: u17-ondraedan-present-participle-agreement48
% mdparagraph:P389:begin
Remove final -e from the participle before adding an overt adjective ending. The strong table follows the grēne type: selected feminine nominative singular and neuter nominative/\allowbreak{}accusative plural -u, and feminine plural -a; the corresponding -e variants are also recorded. The weak forms follow the ordinary weak adjective pattern. The new independent use is describing or naming the specified human groups; no complete personal finite hyngran system is assigned.
% mdparagraph:P389:end

% mdparagraph:P390:begin
\textbf{Strong.}
% mdparagraph:P390:end

% table: u17-ondraedan-present-participle-agreement48-STRONG-DISPLAY
% mdtable:T111:begin
% course-table-style: 2026-10
% course-panel:T111-P01:begin
\par\begingroup
\fontsize{10.5}{13.2}\selectfont
\setlength{\LTleft}{0pt}\setlength{\LTright}{\fill}
\noindent
\begin{tabular}{@{}>{\raggedright\arraybackslash}p{22.0mm}>{\raggedright\arraybackslash}p{31.188mm}>{\raggedright\arraybackslash}p{31.188mm}>{\raggedright\arraybackslash}p{31.188mm}@{}}
\toprule
\textbf{Case} & \textbf{Masculine singular} & \textbf{Neuter singular} & \textbf{Feminine singular} \\
\midrule
Nominative & ondrǣdende & ondrǣdende & ondrǣdendu \\
Genitive & ondrǣdendes & ondrǣdendes & ondrǣdendre \\
Dative & ondrǣdendum & ondrǣdendum & ondrǣdendre \\
Accusative & ondrǣdendne & ondrǣdende & ondrǣdende \\
\bottomrule
\end{tabular}
\par\endgroup
% course-panel:T111-P01:end
\end{minipage}\par
\par\addvspace{5pt}
% course-panel:T111-P02:begin
\par\begingroup
\fontsize{10.5}{13.2}\selectfont
\setlength{\LTleft}{0pt}\setlength{\LTright}{\fill}
\noindent
\begin{tabular}{@{}>{\raggedright\arraybackslash}p{22.0mm}>{\raggedright\arraybackslash}p{31.188mm}>{\raggedright\arraybackslash}p{31.188mm}>{\raggedright\arraybackslash}p{31.188mm}@{}}
\toprule
\textbf{Case} & \textbf{Masculine plural} & \textbf{Neuter plural} & \textbf{Feminine plural} \\
\midrule
Nominative & ondrǣdende & ondrǣdendu & ondrǣdenda \\
Genitive & ondrǣdendra & ondrǣdendra & ondrǣdendra \\
Dative & ondrǣdendum & ondrǣdendum & ondrǣdendum \\
Accusative & ondrǣdende & ondrǣdendu & ondrǣdenda \\
\bottomrule
\end{tabular}
\par\endgroup
% course-panel:T111-P02:end
% mdtable:T111:end

\reserve{18\baselineskip}
\reserve{16\baselineskip}
% mdparagraph:P391:begin
\textbf{Weak.}
% mdparagraph:P391:end

% table: u17-ondraedan-present-participle-agreement48-WEAK-DISPLAY
% mdtable:T112:begin
% course-table-style: 2026-10
% course-panel:T112-P01:begin
\par\begingroup
\fontsize{10.5}{13.2}\selectfont
\setlength{\LTleft}{0pt}\setlength{\LTright}{\fill}
\noindent
\begin{tabular}{@{}>{\raggedright\arraybackslash}p{22.0mm}>{\raggedright\arraybackslash}p{31.188mm}>{\raggedright\arraybackslash}p{31.188mm}>{\raggedright\arraybackslash}p{31.188mm}@{}}
\toprule
\textbf{Case} & \textbf{Masculine singular} & \textbf{Neuter singular} & \textbf{Feminine singular} \\
\midrule
Nominative & ondrǣdenda & ondrǣdende & ondrǣdende \\
Genitive & ondrǣdendan & ondrǣdendan & ondrǣdendan \\
Dative & ondrǣdendan & ondrǣdendan & ondrǣdendan \\
Accusative & ondrǣdendan & ondrǣdende & ondrǣdendan \\
\bottomrule
\end{tabular}
\par\endgroup
% course-panel:T112-P01:end
\par\addvspace{5pt}
% course-panel:T112-P02:begin
\par\begingroup
\fontsize{10.5}{13.2}\selectfont
\setlength{\LTleft}{0pt}\setlength{\LTright}{\fill}
\noindent
\begin{tabular}{@{}>{\raggedright\arraybackslash}p{22.0mm}>{\raggedright\arraybackslash}p{31.188mm}>{\raggedright\arraybackslash}p{31.188mm}>{\raggedright\arraybackslash}p{31.188mm}@{}}
\toprule
\textbf{Case} & \textbf{Masculine plural} & \textbf{Neuter plural} & \textbf{Feminine plural} \\
\midrule
Nominative & ondrǣdendan & ondrǣdendan & ondrǣdendan \\
Genitive & ondrǣdendra & ondrǣdendra & ondrǣdendra \\
Dative & ondrǣdendum & ondrǣdendum & ondrǣdendum \\
Accusative & ondrǣdendan & ondrǣdendan & ondrǣdendan \\
\bottomrule
\end{tabular}
\par\endgroup
% course-panel:T112-P02:end
% mdtable:T112:end

\reserve{7\baselineskip}
% mdparagraph:P392:begin
\subsection*{hyngran: present participle: strong and weak agreement}
% mdparagraph:P392:end

% teaching: NEW-U17-VOC-HYNGRAN-U17-HUNGRY-PARTICIPLE
\par\begin{minipage}{\linewidth}\raggedright
% mdparagraph:P393:begin
\textbf{hyngran} — being hungry; hungering.
% mdparagraph:P393:end

% teaching: u17-hyngran-present-participle-agreement48
% mdparagraph:P394:begin
Remove final -e from the participle before adding an overt adjective ending. The strong table follows the grēne type: selected feminine nominative singular and neuter nominative/\allowbreak{}accusative plural -u, and feminine plural -a; the corresponding -e variants are also recorded. The weak forms follow the ordinary weak adjective pattern. The new independent use is describing or naming the specified human groups; no complete personal finite hyngran system is assigned.
% mdparagraph:P394:end

% mdparagraph:P395:begin
\textbf{Strong.}
% mdparagraph:P395:end

% table: u17-hyngran-present-participle-agreement48-STRONG-DISPLAY
% mdtable:T113:begin
% course-table-style: 2026-10
% course-panel:T113-P01:begin
\par\begingroup
\fontsize{10.5}{13.2}\selectfont
\setlength{\LTleft}{0pt}\setlength{\LTright}{\fill}
\noindent
\begin{tabular}{@{}>{\raggedright\arraybackslash}p{22.0mm}>{\raggedright\arraybackslash}p{31.188mm}>{\raggedright\arraybackslash}p{31.188mm}>{\raggedright\arraybackslash}p{31.188mm}@{}}
\toprule
\textbf{Case} & \textbf{Masculine singular} & \textbf{Neuter singular} & \textbf{Feminine singular} \\
\midrule
Nominative & hyngrende & hyngrende & hyngrendu \\
Genitive & hyngrendes & hyngrendes & hyngrendre \\
Dative & hyngrendum & hyngrendum & hyngrendre \\
Accusative & hyngrendne & hyngrende & hyngrende \\
\bottomrule
\end{tabular}
\par\endgroup
% course-panel:T113-P01:end
\end{minipage}\par
\par\addvspace{5pt}
% course-panel:T113-P02:begin
\par\begingroup
\fontsize{10.5}{13.2}\selectfont
\setlength{\LTleft}{0pt}\setlength{\LTright}{\fill}
\noindent
\begin{tabular}{@{}>{\raggedright\arraybackslash}p{22.0mm}>{\raggedright\arraybackslash}p{31.188mm}>{\raggedright\arraybackslash}p{31.188mm}>{\raggedright\arraybackslash}p{31.188mm}@{}}
\toprule
\textbf{Case} & \textbf{Masculine plural} & \textbf{Neuter plural} & \textbf{Feminine plural} \\
\midrule
Nominative & hyngrende & hyngrendu & hyngrenda \\
Genitive & hyngrendra & hyngrendra & hyngrendra \\
Dative & hyngrendum & hyngrendum & hyngrendum \\
Accusative & hyngrende & hyngrendu & hyngrenda \\
\bottomrule
\end{tabular}
\par\endgroup
% course-panel:T113-P02:end
% mdtable:T113:end

\reserve{18\baselineskip}
\reserve{16\baselineskip}
% mdparagraph:P396:begin
\textbf{Weak.}
% mdparagraph:P396:end

% table: u17-hyngran-present-participle-agreement48-WEAK-DISPLAY
% mdtable:T114:begin
% course-table-style: 2026-10
% course-panel:T114-P01:begin
\par\begingroup
\fontsize{10.5}{13.2}\selectfont
\setlength{\LTleft}{0pt}\setlength{\LTright}{\fill}
\noindent
\begin{tabular}{@{}>{\raggedright\arraybackslash}p{22.0mm}>{\raggedright\arraybackslash}p{31.188mm}>{\raggedright\arraybackslash}p{31.188mm}>{\raggedright\arraybackslash}p{31.188mm}@{}}
\toprule
\textbf{Case} & \textbf{Masculine singular} & \textbf{Neuter singular} & \textbf{Feminine singular} \\
\midrule
Nominative & hyngrenda & hyngrende & hyngrende \\
Genitive & hyngrendan & hyngrendan & hyngrendan \\
Dative & hyngrendan & hyngrendan & hyngrendan \\
Accusative & hyngrendan & hyngrende & hyngrendan \\
\bottomrule
\end{tabular}
\par\endgroup
% course-panel:T114-P01:end
\par\addvspace{5pt}
% course-panel:T114-P02:begin
\par\begingroup
\fontsize{10.5}{13.2}\selectfont
\setlength{\LTleft}{0pt}\setlength{\LTright}{\fill}
\noindent
\begin{tabular}{@{}>{\raggedright\arraybackslash}p{22.0mm}>{\raggedright\arraybackslash}p{31.188mm}>{\raggedright\arraybackslash}p{31.188mm}>{\raggedright\arraybackslash}p{31.188mm}@{}}
\toprule
\textbf{Case} & \textbf{Masculine plural} & \textbf{Neuter plural} & \textbf{Feminine plural} \\
\midrule
Nominative & hyngrendan & hyngrendan & hyngrendan \\
Genitive & hyngrendra & hyngrendra & hyngrendra \\
Dative & hyngrendum & hyngrendum & hyngrendum \\
Accusative & hyngrendan & hyngrendan & hyngrendan \\
\bottomrule
\end{tabular}
\par\endgroup
% course-panel:T114-P02:end
% mdtable:T114:end

\reserve{7\baselineskip}
% mdparagraph:P397:begin
\subsection*{gesēon: preterite subjunctive}
% mdparagraph:P397:end

% teaching: NEW-U17-VOC-GESEON-U17-REGARD
\par\begin{minipage}{\linewidth}\raggedright
% mdparagraph:P398:begin
\textbf{gesēon} — regard; take notice of.
% mdparagraph:P398:end

% teaching: u17-geseon-preterite-subjunctive6
% mdparagraph:P399:begin
Build the past subjunctive on the plural preterite stem gesāw-. All three singular persons have -e and all three plural persons -en. This system supplies the before-clause’s first singular gesāwe.
% mdparagraph:P399:end

% table: u17-geseon-preterite-subjunctive6-SUBJUNCTIVE-DISPLAY
% mdtable:T115:begin
% course-table-style: 2026-10
% course-panel:T115-P01:begin
\par\begingroup
\fontsize{10.5}{13.2}\selectfont
\setlength{\LTleft}{0pt}\setlength{\LTright}{\fill}
\noindent
\begin{tabular}{@{}>{\raggedright\arraybackslash}p{63.342mm}>{\raggedright\arraybackslash}p{57.846mm}@{}}
\toprule
\textbf{Person and number} & \textbf{Preterite} \\
\midrule
1 sg. & gesāwe \\
2 sg. & gesāwe \\
3 sg. & gesāwe \\
1–3 pl. & gesāwen \\
\bottomrule
\end{tabular}
\par\endgroup
% course-panel:T115-P01:end
% mdtable:T115:end
\end{minipage}\par
\reserve{7\baselineskip}
% mdparagraph:P400:begin
\subsection*{sęttan: present subjunctive}
% mdparagraph:P400:end

% teaching: NEW-U17-VOC-SETTAN-U17-ESTABLISH
\par\begin{minipage}{\linewidth}\raggedright
% mdparagraph:P401:begin
\textbf{sęttan} — appoint; establish; make a decree.
% mdparagraph:P401:end

% teaching: u17-settan-present-subjunctive6
% mdparagraph:P402:begin
All three singular persons use -e and all three plural persons -en on the present stem, without an indicative second/\allowbreak{}third-singular vowel change. Only this tense of the subjunctive is added here.
% mdparagraph:P402:end

% table: u17-settan-present-subjunctive6-SUBJUNCTIVE-DISPLAY
% mdtable:T116:begin
% course-table-style: 2026-10
% course-panel:T116-P01:begin
\par\begingroup
\fontsize{10.5}{13.2}\selectfont
\setlength{\LTleft}{0pt}\setlength{\LTright}{\fill}
\noindent
\begin{tabular}{@{}>{\raggedright\arraybackslash}p{63.342mm}>{\raggedright\arraybackslash}p{57.846mm}@{}}
\toprule
\textbf{Person and number} & \textbf{Present} \\
\midrule
1 sg. & sętte \\
2 sg. & sętte \\
3 sg. & sętte \\
1–3 pl. & sętten \\
\bottomrule
\end{tabular}
\par\endgroup
% course-panel:T116-P01:end
% mdtable:T116:end
\end{minipage}\par
\reserve{18\baselineskip}
% mdparagraph:P403:begin
\subsection*{etan: present subjunctive}
% mdparagraph:P403:end

% teaching: u17-etan-present-subjunctive6
% mdparagraph:P404:begin
All three singular persons use -e and all three plural persons -en on the present stem, without an indicative second/\allowbreak{}third-singular vowel change. Only this tense of the subjunctive is added here.
% mdparagraph:P404:end

% table: u17-etan-present-subjunctive6-SUBJUNCTIVE-DISPLAY
% mdtable:T117:begin
% course-table-style: 2026-10
% course-panel:T117-P01:begin
\par\begingroup
\fontsize{10.5}{13.2}\selectfont
\setlength{\LTleft}{0pt}\setlength{\LTright}{\fill}
\noindent
\begin{tabular}{@{}>{\raggedright\arraybackslash}p{63.342mm}>{\raggedright\arraybackslash}p{57.846mm}@{}}
\toprule
\textbf{Person and number} & \textbf{Present} \\
\midrule
1 sg. & ete \\
2 sg. & ete \\
3 sg. & ete \\
1–3 pl. & eten \\
\bottomrule
\end{tabular}
\par\endgroup
% course-panel:T117-P01:end
% mdtable:T117:end
\reserve{18\baselineskip}
% mdparagraph:P405:begin
\subsection*{forgiefan: singular and plural commands}
% mdparagraph:P405:end

% teaching: u17-forgiefan-imperative2
% mdparagraph:P406:begin
Use the unmutated present stem for the singular command and -að for the plural. In the prayer, the meaning is forgive, with the familiar thing/\allowbreak{}person case distinction. Only this complete command pair is added here.
% mdparagraph:P406:end

% table: u17-forgiefan-imperative2-IMPERATIVE-DISPLAY
% mdtable:T118:begin
% course-table-style: 2026-10
% course-panel:T118-P01:begin
\par\begingroup
\fontsize{10.5}{13.2}\selectfont
\setlength{\LTleft}{0pt}\setlength{\LTright}{\fill}
\noindent
\begin{tabular}{@{}>{\raggedright\arraybackslash}p{63.342mm}>{\raggedright\arraybackslash}p{57.846mm}@{}}
\toprule
\textbf{Person and number} & \textbf{Command} \\
\midrule
2 sg. & forgief \\
2 pl. & forgiefað \\
\bottomrule
\end{tabular}
\par\endgroup
% course-panel:T118-P01:end
% mdtable:T118:end
\reserve{18\baselineskip}
% mdparagraph:P407:begin
\subsection*{gesēon: plain and inflected infinitives}
% mdparagraph:P407:end

% teaching: u17-geseon-infinitive-pair2
% mdparagraph:P408:begin
The pair distinguishes the plain infinitive from the inflected infinitive. The selected series uses -anne, with -ianne in leornianne and -onne after contracted gesēon. The alternative -enne series also exists. The modal constructions here require the plain infinitive; this table adds no new purpose-clause use.
% mdparagraph:P408:end

% table: u17-geseon-infinitive-pair2-BOUNDED-NONFINITE-DISPLAY
% mdtable:T119:begin
\par\begingroup
\fontsize{10.5}{13.2}\selectfont
\setlength{\LTleft}{0pt}\setlength{\LTright}{\fill}
\noindent
\begin{tabular}{@{}>{\raggedright\arraybackslash}p{76.37mm}>{\raggedright\arraybackslash}p{44.82mm}@{}}
\toprule
\textbf{Form} & \textbf{Literal surface} \\
\midrule
verb form plain infinitive & gesēon \\
verb form inflected infinitive & gesēonne \\
\bottomrule
\end{tabular}
\par\endgroup

% mdtable:T119:end
\reserve{18\baselineskip}
% mdparagraph:P409:begin
\subsection*{gelīefan: plain and inflected infinitives}
% mdparagraph:P409:end

% teaching: u17-geliefan-infinitive-pair2
% mdparagraph:P410:begin
The pair distinguishes the plain infinitive from the inflected infinitive. The selected series uses -anne, with -ianne in leornianne and -onne after contracted gesēon. The alternative -enne series also exists. The modal constructions here require the plain infinitive; this table adds no new purpose-clause use.
% mdparagraph:P410:end

% table: u17-geliefan-infinitive-pair2-BOUNDED-NONFINITE-DISPLAY
% mdtable:T120:begin
\par\begingroup
\fontsize{10.5}{13.2}\selectfont
\setlength{\LTleft}{0pt}\setlength{\LTright}{\fill}
\noindent
\begin{tabular}{@{}>{\raggedright\arraybackslash}p{73.68mm}>{\raggedright\arraybackslash}p{47.51mm}@{}}
\toprule
\textbf{Form} & \textbf{Literal surface} \\
\midrule
verb form plain infinitive & gelīefan \\
verb form inflected infinitive & gelīefanne \\
\bottomrule
\end{tabular}
\par\endgroup

% mdtable:T120:end
\reserve{18\baselineskip}
% mdparagraph:P411:begin
\subsection*{leornian: plain and inflected infinitives}
% mdparagraph:P411:end

% teaching: u17-leornian-infinitive-pair2
% mdparagraph:P412:begin
The pair distinguishes the plain infinitive from the inflected infinitive. The selected series uses -anne, with -ianne in leornianne and -onne after contracted gesēon. The alternative -enne series also exists. The modal constructions here require the plain infinitive; this table adds no new purpose-clause use.
% mdparagraph:P412:end

% table: u17-leornian-infinitive-pair2-BOUNDED-NONFINITE-DISPLAY
% mdtable:T121:begin
\par\begingroup
\fontsize{10.5}{13.2}\selectfont
\setlength{\LTleft}{0pt}\setlength{\LTright}{\fill}
\noindent
\begin{tabular}{@{}>{\raggedright\arraybackslash}p{73.68mm}>{\raggedright\arraybackslash}p{47.51mm}@{}}
\toprule
\textbf{Form} & \textbf{Literal surface} \\
\midrule
verb form plain infinitive & leornian \\
verb form inflected infinitive & leornianne \\
\bottomrule
\end{tabular}
\par\endgroup

% mdtable:T121:end
\reserve{7\baselineskip}
% mdparagraph:P413:begin
\subsection*{gescieppan: present and past participles}
% mdparagraph:P413:end

% teaching: u17-gescieppan-participle-pair2
% mdparagraph:P414:begin
Retain the lexical ge-. The present participle uses the present stem, while gesceapen has ea and single p.
% mdparagraph:P414:end

\reserve{10\baselineskip}
% mdparagraph:P415:begin
This display assigns the complete present/\allowbreak{}past formation pair. It does not add other finite moods or an unrestricted attributive declension; the past form is used in the stated passive frame.
% mdparagraph:P415:end

% table: u17-gescieppan-participle-pair2-PARTICIPLES-DISPLAY
% mdtable:T122:begin
\par\begingroup
\fontsize{10.5}{13.2}\selectfont
\setlength{\LTleft}{0pt}\setlength{\LTright}{\fill}
\noindent
\begin{tabular}{@{}>{\raggedright\arraybackslash}p{51.9mm}>{\raggedright\arraybackslash}p{69.29mm}@{}}
\toprule
\textbf{Form} & \textbf{Literal surface} \\
\midrule
present & gescieppende \\
past & gesceapen \\
\bottomrule
\end{tabular}
\par\endgroup

% mdtable:T122:end
\reserve{18\baselineskip}
% mdparagraph:P416:begin
\subsection*{wyrcean: base and plural predicate participle}
% mdparagraph:P416:end

% teaching: u17-wyrcean-past-predicate2
% mdparagraph:P417:begin
The base is selected for singular predicates and the -e form for the plural predicate in Cook 37.5. The uninflected base is also permitted with a plural predicate. These are selected alternatives, not a complete attributive declension.
% mdparagraph:P417:end

% table: u17-wyrcean-past-predicate2-DISPLAY
% mdtable:T123:begin
\par\begingroup
\fontsize{10.5}{13.2}\selectfont
\setlength{\LTleft}{0pt}\setlength{\LTright}{\fill}
\noindent
\begin{tabular}{@{}>{\raggedright\arraybackslash}p{87.2mm}>{\raggedright\arraybackslash}p{33.98mm}@{}}
\toprule
\textbf{Predicate form} & \textbf{Literal surface} \\
\midrule
Uninflected predicate base (singular or plural) & worht \\
Agreeing plural predicate & worhte \\
\bottomrule
\end{tabular}
\par\endgroup

% mdtable:T123:end
\reserve{18\baselineskip}
% mdparagraph:P418:begin
\subsection*{gescieppan: base and plural predicate participle}
% mdparagraph:P418:end

% teaching: u17-gescieppan-past-predicate2
% mdparagraph:P419:begin
The base is selected for singular predicates and the -e form for the plural predicate in Cook 37.5. The uninflected base is also permitted with a plural predicate. These are selected alternatives, not a complete attributive declension.
% mdparagraph:P419:end

% table: u17-gescieppan-past-predicate2-DISPLAY
% mdtable:T124:begin
\par\begingroup
\fontsize{10.5}{13.2}\selectfont
\setlength{\LTleft}{0pt}\setlength{\LTright}{\fill}
\noindent
\begin{tabular}{@{}>{\raggedright\arraybackslash}p{83.56mm}>{\raggedright\arraybackslash}p{37.63mm}@{}}
\toprule
\textbf{Predicate form} & \textbf{Literal surface} \\
\midrule
Uninflected predicate base (singular or plural) & gesceapen \\
Agreeing plural predicate & gesceapene \\
\bottomrule
\end{tabular}
\par\endgroup

% mdtable:T124:end
\reserve{7\baselineskip}
% mdparagraph:P420:begin
\section*{4. Meanings and complements}
% mdparagraph:P420:end

% mdparagraph:P421:begin
The following meanings extend familiar words. Use their previously learned forms only in the constructions now stated. The new noun and adjective tables above already give their own meanings; the additional distinctions here prevent a familiar spelling from hiding a new lexical task.
% mdparagraph:P421:end

% teaching: NEW-U17-VOC-HEALDAN-U17-MAINTAIN
% mdparagraph:P422:begin
\textbf{healdan} — maintain; sustain.
% mdparagraph:P422:end

% teaching: NEW-U17-VOC-SAED-U17-DESCENDANTS
% mdparagraph:P423:begin
\textbf{sǣd} — seed; offspring; descendants collectively.
% mdparagraph:P423:end

% teaching: NEW-U17-VOC-FOLC-U17-PEOPLES
% mdparagraph:P424:begin
\textbf{folc} — peoples; distinct peoples or national groups.
% mdparagraph:P424:end

% teaching: NEW-U17-VOC-EALD-U17-ELDERLY
% mdparagraph:P425:begin
\textbf{eald} — old people; the elderly.
% mdparagraph:P425:end

% teaching: NEW-U17-VOC-WER-U17-HUSBAND
% mdparagraph:P426:begin
\textbf{wer} — husband.
% mdparagraph:P426:end

% teaching: NEW-U17-VOC-DON-U17-FULFIL
% mdparagraph:P427:begin
\textbf{dōn} — fulfil; carry out a word or command.
% mdparagraph:P427:end

% teaching: NEW-U17-VOC-SECGAN-U17-CALL
% mdparagraph:P428:begin
\textbf{sęcgan} — call; declare to be.
% mdparagraph:P428:end

% teaching: NEW-U17-VOC-LAEDAN-U17-CARRY
% mdparagraph:P429:begin
\textbf{lǣdan} — take; carry along.
% mdparagraph:P429:end

% teaching: NEW-U17-VOC-ONFON-U17-PERSON
% mdparagraph:P430:begin
\textbf{onfōn} — receive; accept a person into care.
% mdparagraph:P430:end

% teaching: NEW-U17-VOC-FORLAETAN-U17-LET-GO
% mdparagraph:P431:begin
\textbf{forlǣtan} — let go; dismiss.
% mdparagraph:P431:end

\reserve{7\baselineskip}
% mdparagraph:P432:begin
\subsection*{Participles naming people}
% mdparagraph:P432:end

% teaching: NEW-U17-GRAM-U17-PARTICIPLE-HUMAN-NOMINALIZATION
% mdparagraph:P433:begin
A present participle can describe people without an expressed noun. Its ending gives the phrase’s role in the surrounding clause, while the verb inside the participle retains its own complements. In hine ondrǣdendum, the dative plural names the beneficiaries of mercy; accusative hine names the person feared. In ðā hyngrendan, the article and weak accusative plural identify hungry people as a direct object. These are nonfinite descriptions, not present-tense assertions. Learn these human uses of ondrǣdende and hyngrende; no personal finite conjugation of hyngran is added.
% mdparagraph:P433:end

\reserve{7\baselineskip}
% mdparagraph:P434:begin
\subsection*{Past participles in a predicate}
% mdparagraph:P434:end

% teaching: NEW-U17-GRAM-U17-PREDICATE-PARTICIPLE-PLURAL
% mdparagraph:P435:begin
A past participle with a form of be can remain uninflected or take the strong masculine agreement endings, regardless of the subject’s gender. The selected plural predicates in Cook 37.5 are worhte and gesceapene. They belong with wǣron; they are not finite weak preterites. The uninflected bases worht and gesceapen are also possible with plural subjects. The two small displays in the preceding section distinguish the base from the selected agreeing plural.
% mdparagraph:P435:end

\reserve{7\baselineskip}
% mdparagraph:P436:begin
\subsection*{Praise: statement, command and wish}
% mdparagraph:P436:end

% teaching: NEW-U17-GRAM-U17-HERIAN-OBJECT
% mdparagraph:P437:begin
Hęrian takes the praised person or thing in the accusative. A plural command addresses the beings named in the catalogue; a third-person wish has hīe as subject and hęrien as verb. The identical indicative and imperative form hęriað does not make their functions identical.
% mdparagraph:P437:end

\reserve{7\baselineskip}
% mdparagraph:P438:begin
\subsection*{Requesting and seeking in marriage}
% mdparagraph:P438:end

% teaching: NEW-U17-GRAM-U17-GIERNAN-REQUEST-WOO
% mdparagraph:P439:begin
Giernan governs a genitive of the requested thing or marriage partner. Æt with the dative gives the person asked. In Cook 39.3, keep the outer accusative role of eall ðæt separate from the role represented by invariant ðe in its relative clause. In Cook 43.14, seeking the maiden in marriage requires her genitive form.
% mdparagraph:P439:end

\reserve{7\baselineskip}
% mdparagraph:P440:begin
\subsection*{Contents and covering}
% mdparagraph:P440:end

% teaching: NEW-U17-GRAM-U17-FILLING-MATERIAL
% mdparagraph:P441:begin
Āfyllan takes the filled entity in the accusative and its contents after mid with the dative. In a passive, the filled entity becomes the nominative subject. Oferwyrcean has a similar case arrangement for a different act: the throne is overlaid with gold. Learn each verb’s meaning together with its complements.
% mdparagraph:P441:end

\reserve{7\baselineskip}
% mdparagraph:P442:begin
\subsection*{Receiving a thing or a person}
% mdparagraph:P442:end

% teaching: NEW-U17-GRAM-U17-ONFON-PERSON-DATIVE
% mdparagraph:P443:begin
The selected onfōn construction takes the dative: edlēane for a received reward, and his cnihte Israel for the servant accepted into care. Accusative receiving constructions exist elsewhere; these exercises specifically practise the dative frame.
% mdparagraph:P443:end

\reserve{7\baselineskip}
% mdparagraph:P444:begin
\subsection*{Calling someone blessed}
% mdparagraph:P444:end

% teaching: NEW-U17-GRAM-U17-SECGAN-PREDICATE
% mdparagraph:P445:begin
Sęcgan can take an object and an adjective that agrees with it: the generations call Mary blessed. In Cook 40.3, mē is the first-person singular object, ēadige its feminine accusative singular predicate, and sęcgað agrees with the plural generations. The present verb has future reference in the prediction.
% mdparagraph:P445:end

\reserve{7\baselineskip}
% mdparagraph:P446:begin
\subsection*{New meanings of familiar verbs}
% mdparagraph:P446:end

% teaching: NEW-U17-GRAM-U17-VERB-LEXICAL-EXTENSIONS
% mdparagraph:P447:begin
Dōn can mean fulfil or carry out a word or command; healdan can mean sustain all things; sęcgan can mean call or declare someone to be something. Learn those meanings explicitly. Wyrcean retains the familiar work or perform sense with the Father’s will and the make or build sense with a foundation.
% mdparagraph:P447:end

\reserve{7\baselineskip}
% mdparagraph:P448:begin
\subsection*{Indefinite relatives and the subjunctive}
% mdparagraph:P448:end

% teaching: NEW-U17-GRAM-U17-VIRTUAL-RELATIVE
% mdparagraph:P449:begin
A relative clause can describe an indefinite or possible situation: the one you wish to choose, or a place where one might lay one’s head. Such clauses can take the subjunctive. Wille in Cook 43.12 is second singular present subjunctive; āhylde in Cook 42.7 is third singular present subjunctive. The construction does not make every relative subjunctive, and English may need not be translated by an extra modal.
% mdparagraph:P449:end

\reserve{7\baselineskip}
% mdparagraph:P450:begin
\subsection*{Learning from someone}
% mdparagraph:P450:end

% teaching: NEW-U17-GRAM-U17-LEORNIAN-SOURCE
% mdparagraph:P451:begin
Of with the dative can name the person from whom one learns in Cook 43.11. Gewilnian introduces the desired content with ðæt. The embedded ic takes first singular leornie, while plural gē takes habbað. The separate clauses keep their own subjects and agreement.
% mdparagraph:P451:end

\reserve{7\baselineskip}
% mdparagraph:P452:begin
\subsection*{Playing with an object}
% mdparagraph:P452:end

% teaching: NEW-U17-GRAM-U17-PLEGIAN-INSTRUMENT
% mdparagraph:P453:begin
Plegian in the ordinary play sense takes mid with the dative instrument. The ball exercise therefore uses the familiar dative ðoðore. The reading’s act or perform sense is supplied separately.
% mdparagraph:P453:end

\reserve{7\baselineskip}
% mdparagraph:P454:begin
\subsection*{Before a past event}
% mdparagraph:P454:end

% teaching: NEW-U17-GRAM-U17-TEMPORAL-BEFORE
% mdparagraph:P455:begin
In Cook 38.8, nolde expresses a former unwillingness to believe. The seeing is still prospective from that earlier viewpoint: the before-clause uses first singular past subjunctive gesāwe. Its spelling alone does not make it second singular indicative. Retain self and follow the requested final position of the verb.
% mdparagraph:P455:end

\reserve{7\baselineskip}
% mdparagraph:P456:begin
\subsection*{Negative concord}
% mdparagraph:P456:end

% teaching: NEW-U17-GRAM-U15-NEGATIVE-CONCORD
% mdparagraph:P457:begin
Preverbal ne can combine with nān or another negative word to express one negation. In the coordinated eye, ear and heart sentence, the first mæg can govern the following infinitives as well. Identify the subject of each part; repeated negative words do not make it affirmative. Use the ordinary ne witon for we do not know. The reading’s contracted nāst is supplied separately.
% mdparagraph:P457:end

\reserve{7\baselineskip}
% mdparagraph:P458:begin
\subsection*{Choosing a past tense}
% mdparagraph:P458:end

% teaching: NEW-U17-GRAM-U17-PAST-TENSE-SELECTION
% mdparagraph:P459:begin
The simple preterite can correspond to English simple past, present perfect or past perfect when the context establishes the time relation. Cook §188 explicitly permits this range. An English has or had therefore does not always require a separate auxiliary in Old English. Follow an explicit omission instruction; elsewhere choose a construction whose exact forms and intended time relation you know.
% mdparagraph:P459:end

\reserve{7\baselineskip}
% mdparagraph:P460:begin
\subsection*{The grounds for a reward}
% mdparagraph:P460:end

% teaching: NEW-U17-GRAM-U17-GENITIVE-REWARD-CAUSE
% mdparagraph:P461:begin
A genitive dependent on a reward can express its grounds: his weorces, for his work. In Cook 39.6, edlēane is dative because it is received with onfōn; the work-genitive depends on the reward, and æt Gode gives its source.
% mdparagraph:P461:end

% mdparagraph:P462:begin
For sǣd in the promise to Abraham, the descendants are expressed collectively. The assigned dative sǣde is singular. Folc can denote distinct peoples in the plural catalogue, while the collective people in the dense horn passage remains singular. Wer now also means husband. Keep these contextual meanings separate from the earlier plant seed, people and man meanings.
% mdparagraph:P462:end

\reserve{7\baselineskip}
% mdparagraph:P463:begin
\section*{5. Form contrasts and composition}
% mdparagraph:P463:end

% mdparagraph:P464:begin
Complete each requested response. The form contrasts test case, number and clause function in meaningful expressions. Keep supplied dictionary forms available, but choose the required endings from the tables and explanations. In the compositions, preserve every participant and the stated clause relationship.
% mdparagraph:P464:end

\reserve{8\baselineskip}
% mdparagraph:P465:begin
\subsection*{F01}
% mdparagraph:P465:end

% item: NEW-U17-F01
% mdparagraph:P466:begin
Use cniht in its explicitly taught servant sense and bōc book. Translate the four expressions. In each, choose the noun's case from its role rather than from an English word ending.
% mdparagraph:P466:end

% mdparagraph:P467:begin
cniht
% mdparagraph:P467:end

% mdparagraph:P468:begin
bōc
% mdparagraph:P468:end

% output: NEW-U17-F01-O001
% mdparagraph:P469:begin
\itempara{1}{the servant's book}
% mdparagraph:P469:end

% output: NEW-U17-F01-O002
% mdparagraph:P470:begin
\itempara{2}{the servants' books}
% mdparagraph:P470:end

% output: NEW-U17-F01-O003
% mdparagraph:P471:begin
\itempara{3}{I give the servant the book.}
% mdparagraph:P471:end

% output: NEW-U17-F01-O004
% mdparagraph:P472:begin
\itempara{4}{We give the servants the books.}
% mdparagraph:P472:end

\reserve{8\baselineskip}
% mdparagraph:P473:begin
\subsection*{F02}
% mdparagraph:P473:end

% item: NEW-U17-F02
% mdparagraph:P474:begin
Use the selected endingless accusative of cwēn queen, while retaining its genitive/\allowbreak{}dative ending. Translate the three expressions, then explain the contrast.
% mdparagraph:P474:end

% mdparagraph:P475:begin
cwēn
% mdparagraph:P475:end

% mdparagraph:P476:begin
bōc
% mdparagraph:P476:end

% output: NEW-U17-F02-O001
% mdparagraph:P477:begin
\itempara{1}{I saw the queen.}
% mdparagraph:P477:end

% output: NEW-U17-F02-O002
% mdparagraph:P478:begin
\itempara{2}{I gave the queen the book.}
% mdparagraph:P478:end

% output: NEW-U17-F02-O003
% mdparagraph:P479:begin
\itempara{3}{the queen's book}
% mdparagraph:P479:end

% output: NEW-U17-F02-O004
% mdparagraph:P480:begin
\itempara{4}{Explain what the three queen forms show about identifying case from an ending alone.}
% mdparagraph:P480:end

\reserve{8\baselineskip}
% mdparagraph:P481:begin
\subsection*{F03}
% mdparagraph:P481:end

% item: NEW-U17-F03
% mdparagraph:P482:begin
Use fox fox, hol hole and nest nest. Translate each sentence. Keep ordinary locative on with the dative and distinguish the two neuter plural patterns.
% mdparagraph:P482:end

% mdparagraph:P483:begin
fox
% mdparagraph:P483:end

% mdparagraph:P484:begin
hol
% mdparagraph:P484:end

% mdparagraph:P485:begin
nest
% mdparagraph:P485:end

% output: NEW-U17-F03-O001
% mdparagraph:P486:begin
\itempara{1}{The fox is in the hole.}
% mdparagraph:P486:end

% output: NEW-U17-F03-O002
% mdparagraph:P487:begin
\itempara{2}{The foxes are in the holes.}
% mdparagraph:P487:end

% output: NEW-U17-F03-O003
% mdparagraph:P488:begin
\itempara{3}{The nest is in the tree.}
% mdparagraph:P488:end

% output: NEW-U17-F03-O004
% mdparagraph:P489:begin
\itempara{4}{The nests are in the trees.}
% mdparagraph:P489:end

% output: NEW-U17-F03-O005
% mdparagraph:P490:begin
\itempara{5}{The holes are large.}
% mdparagraph:P490:end

\reserve{8\baselineskip}
% mdparagraph:P491:begin
\subsection*{F04}
% mdparagraph:P491:end

% item: NEW-U17-F04
% mdparagraph:P492:begin
Use hęrian praise. Supply only the finite form required in each stated setting, then explain the contrast. The dictionary form is visible reference help.
% mdparagraph:P492:end

% mdparagraph:P493:begin
hęrian
% mdparagraph:P493:end

% output: NEW-U17-F04-O001
% mdparagraph:P494:begin
\itempara{1}{A narrator states: they praise. Give present indicative third plural.}
% mdparagraph:P494:end

% output: NEW-U17-F04-O002
% mdparagraph:P495:begin
\itempara{2}{A speaker commands several listeners: praise! Give second-plural imperative.}
% mdparagraph:P495:end

% output: NEW-U17-F04-O003
% mdparagraph:P496:begin
\itempara{3}{A speaker wishes: may they praise. Give present subjunctive third plural.}
% mdparagraph:P496:end

% output: NEW-U17-F04-O004
% mdparagraph:P497:begin
\itempara{4}{Which answers are spelled alike, and what distinguishes their functions?}
% mdparagraph:P497:end

\reserve{8\baselineskip}
% mdparagraph:P498:begin
\subsection*{F05}
% mdparagraph:P498:end

% item: NEW-U17-F05
% mdparagraph:P499:begin
Use the explicitly taught present participles of ondrǣdan fear and hyngran be hungry. Translate the two nominal phrases, then distinguish each phrase's outer case from any object governed inside it.
% mdparagraph:P499:end

% mdparagraph:P500:begin
ondrǣdan
% mdparagraph:P500:end

% mdparagraph:P501:begin
hyngran
% mdparagraph:P501:end

% output: NEW-U17-F05-O001
% mdparagraph:P502:begin
\itempara{1}{to those fearing him: use a dative plural participle without an expressed people noun; him is the one feared}
% mdparagraph:P502:end

% output: NEW-U17-F05-O002
% mdparagraph:P503:begin
\itempara{2}{the hungry: give the weak accusative plural phrase used as a direct object}
% mdparagraph:P503:end

% output: NEW-U17-F05-O003
% mdparagraph:P504:begin
\itempara{3}{State the two grammatical relations in the fear phrase.}
% mdparagraph:P504:end

% output: NEW-U17-F05-O004
% mdparagraph:P505:begin
\itempara{4}{Explain why a present participle can name people here without being a finite verb.}
% mdparagraph:P505:end

\reserve{8\baselineskip}
% mdparagraph:P506:begin
\subsection*{C01}
% mdparagraph:P506:end

% item: NEW-U17-C01
% mdparagraph:P507:begin
Translate: The king saw the men who sang. Use an invariant relative particle.
% mdparagraph:P507:end

% output: NEW-U17-C01-O001
% mdparagraph:P508:begin
\itempara{1}{Write one complete Old English sentence.}
% mdparagraph:P508:end

\reserve{8\baselineskip}
% mdparagraph:P509:begin
\subsection*{C02}
% mdparagraph:P509:end

% item: NEW-U17-C02
% mdparagraph:P510:begin
Translate: The king is as good as the men are. Express equal degree with swā … swā.
% mdparagraph:P510:end

% output: NEW-U17-C02-O001
% mdparagraph:P511:begin
\itempara{1}{Write one complete Old English sentence.}
% mdparagraph:P511:end

\reserve{8\baselineskip}
% mdparagraph:P512:begin
\subsection*{C03}
% mdparagraph:P512:end

% item: NEW-U17-C03
% mdparagraph:P513:begin
A visit is possible tomorrow. Translate: If he comes, give (plural) him the book.
% mdparagraph:P513:end

% output: NEW-U17-C03-O001
% mdparagraph:P514:begin
\itempara{1}{Write one complete Old English sentence.}
% mdparagraph:P514:end

\reserve{8\baselineskip}
% mdparagraph:P515:begin
\subsection*{C04}
% mdparagraph:P515:end

% item: NEW-U17-C04
% mdparagraph:P516:begin
Translate: The man says that the king is good. Present the content as a report, using the present subjunctive.
% mdparagraph:P516:end

% output: NEW-U17-C04-O001
% mdparagraph:P517:begin
\itempara{1}{Write one complete Old English sentence.}
% mdparagraph:P517:end

\reserve{7\baselineskip}
% mdparagraph:P518:begin
\section*{6. Cook: connected composition}
% mdparagraph:P518:end

% teaching: NEW-U17-SOURCE-HINTS
% teaching: NEW-U17-VERBAL-CONTEXT-HELP
% mdparagraph:P519:begin
Translate the complete English prompts. Parenthetic Old English is visible lexical or phrase help; a headword may require inflection, while an exact supplied expression is used as given. Parenthetic numbers refer to Cook’s grammar: the relevant constructions are explained here, so the locator is not a substitute for instruction. An omission control removes the specified English word from the Old English wording, while preserving its intended relation. Brackets supply the requested paraphrase. Distinguish a grammatical alternative from compliance with an explicit word, order, number or omission target.
% mdparagraph:P519:end

% mdparagraph:P520:begin
All supplied expressions and meanings remain available during these exercises. Copying an exact given form is a reference response. When the hint gives a headword but the task requires a different inflection, its lexical meaning is still given; the new ending is the part you must choose.
% mdparagraph:P520:end

\reserve{7\baselineskip}
% mdparagraph:P521:begin
\section*{Cook 37. Praise and enumeration}
% mdparagraph:P521:end

% mdparagraph:P522:begin
The catalogue alternates addressed beings with statements and wishes about them. Track who is addressed, who acts, and who receives praise. A list can leave a governing verb understood from its setting; keep the listed items complete.
% mdparagraph:P522:end

% mdparagraph:P523:begin
[In this exercise, omit the article and the second personal pronoun throughout.]
% mdparagraph:P523:end

% mdparagraph:P524:begin
This applies to all fourteen items: omit the article and second-person pronoun. It does not remove a third-person subject needed for “let them praise”.
% mdparagraph:P524:end

\par\begin{minipage}{\linewidth}\raggedright
% mdparagraph:P525:begin
\subsection*{Cook 37.1}
% mdparagraph:P525:end

% item: EX-XII-XIII-37-01
% mdparagraph:P526:begin
Praise ye the Lord from (of) the heavens; praise him in the heights.
% mdparagraph:P526:end

% output: EX-XII-XIII-37-01-O001
% mdparagraph:P527:begin
Write one complete Old English translation, keeping the printed controls and the supplied expressions available.
% mdparagraph:P527:end

\end{minipage}\par\addvspace{4pt}

\par\begin{minipage}{\linewidth}\raggedright
% mdparagraph:P528:begin
\subsection*{Cook 37.2}
% mdparagraph:P528:end

% item: EX-XII-XIII-37-02
% mdparagraph:P529:begin
Praise ye him, all his angels; praise him, all his powers (mægen).
% mdparagraph:P529:end

% output: EX-XII-XIII-37-02-O001
% mdparagraph:P530:begin
Write one complete Old English translation, keeping the printed controls and the supplied expressions available.
% mdparagraph:P530:end

\end{minipage}\par\addvspace{4pt}

\par\begin{minipage}{\linewidth}\raggedright
% mdparagraph:P531:begin
\subsection*{Cook 37.3}
% mdparagraph:P531:end

% item: EX-XII-XIII-37-03
% mdparagraph:P532:begin
Praise ye him, sun and moon; praise him, all ye stars and light (lēoht).
% mdparagraph:P532:end

% output: EX-XII-XIII-37-03-O001
% mdparagraph:P533:begin
Write one complete Old English translation, keeping the printed controls and the supplied expressions available.
% mdparagraph:P533:end

\end{minipage}\par\addvspace{4pt}

\par\begin{minipage}{\linewidth}\raggedright
% mdparagraph:P534:begin
\subsection*{Cook 37.4}
% mdparagraph:P534:end

% item: EX-XII-XIII-37-04
% mdparagraph:P535:begin
Praise him, ye heavens of heavens, and ye waters that are above (ofer, 172) the heavens.
% mdparagraph:P535:end

% output: EX-XII-XIII-37-04-O001
% mdparagraph:P536:begin
Write one complete Old English translation, keeping the printed controls and the supplied expressions available.
% mdparagraph:P536:end

\end{minipage}\par\addvspace{4pt}

\par\begin{minipage}{\linewidth}\raggedright
% mdparagraph:P537:begin
\subsection*{Cook 37.5}
% mdparagraph:P537:end

% item: EX-XII-XIII-37-05
% mdparagraph:P538:begin
Let them praise [praise they] the name of the Lord; because he spake (gecweðan), and they were made (wyrcean); he commanded, and they were created (p. 190, l. 19).
% mdparagraph:P538:end

% mdparagraph:P539:begin
\textbf{Construction.} “Praise they” calls for a third-person wish with an expressed subject. Cook’s page 190, line 19 begins the said/\allowbreak{}made example; commanded/\allowbreak{}created follows on line 20. Use the passive auxiliary with the appropriate predicate participles.
% mdparagraph:P539:end

% output: EX-XII-XIII-37-05-O001
% mdparagraph:P540:begin
Write one complete Old English translation, keeping the printed controls and the supplied expressions available.
% mdparagraph:P540:end

\end{minipage}\par\addvspace{4pt}

\par\begin{minipage}{\linewidth}\raggedright
% mdparagraph:P541:begin
\subsection*{Cook 37.6}
% mdparagraph:P541:end

% item: EX-XII-XIII-37-06
% mdparagraph:P542:begin
He hath appointed them for ever and ever; he (om. hath) made (sęttan) a decree (bebod, 47), and (om. it) shall not pass (gewītan).
% mdparagraph:P542:end

\end{minipage}\par\addvspace{4pt}

% teaching: NEW-U17-GRAM-U17-H-37-06
% mdparagraph:P543:begin
\textbf{Construction.} Omit a separate word for hath. In the final coordinated predicate, the decree is the understood subject; God is not the thing that passes away. The supplied present gewīt has future reference.
% mdparagraph:P543:end

% teaching: NEW-U17-VOC-GEWITAN-U17-DECREE-PASS
% mdparagraph:P544:begin
\textbf{Supplied:} hē sętte bebod and ne gewīt. — He established a decree, and it will not pass away. The second predicate has the understood decree as subject. Keep the omitted it unexpressed and preserve the future reference of present gewīt.
% mdparagraph:P544:end

% mdparagraph:P545:begin
\textbf{Supplied:} ā — ever; always.
% mdparagraph:P545:end

% mdparagraph:P546:begin
\textbf{Supplied:} ā — ever; always.
% mdparagraph:P546:end

% mdparagraph:P547:begin
\textbf{Supplied:} ā — ever; always.
% mdparagraph:P547:end

% mdparagraph:P548:begin
\textbf{Supplied:} gewīt — pass away; cease to hold; indicative, present, third person, singular.
% mdparagraph:P548:end

% output: EX-XII-XIII-37-06-O001
% mdparagraph:P549:begin
Write one complete Old English translation, keeping the printed controls and the supplied expressions available.
% mdparagraph:P549:end

\par\begin{minipage}{\linewidth}\raggedright
% mdparagraph:P550:begin
\subsection*{Cook 37.7}
% mdparagraph:P550:end

% item: EX-XII-XIII-37-07
% mdparagraph:P551:begin
Praise the Lord from the earth, ye dragons (draca, 53), and all deeps (nēowolnes).
% mdparagraph:P551:end

% mdparagraph:P552:begin
\textbf{Supplied:} nēowolnessa — abyss; the deep; nominative, plural.
% mdparagraph:P552:end

% output: EX-XII-XIII-37-07-O001
% mdparagraph:P553:begin
Write one complete Old English translation, keeping the printed controls and the supplied expressions available.
% mdparagraph:P553:end

\end{minipage}\par\addvspace{4pt}

\par\begin{minipage}{\linewidth}\raggedright
% mdparagraph:P554:begin
\subsection*{Cook 37.8}
% mdparagraph:P554:end

% item: EX-XII-XIII-37-08
% mdparagraph:P555:begin
Fire and hail (hægel, 43); snow (snāw, 43) and ice (īs, 47); spirits of the storms (ȳst, 51. b) which fulfil (dōn) his word;
% mdparagraph:P555:end

% mdparagraph:P556:begin
\textbf{Construction.} The spirits of the storms fulfil his word; keep this fulfilment sense of dōn distinct from a general act of making.
% mdparagraph:P556:end

\end{minipage}\par\addvspace{4pt}

% teaching: U17-SUP-HAEGEL-NOM-SG
% teaching: U17-SUP-SNAW-NOM-SG
% teaching: U17-SUP-IS-ICE-NOM-SG
% mdparagraph:P557:begin
\textbf{Supplied:} hægel — hail; nominative, masculine, singular.
% mdparagraph:P557:end

% mdparagraph:P558:begin
\textbf{Supplied:} snāw — snow; nominative, masculine, singular.
% mdparagraph:P558:end

% mdparagraph:P559:begin
\textbf{Supplied:} īs — ice; nominative, neuter, singular.
% mdparagraph:P559:end

% output: EX-XII-XIII-37-08-O001
% mdparagraph:P560:begin
Write one complete Old English translation, keeping the printed controls and the supplied expressions available.
% mdparagraph:P560:end

\par\begin{minipage}{\linewidth}\raggedright
% mdparagraph:P561:begin
\subsection*{Cook 37.9}
% mdparagraph:P561:end

% item: EX-XII-XIII-37-09
% mdparagraph:P562:begin
Mountains, and all hills (hyll, 43); fruitful (wæstmbǣre, 59) trees, and all cedars;
% mdparagraph:P562:end

% mdparagraph:P563:begin
\textbf{Construction.} The described trees are addressed directly; use the weak plural adjective form. The quantifier all keeps its own agreement.
% mdparagraph:P563:end

% output: EX-XII-XIII-37-09-O001
% mdparagraph:P564:begin
Write one complete Old English translation, keeping the printed controls and the supplied expressions available.
% mdparagraph:P564:end

\end{minipage}\par\addvspace{4pt}

\par\begin{minipage}{\linewidth}\raggedright
% mdparagraph:P565:begin
\subsection*{Cook 37.10}
% mdparagraph:P565:end

% item: EX-XII-XIII-37-10
% mdparagraph:P566:begin
(Wild) beasts, and all cattle; serpents, and winged (gefiðred) fowl;
% mdparagraph:P566:end

% mdparagraph:P567:begin
\textbf{Construction.} “(Wild) beasts” specifies the intended animal sense. The winged birds are directly addressed and take a weak plural adjective.
% mdparagraph:P567:end

\end{minipage}\par\addvspace{4pt}

% teaching: U17-SUP-GEFITHRED-NOM-PL-WEAK
% mdparagraph:P568:begin
\textbf{Supplied:} gefiðredan — feathered; nominative, masculine, plural, weak.
% mdparagraph:P568:end

% output: EX-XII-XIII-37-10-O001
% mdparagraph:P569:begin
Write one complete Old English translation, keeping the printed controls and the supplied expressions available.
% mdparagraph:P569:end

\par\begin{minipage}{\linewidth}\raggedright
% mdparagraph:P570:begin
\subsection*{Cook 37.11}
% mdparagraph:P570:end

% item: EX-XII-XIII-37-11
% mdparagraph:P571:begin
Kings of the earth, and all peoples; princes, and all judges of the earth;
% mdparagraph:P571:end

% mdparagraph:P572:begin
\textbf{Supplied:} folc — peoples; distinct peoples or national groups; nominative, plural.
% mdparagraph:P572:end

% output: EX-XII-XIII-37-11-O001
% mdparagraph:P573:begin
Write one complete Old English translation, keeping the printed controls and the supplied expressions available.
% mdparagraph:P573:end

\end{minipage}\par\addvspace{4pt}

\par\begin{minipage}{\linewidth}\raggedright
% mdparagraph:P574:begin
\subsection*{Cook 37.12}
% mdparagraph:P574:end

% item: EX-XII-XIII-37-12
% mdparagraph:P575:begin
Young (om. men), and maidens (fǣmne); old (om. men), with children;
% mdparagraph:P575:end

% mdparagraph:P576:begin
\textbf{Construction.} Omit men but retain the people meant by young and old. Their weak masculine plural adjectives carry the direct address.
% mdparagraph:P576:end

% output: EX-XII-XIII-37-12-O001
% mdparagraph:P577:begin
Write one complete Old English translation, keeping the printed controls and the supplied expressions available.
% mdparagraph:P577:end

\end{minipage}\par\addvspace{4pt}

\par\begin{minipage}{\linewidth}\raggedright
% mdparagraph:P578:begin
\subsection*{Cook 37.13}
% mdparagraph:P578:end

% item: EX-XII-XIII-37-13
% mdparagraph:P579:begin
Let them praise the name of the Lord; for the name of him alone is exalted (ūpāhębban), his glory (wuldor) is above (ofer, 172) the earth and heaven.
% mdparagraph:P579:end

% mdparagraph:P580:begin
\textbf{Supplied:} ānes — alone; genitive, masculine, singular, strong.
% mdparagraph:P580:end

% output: EX-XII-XIII-37-13-O001
% mdparagraph:P581:begin
Write one complete Old English translation, keeping the printed controls and the supplied expressions available.
% mdparagraph:P581:end

\end{minipage}\par\addvspace{4pt}

\par\begin{minipage}{\linewidth}\raggedright
% mdparagraph:P582:begin
\subsection*{Cook 37.14}
% mdparagraph:P582:end

% item: EX-XII-XIII-37-14
% mdparagraph:P583:begin
He exalteth the horn (horn, 43) of his people, the praise of all his saints (hālga, 53), of the children (bearn) of Israel (Israel), of the people drawing nigh (om. unto) him.
% mdparagraph:P583:end

\end{minipage}\par\addvspace{4pt}

% teaching: NEW-U17-GRAM-U17-H-37-14
% teaching: U17-SUP-ISRAEL-DAT-SG
% mdparagraph:P584:begin
\textbf{Supplied:} horn his folces, lof ealra his hālgena, bearna Israeles, folces him nēalǣcendes. — The horn of his people, the praise of all his saints, of the children of Israel, of a people approaching him. Horn is the accusative thing exalted and represents strength or dignity; lof is appositional praise. The successive genitives identify the people concerned. Hālgena belongs to the noun hālga; folces is singular collective, nēalǣcendes agrees with it, and him is the dative person approached. Use the complete phrase as supplied.
% mdparagraph:P584:end

% teaching: U17-SUP-ISRAEL-GEN-SG
% mdparagraph:P585:begin
\textbf{Supplied:} Israeles — Israel; genitive, singular.
% mdparagraph:P585:end

% mdparagraph:P586:begin
\textbf{Supplied:} nēalǣcendes — approach; draw near; present participle, genitive, neuter, singular, strong.
% mdparagraph:P586:end

% output: EX-XII-XIII-37-14-O001
% mdparagraph:P587:begin
Write one complete Old English translation, keeping the printed controls and the supplied expressions available.
% mdparagraph:P587:end

\reserve{7\baselineskip}
% mdparagraph:P588:begin
\section*{Cook 38. The queen’s visit}
% mdparagraph:P588:end

% mdparagraph:P589:begin
Follow the sequence of hearing, arrival, conversation, observation and judgment. Comparison and reporting clauses keep their own subjects and times.
% mdparagraph:P589:end

\par\begin{minipage}{\linewidth}\raggedright
% mdparagraph:P590:begin
\subsection*{Cook 38.1}
% mdparagraph:P590:end

% item: EX-XII-XIII-38-01
% mdparagraph:P591:begin
The king made a great ivory (elpenbǣnen, 57. 3) throne, and overlaid (oferwyrcean) it with (mid) the best gold.
% mdparagraph:P591:end

\end{minipage}\par\addvspace{4pt}

% teaching: U17-SUP-ELPENBAENEN-ACC-N-SG
% mdparagraph:P592:begin
\textbf{Supplied:} elpenbǣnen — ivory; made of ivory; accusative, neuter, singular, strong.
% mdparagraph:P592:end

% output: EX-XII-XIII-38-01-O001
% mdparagraph:P593:begin
Write one complete Old English translation, keeping the printed controls and the supplied expressions available.
% mdparagraph:P593:end

\par\begin{minipage}{\linewidth}\raggedright
% mdparagraph:P594:begin
\subsection*{Cook 38.2}
% mdparagraph:P594:end

% item: EX-XII-XIII-38-02
% mdparagraph:P595:begin
There (om.) was a certain queen in those days in the south, called (hātan) Sheba (Saba).
% mdparagraph:P595:end

\end{minipage}\par\addvspace{4pt}

% teaching: NEW-U17-LOCAL-SENSE-ON-U15-SOURCE-DURING
% teaching: NEW-U17-LOCAL-SENSE-BE-U17-SOUTH
% teaching: NEW-U17-GRAM-U17-SUPPLIED-DAYS
% teaching: NEW-U17-GRAM-U17-SUPPLIED-SOUTH
% mdparagraph:P596:begin
\textbf{Construction.} Omit the English dummy there. The complete time and south expressions below remain available.
% mdparagraph:P596:end

% teaching: U17-SUP-SABA-NOM-SG
% mdparagraph:P597:begin
\textbf{Supplied:} on ðǣm dagum — In those days; a complete dative plural time expression.
% mdparagraph:P597:end

% mdparagraph:P598:begin
\textbf{Supplied:} be sūðan, — In the south. Preserve the short vowel of be and use the complete expression.
% mdparagraph:P598:end

% teaching: U17-SUP-SUTHAN-INVARIANT
% mdparagraph:P599:begin
\textbf{Supplied:} on — in those days; during the stated time.
% mdparagraph:P599:end

% mdparagraph:P600:begin
\textbf{Supplied:} be — in the south, only in supplied be sūðan.
% mdparagraph:P600:end

% mdparagraph:P601:begin
\textbf{Supplied:} sūðan — in/\allowbreak{}from the south.
% mdparagraph:P601:end

% mdparagraph:P602:begin
\textbf{Supplied:} Saba — Sheba; nominative, singular.
% mdparagraph:P602:end

% output: EX-XII-XIII-38-02-O001
% mdparagraph:P603:begin
Write one complete Old English translation, keeping the printed controls and the supplied expressions available.
% mdparagraph:P603:end

\par\begin{minipage}{\linewidth}\raggedright
% mdparagraph:P604:begin
\subsection*{Cook 38.3}
% mdparagraph:P604:end

% item: EX-XII-XIII-38-03
% mdparagraph:P605:begin
Then she heard the fame (hlīsa) of Solomon (Salomones), and came from the southern (sūðerne, 59) boundaries to Solomon within Jerusalem, and her camels bore southern spices [herbs], and precious gems, and a quantity (unrīm) of gold (154. a).
% mdparagraph:P605:end

\end{minipage}\par\addvspace{4pt}

% teaching: NEW-U17-GRAM-U17-SUPPLIED-WITHIN-PLACE
% teaching: NEW-U17-GRAM-U17-SUPPLIED-UNRIM-GOLD
% mdparagraph:P606:begin
\textbf{Construction.} Use the bracketed herbs meaning. Cook §154.a concerns the partitive genitive after a quantity: goldes depends on unrīm. The feminine accusative plural adjective for southern follows its taught inflection; its dictionary form in the hint is not automatically the answer form.
% mdparagraph:P606:end

% teaching: U17-SUP-SALOMON-GEN-SG
% teaching: U17-SUP-SUTHERNE-ACC-F-PL
% teaching: U17-SUP-UNRIM-ACC-SG
% mdparagraph:P607:begin
\textbf{Supplied:} binnan Hierusalem, — Within Jerusalem. The foreign place form and the complete prepositional expression are supplied.
% mdparagraph:P607:end

% teaching: U17-SUP-BINNAN-INVARIANT
% teaching: U17-SUP-HIERUSALEM-DAT-SG
% mdparagraph:P608:begin
\textbf{Supplied:} unrīm goldes. — An immense quantity of gold. Genitive singular goldes depends on the quantity noun unrīm.
% mdparagraph:P608:end

% mdparagraph:P609:begin
\textbf{Supplied:} Salomones — Solomon; genitive, singular.
% mdparagraph:P609:end

% teaching: U17-SUP-SUTHERNE-DAT-PL
% mdparagraph:P610:begin
\textbf{Supplied:} sūðernum — southern; dative, neuter, plural, strong.
% mdparagraph:P610:end

% teaching: U17-SUP-SALOMON-DAT-SG
% mdparagraph:P611:begin
\textbf{Supplied:} Salomone — Solomon; dative, singular.
% mdparagraph:P611:end

% mdparagraph:P612:begin
\textbf{Supplied:} binnan — within.
% mdparagraph:P612:end

% mdparagraph:P613:begin
\textbf{Supplied:} Hierusalem — Jerusalem; dative, singular.
% mdparagraph:P613:end

% mdparagraph:P614:begin
\textbf{Supplied:} sūðerne — southern; accusative, feminine, plural, strong.
% mdparagraph:P614:end

% mdparagraph:P615:begin
\textbf{Supplied:} unrīm — an uncounted quantity; abundance; accusative, neuter, singular.
% mdparagraph:P615:end

% output: EX-XII-XIII-38-03-O001
% mdparagraph:P616:begin
Write one complete Old English translation, keeping the printed controls and the supplied expressions available.
% mdparagraph:P616:end

\par\begin{minipage}{\linewidth}\raggedright
% mdparagraph:P617:begin
\subsection*{Cook 38.4}
% mdparagraph:P617:end

% item: EX-XII-XIII-38-04
% mdparagraph:P618:begin
The queen then (ðā) had speech with (wið) Solomon, and said (om. to) him whatsoever she thought in her heart (verb last).
% mdparagraph:P618:end

\end{minipage}\par\addvspace{4pt}

% teaching: NEW-U17-GRAM-U17-H-38-04
% teaching: NEW-U17-LOCAL-SENSE-WITH-U17-CONVERSATION
% mdparagraph:P619:begin
\textbf{Construction.} Omit a separate to with the following object and put the finite verb last in the specified dependent clause. The supplied conversation frame gives wið its local meaning.
% mdparagraph:P619:end

% mdparagraph:P620:begin
\textbf{Supplied:} Sēo cwēn hæfde ðā sprǣce wið Salomon — The queen then had speech with Solomon. Wið identifies the conversational partner. This supplied arrangement selects accusative Salomon; dative Salomone is also supported for this conversation frame.
% mdparagraph:P620:end

% teaching: U17-SUP-SALOMON-ACC-SG
% mdparagraph:P621:begin
\textbf{Supplied:} wið — with, the conversational counterpart in the supplied speech frame.
% mdparagraph:P621:end

% mdparagraph:P622:begin
\textbf{Supplied:} Salomon — Solomon; accusative, singular.
% mdparagraph:P622:end

% output: EX-XII-XIII-38-04-O001
% mdparagraph:P623:begin
Write one complete Old English translation, keeping the printed controls and the supplied expressions available.
% mdparagraph:P623:end

\par\begin{minipage}{\linewidth}\raggedright
% mdparagraph:P624:begin
\subsection*{Cook 38.5}
% mdparagraph:P624:end

% item: EX-XII-XIII-38-05
% mdparagraph:P625:begin
Solomon then taught her, and told (sęcgean) her the meaning of all the words that (87. d) she asked (āscian) him.
% mdparagraph:P625:end

% mdparagraph:P626:begin
\textbf{Construction.} Cook §87.d points to invariant relative ðe. Work out its clause role from the verb and antecedent; its shape does not show a case ending.
% mdparagraph:P626:end

\end{minipage}\par\addvspace{4pt}

% teaching: U17-SUP-SALOMON-NOM-SG
% mdparagraph:P627:begin
\textbf{Supplied:} Salomon — Solomon; nominative, singular.
% mdparagraph:P627:end

% output: EX-XII-XIII-38-05-O001
% mdparagraph:P628:begin
Write one complete Old English translation, keeping the printed controls and the supplied expressions available.
% mdparagraph:P628:end

\par\begin{minipage}{\linewidth}\raggedright
% mdparagraph:P629:begin
\subsection*{Cook 38.6}
% mdparagraph:P629:end

% item: EX-XII-XIII-38-06
% mdparagraph:P630:begin
Then the queen saw the wisdom of Solomon, and the great (mǣre) temple that (87. d) he (om. had) built, and the king’s manifold services.
% mdparagraph:P630:end

% mdparagraph:P631:begin
\textbf{Construction.} Use invariant relative ðe and retain the complete list of what the queen saw. Omit the separately marked had.
% mdparagraph:P631:end

% mdparagraph:P632:begin
\textbf{Supplied:} Salomones — Solomon; genitive, singular.
% mdparagraph:P632:end

% output: EX-XII-XIII-38-06-O001
% mdparagraph:P633:begin
Write one complete Old English translation, keeping the printed controls and the supplied expressions available.
% mdparagraph:P633:end

\end{minipage}\par\addvspace{4pt}

\par\begin{minipage}{\linewidth}\raggedright
% mdparagraph:P634:begin
\subsection*{Cook 38.7}
% mdparagraph:P634:end

% item: EX-XII-XIII-38-07
% mdparagraph:P635:begin
She was so (tō ðan) greatly astonished (ofwundrian, 118), that she had no spirit more (furðor).
% mdparagraph:P635:end

\end{minipage}\par\addvspace{4pt}

% teaching: NEW-U17-LOCAL-SENSE-TO-U17-DEGREE-RESULT
% teaching: NEW-U17-LOCAL-REFERENCE-SE-U17-DEGREE-REFERENCE-THAN
% teaching: NEW-U17-GRAM-U17-SUPPLIED-DEGREE-RESULT
% teaching: U17-SUP-FURTHOR-INVARIANT
% mdparagraph:P636:begin
\textbf{Supplied:} tō ðan swīðe ofwundrod, — So greatly astonished. Tō ðan gives the degree and the following ðæt clause its result. Keep the exact supplied ðan.
% mdparagraph:P636:end

% mdparagraph:P637:begin
\textbf{Supplied:} tō — to that degree, within supplied tō ðan … ðæt.
% mdparagraph:P637:end

% mdparagraph:P638:begin
\textbf{Supplied:} ðan — that degree, in tō ðan.
% mdparagraph:P638:end

% mdparagraph:P639:begin
\textbf{Supplied:} furðor — further; any more.
% mdparagraph:P639:end

% output: EX-XII-XIII-38-07-O001
% mdparagraph:P640:begin
Write one complete Old English translation, keeping the printed controls and the supplied expressions available.
% mdparagraph:P640:end

\par\begin{minipage}{\linewidth}\raggedright
% mdparagraph:P641:begin
\subsection*{Cook 38.8}
% mdparagraph:P641:end

% item: EX-XII-XIII-38-08
% mdparagraph:P642:begin
She said then to (tō) the king, True is the word that (87. d) I heard in my country concerning thee and thy wisdom (wīsdōm); but I would not believe before I (om. my)self saw (196. b) it (verb last).
% mdparagraph:P642:end

% mdparagraph:P643:begin
\textbf{Construction.} Keep self while omitting my. Cook §196.b supports the past subjunctive in the before-clause, viewed from the earlier unwillingness. Put that finite verb last.
% mdparagraph:P643:end

% mdparagraph:P644:begin
\textbf{Supplied:} tō — to, the person addressed in speech.
% mdparagraph:P644:end

\end{minipage}\par\addvspace{4pt}

% teaching: U17-SUP-SOTH-CITATION-NOM-N-SG
% mdparagraph:P645:begin
\textbf{Supplied:} Sōð — true; nominative, neuter, singular, strong.
% mdparagraph:P645:end

% output: EX-XII-XIII-38-08-O001
% mdparagraph:P646:begin
Write one complete Old English translation, keeping the printed controls and the supplied expressions available.
% mdparagraph:P646:end

\par\begin{minipage}{\linewidth}\raggedright
% mdparagraph:P647:begin
\subsection*{Cook 38.9}
% mdparagraph:P647:end

% item: EX-XII-XIII-38-09
% mdparagraph:P648:begin
Greater (micel, 66) is thy wisdom and thy work than the report (hlīsa) was (wǣre) that I heard.
% mdparagraph:P648:end

\end{minipage}\par\addvspace{4pt}

% teaching: NEW-U17-GRAM-U17-H-38-09
% mdparagraph:P649:begin
\textbf{Construction.} The printed wǣre is supplied exactly. It belongs to the report in the comparison; it does not agree with wisdom and works. The whole expression below construes wisdom and its work together.
% mdparagraph:P649:end

% mdparagraph:P650:begin
\textbf{Supplied:} Māra is ðīn wīsdōm and ðīn weorc ðonne se hlīsa wǣre ðe ic hīerde. — Your wisdom and your work are greater than the report I heard. In this supplied arrangement, singular Māra is construes wisdom and its work together. The exact wǣre is third singular past subjunctive in the comparison, referring to the report. Keep that given form available; a defensible plural arrangement of the main clause can also be grammatical.
% mdparagraph:P650:end

% output: EX-XII-XIII-38-09-O001
% mdparagraph:P651:begin
Write one complete Old English translation, keeping the printed controls and the supplied expressions available.
% mdparagraph:P651:end

\par\begin{minipage}{\linewidth}\raggedright
% mdparagraph:P652:begin
\subsection*{Cook 38.10}
% mdparagraph:P652:end

% item: EX-XII-XIII-38-10
% mdparagraph:P653:begin
Blessed (ēadig) are thy vassals and thy servants (ðēowa, 53), that always (simle) stand before (ætforan) thee, and hear thy wisdom.
% mdparagraph:P653:end

% mdparagraph:P654:begin
\textbf{Supplied:} simle — always.
% mdparagraph:P654:end

% mdparagraph:P655:begin
\textbf{Supplied:} ætforan — before; in front of.
% mdparagraph:P655:end

% output: EX-XII-XIII-38-10-O001
% mdparagraph:P656:begin
Write one complete Old English translation, keeping the printed controls and the supplied expressions available.
% mdparagraph:P656:end

\end{minipage}\par\addvspace{4pt}

\reserve{7\baselineskip}
% mdparagraph:P657:begin
\section*{Cook 39. Praise, gifts and a change of subject}
% mdparagraph:P657:end

% mdparagraph:P658:begin
The first four items continue praise and reciprocal gifts. Items 5–7 turn to building and to what eye, ear and heart can apprehend; they are a change of material, not a continuation of the queen’s visit.
% mdparagraph:P658:end

\par\begin{minipage}{\linewidth}\raggedright
% mdparagraph:P659:begin
\subsection*{Cook 39.1}
% mdparagraph:P659:end

% item: EX-XII-XIII-39-01
% mdparagraph:P660:begin
Blessed (blētsian) be the almighty God, who (87. d) selected thee, and placed (gesęttan) thee over (ofer, 172) the kingdom of the Israelites (Israela), that thou (om. mightest) establish (sęttan) judgments and righteousness.
% mdparagraph:P660:end

% mdparagraph:P661:begin
\textbf{Construction.} Use the relative particle for “which”. In the that-clause, omit a separate mightest and express the intended purpose with the subjunctive.
% mdparagraph:P661:end

\end{minipage}\par\addvspace{4pt}

% teaching: U17-SUP-ISRAEL-GEN-PL
% teaching: U17-SUP-AELMIHTIG-NOM-M-SG-WEAK
% mdparagraph:P662:begin
\textbf{Supplied:} ælmihtiga — almighty; nominative, masculine, singular, weak.
% mdparagraph:P662:end

% mdparagraph:P663:begin
\textbf{Supplied:} Israela — Israelites, in the supplied genitive plural; genitive, plural.
% mdparagraph:P663:end

% mdparagraph:P664:begin
\textbf{Supplied:} dōmas — judgment; a judgment made in judging; accusative, plural.
% mdparagraph:P664:end

% output: EX-XII-XIII-39-01-O001
% mdparagraph:P665:begin
Write one complete Old English translation, keeping the printed controls and the supplied expressions available.
% mdparagraph:P665:end

\par\begin{minipage}{\linewidth}\raggedright
% mdparagraph:P666:begin
\subsection*{Cook 39.2}
% mdparagraph:P666:end

% item: EX-XII-XIII-39-02
% mdparagraph:P667:begin
She gave (forgiefan) the king (om. a) hundred and twenty pounds (154. c) of gold (154. a), and (om. a) multitude of precious spices (wyrt) and of precious stones.
% mdparagraph:P667:end

\end{minipage}\par\addvspace{4pt}

% teaching: NEW-U17-GRAM-U17-SUPPLIED-120-GOLD
% teaching: NEW-U17-GRAM-U17-SUPPLIED-MENIGU-QUANTITY
% mdparagraph:P668:begin
\textbf{Construction.} Cook §154.c concerns a genitive after a higher numeral; §154.a concerns quantity plus genitive. Both complete quantity expressions are supplied below. Omit the marked indefinite articles.
% mdparagraph:P668:end

% mdparagraph:P669:begin
\textbf{Supplied:} hund and twēntig punda goldes — One hundred and twenty pounds of gold. Punda is genitive plural after the higher numeral; goldes is genitive singular naming the material. Use the complete quantity phrase.
% mdparagraph:P669:end

% mdparagraph:P670:begin
\textbf{Supplied:} męnigeo dēorwierðra wyrta and dēorwierðra stāna. — A great quantity of costly herbs and precious stones. The quantity noun is accusative singular, followed by both groups in the genitive plural. The complete phrase is supplied.
% mdparagraph:P670:end

% teaching: U17-SUP-MENIGU-ACC-SG-VARIANT
% mdparagraph:P671:begin
\textbf{Supplied:} hund — hundred, within the complete supplied quantity of gold.
% mdparagraph:P671:end

% mdparagraph:P672:begin
\textbf{Supplied:} twēntig — twenty.
% mdparagraph:P672:end

% mdparagraph:P673:begin
\textbf{Supplied:} męnigeo — company; crowd; large quantity; accusative, feminine, singular.
% mdparagraph:P673:end

% output: EX-XII-XIII-39-02-O001
% mdparagraph:P674:begin
Write one complete Old English translation, keeping the printed controls and the supplied expressions available.
% mdparagraph:P674:end

\par\begin{minipage}{\linewidth}\raggedright
% mdparagraph:P675:begin
\subsection*{Cook 39.3}
% mdparagraph:P675:end

% item: EX-XII-XIII-39-03
% mdparagraph:P676:begin
Solomon also gave the queen whatsoever she desired (giernan) of (æt) him; and she returned (gewęndan) again (ongēan) to her country with her retainers.
% mdparagraph:P676:end

% mdparagraph:P677:begin
\textbf{Construction.} The requested thing is genitive with giernan; æt with the dative marks the person asked. Keep the main object’s role separate from that of the relative clause.
% mdparagraph:P677:end

\end{minipage}\par\addvspace{4pt}

% teaching: U17-SUP-ONGEAN-BACK
% mdparagraph:P678:begin
\textbf{Supplied:} Salomon — Solomon; nominative, singular.
% mdparagraph:P678:end

% mdparagraph:P679:begin
\textbf{Supplied:} ongēan — back; again.
% mdparagraph:P679:end

% output: EX-XII-XIII-39-03-O001
% mdparagraph:P680:begin
Write one complete Old English translation, keeping the printed controls and the supplied expressions available.
% mdparagraph:P680:end

\par\begin{minipage}{\linewidth}\raggedright
% mdparagraph:P681:begin
\subsection*{Cook 39.4}
% mdparagraph:P681:end

% item: EX-XII-XIII-39-04
% mdparagraph:P682:begin
Solomon then was magnified (mǣrsian, 118) above (ofer, 166) all earthly kings; and all nations (ðēod) desired (gewilnian) to see [that they might see, 194. b] him, and hear his wisdom.
% mdparagraph:P682:end

\end{minipage}\par\addvspace{4pt}

% teaching: NEW-U17-GRAM-U17-H-39-04
% teaching: NEW-U17-LOCAL-SENSE-OFER-U17-RANK-DATIVE
% mdparagraph:P683:begin
\textbf{Construction.} The bracket replaces the infinitive “to see” with the finite clause “that they might see”, expressing what all nations desired. This desired-content clause takes the subjunctive (Cook §194.b). The rank comparison after ofer is dative as Cook §166 directs.
% mdparagraph:P683:end

% mdparagraph:P684:begin
\textbf{Supplied:} Salomon wæs ðā mǣrsod ofer eallum eorðlicum cyningum; — Solomon was then celebrated above all earthly kings. Ofer here expresses rank and takes the explicitly requested dative complement. Keep that local case choice separate from the accusative rank construction encountered earlier.
% mdparagraph:P684:end

% mdparagraph:P685:begin
\textbf{Supplied:} Salomon — Solomon; nominative, singular.
% mdparagraph:P685:end

% mdparagraph:P686:begin
\textbf{Supplied:} ofer — above in rank, with the explicitly supplied dative complement.
% mdparagraph:P686:end

% output: EX-XII-XIII-39-04-O001
% mdparagraph:P687:begin
Write one complete Old English translation, keeping the printed controls and the supplied expressions available.
% mdparagraph:P687:end

\par\begin{minipage}{\linewidth}\raggedright
% mdparagraph:P688:begin
\subsection*{Cook 39.5}
% mdparagraph:P688:end

% item: EX-XII-XIII-39-05
% mdparagraph:P689:begin
Whosoever buildeth upon (ofer, 166) this foundation (grundweall, 43) gold, or silver, or precious stones, or trees, straw (strēaw, 47) or chaff (ceaf, 47), every one’s (ānes gehwilces) work shall be manifested.
% mdparagraph:P689:end

\end{minipage}\par\addvspace{4pt}

% teaching: U17-SUP-STREAW-ACC-SG
% teaching: U17-SUP-CEAF-ACC-SG
% mdparagraph:P690:begin
\textbf{Supplied:} strēaw — straw; accusative, neuter, singular.
% mdparagraph:P690:end

% mdparagraph:P691:begin
\textbf{Supplied:} ceaf — chaff; accusative, neuter, singular.
% mdparagraph:P691:end

% output: EX-XII-XIII-39-05-O001
% mdparagraph:P692:begin
Write one complete Old English translation, keeping the printed controls and the supplied expressions available.
% mdparagraph:P692:end

\par\begin{minipage}{\linewidth}\raggedright
% mdparagraph:P693:begin
\subsection*{Cook 39.6}
% mdparagraph:P693:end

% item: EX-XII-XIII-39-06
% mdparagraph:P694:begin
If the building (getimbrung, 51. 3) continues (ðurhwunian, 118), and withstands the fire (164. m), then shall the workman receive (om. a) reward (edlēan, 47) from (æt) God for his work (153. e).
% mdparagraph:P694:end

% mdparagraph:P695:begin
\textbf{Construction.} Cook §164.m calls for the dative thing resisted. The received reward is also dative in the selected onfōn frame. Under §153.e the work is a genitive dependent on reward, expressing its grounds.
% mdparagraph:P695:end

% output: EX-XII-XIII-39-06-O001
% mdparagraph:P696:begin
Write one complete Old English translation, keeping the printed controls and the supplied expressions available.
% mdparagraph:P696:end

\end{minipage}\par\addvspace{4pt}

\par\begin{minipage}{\linewidth}\raggedright
% mdparagraph:P697:begin
\subsection*{Cook 39.7}
% mdparagraph:P697:end

% item: EX-XII-XIII-39-07
% mdparagraph:P698:begin
No (183) eye can in this life see, nor any (183) ear (53. 2) hear, nor any (183) man’s heart conceive (āsmēan, 113), the things that God prepares (gearcian, 118) (om. for) those that love him.
% mdparagraph:P698:end

\end{minipage}\par\addvspace{4pt}

% teaching: NEW-U17-GRAM-U17-H-39-07
% mdparagraph:P699:begin
\textbf{Construction.} Cook §183 calls for negative concord. Keep the subjects eye, ear and heart distinct; the first mæg can govern all three following infinitives. Ēare is singular weak neuter (§53.2). The beneficiaries are dative without a separate for.
% mdparagraph:P699:end

% teaching: NEW-U17-VOC-ASMEAN-U17-SOURCE-SENSE
% mdparagraph:P700:begin
\textbf{Supplied:} ne nānes mannes heorte āsmēan — Nor can any man’s heart conceive: the earlier mæg also governs this last infinitive. The heart is the final subject, and the things mentioned are its shared object. The exact printed āsmēan and this frame are supplied.
% mdparagraph:P700:end

% mdparagraph:P701:begin
\textbf{Supplied:} on — in worldly life, temporal/\allowbreak{}situational expression.
% mdparagraph:P701:end

% mdparagraph:P702:begin
\textbf{Supplied:} āsmēan — conceive; plain infinitive under the shared mæg, with the heart as subject.
% mdparagraph:P702:end

% output: EX-XII-XIII-39-07-O001
% mdparagraph:P703:begin
Write one complete Old English translation, keeping the printed controls and the supplied expressions available.
% mdparagraph:P703:end

\reserve{7\baselineskip}
% mdparagraph:P704:begin
\section*{Cook 40. Blessing and reversal}
% mdparagraph:P704:end

% mdparagraph:P705:begin
The prayer moves between praise, remembered acts, descriptions of people and a promise to descendants. Keep the beneficiaries, objects and predicates distinct.
% mdparagraph:P705:end

\par\begin{minipage}{\linewidth}\raggedright
% mdparagraph:P706:begin
\subsection*{Cook 40.1}
% mdparagraph:P706:end

% item: EX-XII-XIII-40-01
% mdparagraph:P707:begin
The angel saluted Mary, and said that she was (194. b) filled (āfyllan, 113) with God’s grace [gift], and that God was with her (preposition last); and she was blessed (blētsian) among women.
% mdparagraph:P707:end

% mdparagraph:P708:begin
\textbf{Construction.} Cook §194.b calls for the subjunctive in the reported content. Use the bracketed gift meaning and put the specified preposition after its complement; its government is unchanged.
% mdparagraph:P708:end

\end{minipage}\par\addvspace{4pt}

% teaching: U17-SUP-MARIA-ACC-SG
% mdparagraph:P709:begin
\textbf{Supplied:} Marian — Mary; accusative, singular.
% mdparagraph:P709:end

% teaching: U17-SUP-GIEFU-DAT-SG-GRACE
% mdparagraph:P710:begin
\textbf{Supplied:} giefe — gift of grace; divine favour; dative, feminine, singular.
% mdparagraph:P710:end

% teaching: U17-SUP-BETWEOX-INVARIANT
% mdparagraph:P711:begin
\textbf{Supplied:} betwēox wīfum — Among women. Wīfum is dative plural after betwēox in this expression. Keep the whole phrase available; its government is supplied here.
% mdparagraph:P711:end

% output: EX-XII-XIII-40-01-O001
% mdparagraph:P712:begin
Write one complete Old English translation, keeping the printed controls and the supplied expressions available.
% mdparagraph:P712:end

\par\begin{minipage}{\linewidth}\raggedright
% mdparagraph:P713:begin
\subsection*{Cook 40.2}
% mdparagraph:P713:end

% item: EX-XII-XIII-40-02
% mdparagraph:P714:begin
My soul (om. doth) magnify (mǣrsian) the Lord, and my spirit (om. hath) rejoiced in God my Saviour.
% mdparagraph:P714:end

\end{minipage}\par\addvspace{4pt}

% teaching: NEW-U17-GRAM-U17-H-40-02
% teaching: NEW-U17-LOCAL-SENSE-ON-U17-REJOICING
% mdparagraph:P715:begin
\textbf{Construction.} Omit separate hath and doth. The supplied rejoice-in expression gives on its local relation to the source of joy.
% mdparagraph:P715:end

% mdparagraph:P716:begin
\textbf{Supplied:} mīn gāst blissode on Gode mīnum Hǣlende. — My spirit rejoiced in God my Saviour. On with the dative names the source or object of joy in this supplied expression.
% mdparagraph:P716:end

% mdparagraph:P717:begin
\textbf{Supplied:} on — in, the source or object of joy in the supplied rejoice-in frame.
% mdparagraph:P717:end

% output: EX-XII-XIII-40-02-O001
% mdparagraph:P718:begin
Write one complete Old English translation, keeping the printed controls and the supplied expressions available.
% mdparagraph:P718:end

\par\begin{minipage}{\linewidth}\raggedright
% mdparagraph:P719:begin
\subsection*{Cook 40.3}
% mdparagraph:P719:end

% item: EX-XII-XIII-40-03
% mdparagraph:P720:begin
For he regarded (gesēon) the humility of his handmaiden (ðīnen, 51. 5); verily (sōðlice), hence forth all generations (cnēoris) shall call (sęcgean) me blessed (ēadig).
% mdparagraph:P720:end

% mdparagraph:P721:begin
\textbf{Construction.} The plural generations are the subject of call. The first-person object has its own agreeing feminine predicate adjective.
% mdparagraph:P721:end

\end{minipage}\par\addvspace{4pt}

% teaching: U17-SUP-HEONAN-FORTH-INVARIANT
% mdparagraph:P722:begin
\textbf{Supplied:} heonanforð — henceforth.
% mdparagraph:P722:end

% output: EX-XII-XIII-40-03-O001
% mdparagraph:P723:begin
Write one complete Old English translation, keeping the printed controls and the supplied expressions available.
% mdparagraph:P723:end

\par\begin{minipage}{\linewidth}\raggedright
% mdparagraph:P724:begin
\subsection*{Cook 40.4}
% mdparagraph:P724:end

% item: EX-XII-XIII-40-04
% mdparagraph:P725:begin
For he that (sē ðe) is mighty hath done (om. to) me great things; and holy is his name.
% mdparagraph:P725:end

% output: EX-XII-XIII-40-04-O001
% mdparagraph:P726:begin
Write one complete Old English translation, keeping the printed controls and the supplied expressions available.
% mdparagraph:P726:end

\end{minipage}\par\addvspace{4pt}

\par\begin{minipage}{\linewidth}\raggedright
% mdparagraph:P727:begin
\subsection*{Cook 40.5}
% mdparagraph:P727:end

% item: EX-XII-XIII-40-05
% mdparagraph:P728:begin
And his mercy is (om. on those) fearing (161) him from (of) generation to (on) generation.
% mdparagraph:P728:end

\end{minipage}\par\addvspace{4pt}

% teaching: NEW-U17-GRAM-U17-H-40-05
% teaching: NEW-U17-LOCAL-SENSE-OF-U17-GENERATION-RANGE
% teaching: NEW-U17-LOCAL-SENSE-ON-U17-GENERATION-RANGE
% mdparagraph:P729:begin
\textbf{Construction.} Cook §161 marks the dative beneficiaries. Omit on those, but retain the people through the participle. Their fearing still takes accusative him. The generation-to-generation expression is supplied as a whole.
% mdparagraph:P729:end

% mdparagraph:P730:begin
\textbf{Supplied:} of cnēorisse on cnēorisse. — From generation to generation. Both cnēorisse forms are dative singular; of marks the starting generation and on the succeeding extent. Use the whole supplied temporal expression.
% mdparagraph:P730:end

% mdparagraph:P731:begin
\textbf{Supplied:} of — from, the starting generation in the supplied temporal range.
% mdparagraph:P731:end

% mdparagraph:P732:begin
\textbf{Supplied:} on — to, the succeeding generation in the supplied temporal range.
% mdparagraph:P732:end

% output: EX-XII-XIII-40-05-O001
% mdparagraph:P733:begin
Write one complete Old English translation, keeping the printed controls and the supplied expressions available.
% mdparagraph:P733:end

\par\begin{minipage}{\linewidth}\raggedright
% mdparagraph:P734:begin
\subsection*{Cook 40.6}
% mdparagraph:P734:end

% item: EX-XII-XIII-40-06
% mdparagraph:P735:begin
He hath put down (āweorpan) the powerful from (of) (om. their) seat, and lifted up (ūpāhębban) the lowly.
% mdparagraph:P735:end

% mdparagraph:P736:begin
\textbf{Construction.} The mighty and the humble are weak masculine plural adjectives naming people. Omit the marked their without deleting the seat.
% mdparagraph:P736:end

% output: EX-XII-XIII-40-06-O001
% mdparagraph:P737:begin
Write one complete Old English translation, keeping the printed controls and the supplied expressions available.
% mdparagraph:P737:end

\end{minipage}\par\addvspace{4pt}

\par\begin{minipage}{\linewidth}\raggedright
% mdparagraph:P738:begin
\subsection*{Cook 40.7}
% mdparagraph:P738:end

% item: EX-XII-XIII-40-07
% mdparagraph:P739:begin
He hath filled the hungering (om. ones) with good (om. things), and the rich (rīce) he hath let go empty.
% mdparagraph:P739:end

% mdparagraph:P740:begin
\textbf{Construction.} The hungry and the rich name people without added ones. Good things uses the familiar plural noun for goods; omit a separate things.
% mdparagraph:P740:end

% output: EX-XII-XIII-40-07-O001
% mdparagraph:P741:begin
Write one complete Old English translation, keeping the printed controls and the supplied expressions available.
% mdparagraph:P741:end

\end{minipage}\par\addvspace{4pt}

\par\begin{minipage}{\linewidth}\raggedright
% mdparagraph:P742:begin
\subsection*{Cook 40.8}
% mdparagraph:P742:end

% item: EX-XII-XIII-40-08
% mdparagraph:P743:begin
He hath accepted his servant (cniht) Israel, and remembered his mercy.
% mdparagraph:P743:end

% mdparagraph:P744:begin
\textbf{Supplied:} Israel — Israel; dative, singular.
% mdparagraph:P744:end

% output: EX-XII-XIII-40-08-O001
% mdparagraph:P745:begin
Write one complete Old English translation, keeping the printed controls and the supplied expressions available.
% mdparagraph:P745:end

\end{minipage}\par\addvspace{4pt}

\par\begin{minipage}{\linewidth}\raggedright
% mdparagraph:P746:begin
\subsection*{Cook 40.9}
% mdparagraph:P746:end

% item: EX-XII-XIII-40-09
% mdparagraph:P747:begin
As he spake (sprecan) to (tō) our fathers, Abraham and his seed, for ever and ever.
% mdparagraph:P747:end

% mdparagraph:P748:begin
\textbf{Supplied:} tō — to, the person addressed in speech.
% mdparagraph:P748:end

\end{minipage}\par\addvspace{4pt}

% teaching: U17-SUP-ABRAHAM-DAT-SG
% mdparagraph:P749:begin
\textbf{Supplied:} Abrahame — Abraham; dative, singular.
% mdparagraph:P749:end

% mdparagraph:P750:begin
\textbf{Supplied:} ā — ever; always.
% mdparagraph:P750:end

% mdparagraph:P751:begin
\textbf{Supplied:} ā — ever; always.
% mdparagraph:P751:end

% mdparagraph:P752:begin
\textbf{Supplied:} ā — ever; always.
% mdparagraph:P752:end

% output: EX-XII-XIII-40-09-O001
% mdparagraph:P753:begin
Write one complete Old English translation, keeping the printed controls and the supplied expressions available.
% mdparagraph:P753:end

\par\begin{minipage}{\linewidth}\raggedright
% mdparagraph:P754:begin
\subsection*{Cook 40.10}
% mdparagraph:P754:end

% item: EX-XII-XIII-40-10
% mdparagraph:P755:begin
Mary dwelt with her about (swilce) three months, and returned (gewęndan) then to her own house.
% mdparagraph:P755:end

\end{minipage}\par\addvspace{4pt}

% teaching: NEW-U17-GRAM-U17-SUPPLIED-THREE-MONTHS
% teaching: U17-SUP-SWILCE-APPROXIMATE
% mdparagraph:P756:begin
\textbf{Supplied:} swilce ðrīe mōnaðas — About three months. Swilce gives approximation and mōnaðas is accusative plural of duration; use this complete supplied phrase.
% mdparagraph:P756:end

% teaching: U17-SUP-MONATH-ACC-PL
% teaching: U17-SUP-MARIA-NOM-SG
% mdparagraph:P757:begin
\textbf{Supplied:} Maria — Mary; nominative, singular.
% mdparagraph:P757:end

% mdparagraph:P758:begin
\textbf{Supplied:} swilce — about; approximately.
% mdparagraph:P758:end

% mdparagraph:P759:begin
\textbf{Supplied:} mōnaðas — months; accusative, masculine, plural.
% mdparagraph:P759:end

% mdparagraph:P760:begin
\textbf{Supplied:} āgnum — own; dative, neuter, singular, strong.
% mdparagraph:P760:end

% output: EX-XII-XIII-40-10-O001
% mdparagraph:P761:begin
Write one complete Old English translation, keeping the printed controls and the supplied expressions available.
% mdparagraph:P761:end

\reserve{7\baselineskip}
% mdparagraph:P762:begin
\section*{Cook 41. Petitions and explanation}
% mdparagraph:P762:end

% mdparagraph:P763:begin
Identify each addressee, the requested act and its recipients. The asking and receiving constructions have their own case patterns.
% mdparagraph:P763:end

\par\begin{minipage}{\linewidth}\raggedright
% mdparagraph:P764:begin
\subsection*{Cook 41.1}
% mdparagraph:P764:end

% item: EX-XII-XIII-41-01
% mdparagraph:P765:begin
God was ever, and ever was his wisdom begotten of (of) him, and the wisdom is his Son, just as mighty as the Father.
% mdparagraph:P765:end

\end{minipage}\par\addvspace{4pt}

% teaching: NEW-U17-GRAM-U17-H-41-01
% teaching: NEW-U17-LOCAL-SENSE-OF-U17-BEGETTING
% mdparagraph:P766:begin
\textbf{Supplied:} ǣfre wæs his wīsdōm of him ācęnned; — His wisdom was always begotten of him. The complete theological relation is supplied: of him gives the origin in this expression. It does not establish a general passive-agent use of of.
% mdparagraph:P766:end

% teaching: U17-SUP-AEFRE-INVARIANT
% mdparagraph:P767:begin
\textbf{Supplied:} ǣfre — ever; always.
% mdparagraph:P767:end

% mdparagraph:P768:begin
\textbf{Supplied:} ǣfre — ever; always.
% mdparagraph:P768:end

% mdparagraph:P769:begin
\textbf{Supplied:} of — of, the origin of wisdom in the supplied begetting statement.
% mdparagraph:P769:end

% output: EX-XII-XIII-41-01-O001
% mdparagraph:P770:begin
Write one complete Old English translation, keeping the printed controls and the supplied expressions available.
% mdparagraph:P770:end

\par\begin{minipage}{\linewidth}\raggedright
% mdparagraph:P771:begin
\subsection*{Cook 41.2}
% mdparagraph:P771:end

% item: EX-XII-XIII-41-02
% mdparagraph:P772:begin
The Spirit of God filleth all the circuit of the world, and he maintaineth all things.
% mdparagraph:P772:end

% output: EX-XII-XIII-41-02-O001
% mdparagraph:P773:begin
Write one complete Old English translation, keeping the printed controls and the supplied expressions available.
% mdparagraph:P773:end

\end{minipage}\par\addvspace{4pt}

\par\begin{minipage}{\linewidth}\raggedright
% mdparagraph:P774:begin
\subsection*{Cook 41.3}
% mdparagraph:P774:end

% item: EX-XII-XIII-41-03
% mdparagraph:P775:begin
He that doeth (wyrcean) my Father’s will, he is my brother, and my mother, and my sister.
% mdparagraph:P775:end

% output: EX-XII-XIII-41-03-O001
% mdparagraph:P776:begin
Write one complete Old English translation, keeping the printed controls and the supplied expressions available.
% mdparagraph:P776:end

\end{minipage}\par\addvspace{4pt}

\par\begin{minipage}{\linewidth}\raggedright
% mdparagraph:P777:begin
\subsection*{Cook 41.4}
% mdparagraph:P777:end

% item: EX-XII-XIII-41-04
% mdparagraph:P778:begin
We brought (bringan) no thing to this world, nor can we take (lǣdan) any (183) thing hence with us.
% mdparagraph:P778:end

% mdparagraph:P779:begin
\textbf{Construction.} Cook §183 calls for negative concord. Retain the specified negative relation through both parts of the sentence.
% mdparagraph:P779:end

\end{minipage}\par\addvspace{4pt}

% teaching: U17-SUP-HEONAN-INVARIANT-VARIANT
% mdparagraph:P780:begin
\textbf{Supplied:} heonon — hence; from here.
% mdparagraph:P780:end

% output: EX-XII-XIII-41-04-O001
% mdparagraph:P781:begin
Write one complete Old English translation, keeping the printed controls and the supplied expressions available.
% mdparagraph:P781:end

\par\begin{minipage}{\linewidth}\raggedright
% mdparagraph:P782:begin
\subsection*{Cook 41.5}
% mdparagraph:P782:end

% item: EX-XII-XIII-41-05
% mdparagraph:P783:begin
Ask (biddan), and (om. it shall) be given you; seek, and ye shall find; knock (cnucian, 118), and it shall be opened (om. to) you.
% mdparagraph:P783:end

\end{minipage}\par\addvspace{4pt}

% teaching: NEW-U17-GRAM-U17-H-41-05
% mdparagraph:P784:begin
\textbf{Construction.} Use plural commands. The supplied impersonal passive expressions take dative addressees; do not insert the omitted it shall.
% mdparagraph:P784:end

% mdparagraph:P785:begin
\textbf{Supplied:} ēow bið seald; — It will be given to you. The addressees are dative plural; singular bið has future reference in this supplied impersonal passive.
% mdparagraph:P785:end

% mdparagraph:P786:begin
\textbf{Supplied:} ēow bið geopenod. — It will be opened to you. The addressees are dative plural; singular bið has future reference in this supplied impersonal passive.
% mdparagraph:P786:end

% output: EX-XII-XIII-41-05-O001
% mdparagraph:P787:begin
Write one complete Old English translation, keeping the printed controls and the supplied expressions available.
% mdparagraph:P787:end

\par\begin{minipage}{\linewidth}\raggedright
% mdparagraph:P788:begin
\subsection*{Cook 41.6}
% mdparagraph:P788:end

% item: EX-XII-XIII-41-06
% mdparagraph:P789:begin
Our Father, thou which (87. d) art in heaven(s), be thy name hallowed.
% mdparagraph:P789:end

% mdparagraph:P790:begin
\textbf{Construction.} Use the indicated plural heavens. In “thou who art”, invariant ðe leaves the verb second singular with the addressed thou.
% mdparagraph:P790:end

% output: EX-XII-XIII-41-06-O001
% mdparagraph:P791:begin
Write one complete Old English translation, keeping the printed controls and the supplied expressions available.
% mdparagraph:P791:end

\end{minipage}\par\addvspace{4pt}

\par\begin{minipage}{\linewidth}\raggedright
% mdparagraph:P792:begin
\subsection*{Cook 41.7}
% mdparagraph:P792:end

% item: EX-XII-XIII-41-07
% mdparagraph:P793:begin
Thy kingdom (rīce) come (tōcuman; verb first); thy will befall (geweorðan) on earth as (swā-swā) in heaven(s) (verb first).
% mdparagraph:P793:end

% mdparagraph:P794:begin
\textbf{Construction.} Put the finite verbs first as directed. The wishes are expressed by present subjunctives. The exact swā-swā comparison remains visible reference help.
% mdparagraph:P794:end

% mdparagraph:P795:begin
\textbf{Supplied:} swā-swā — just as; like.
% mdparagraph:P795:end

% output: EX-XII-XIII-41-07-O001
% mdparagraph:P796:begin
Write one complete Old English translation, keeping the printed controls and the supplied expressions available.
% mdparagraph:P796:end

\end{minipage}\par\addvspace{4pt}

\par\begin{minipage}{\linewidth}\raggedright
% mdparagraph:P797:begin
\subsection*{Cook 41.8}
% mdparagraph:P797:end

% item: EX-XII-XIII-41-08
% mdparagraph:P798:begin
Give (sęllan) us to-day our daily (dæghwāmlīc, 55) bread; and forgive (forgiefan) us our debts (gylt, 43), as we forgive (164) our debtors (gyltend, 43. 6).
% mdparagraph:P798:end

% mdparagraph:P799:begin
\textbf{Construction.} Use a singular command to the addressee. Debts are the thing forgiven; the debtors are dative people forgiven (§164). Preserve the requested forgiveness relation.
% mdparagraph:P799:end

\end{minipage}\par\addvspace{4pt}

% teaching: U17-SUP-TODAEG-INVARIANT
% mdparagraph:P800:begin
\textbf{Supplied:} tōdæg — today.
% mdparagraph:P800:end

% teaching: U17-SUP-HLAF-ACC-SG
% mdparagraph:P801:begin
\textbf{Supplied:} hlāf — bread; accusative, masculine, singular.
% mdparagraph:P801:end

% output: EX-XII-XIII-41-08-O001
% mdparagraph:P802:begin
Write one complete Old English translation, keeping the printed controls and the supplied expressions available.
% mdparagraph:P802:end

\par\begin{minipage}{\linewidth}\raggedright
% mdparagraph:P803:begin
\subsection*{Cook 41.9}
% mdparagraph:P803:end

% item: EX-XII-XIII-41-09
% mdparagraph:P804:begin
And lead thou us not into (on, 172) temptation (costnung), but deliver us from evil.
% mdparagraph:P804:end

\end{minipage}\par\addvspace{4pt}

% teaching: U17-SUP-YFEL-DAT-SG-EVIL
% mdparagraph:P805:begin
\textbf{Supplied:} yfele — evil; dative, neuter, singular.
% mdparagraph:P805:end

% output: EX-XII-XIII-41-09-O001
% mdparagraph:P806:begin
Write one complete Old English translation, keeping the printed controls and the supplied expressions available.
% mdparagraph:P806:end

\par\begin{minipage}{\linewidth}\raggedright
% mdparagraph:P807:begin
\subsection*{Cook 41.10}
% mdparagraph:P807:end

% item: EX-XII-XIII-41-10
% mdparagraph:P808:begin
Be it so.
% mdparagraph:P808:end

% output: EX-XII-XIII-41-10-O001
% mdparagraph:P809:begin
Write one complete Old English translation, keeping the printed controls and the supplied expressions available.
% mdparagraph:P809:end

\end{minipage}\par\addvspace{4pt}

\par\begin{minipage}{\linewidth}\raggedright
% mdparagraph:P810:begin
\subsection*{Cook 41.11}
% mdparagraph:P810:end

% item: EX-XII-XIII-41-11
% mdparagraph:P811:begin
What (hwilc) father will give to his child (164. a) a stone, if it asks (156. b; 159. b) bread of him?
% mdparagraph:P811:end

% mdparagraph:P812:begin
\textbf{Construction.} The recipient of give is dative (§164.a). Ask has an accusative person asked and a genitive thing requested (§156.b; §159.b); the child is the subject of that clause.
% mdparagraph:P812:end

\end{minipage}\par\addvspace{4pt}

% teaching: U17-SUP-HLAF-GEN-SG
% mdparagraph:P813:begin
\textbf{Supplied:} hlāfes — bread; genitive, masculine, singular.
% mdparagraph:P813:end

% output: EX-XII-XIII-41-11-O001
% mdparagraph:P814:begin
Write one complete Old English translation, keeping the printed controls and the supplied expressions available.
% mdparagraph:P814:end

\par\begin{minipage}{\linewidth}\raggedright
% mdparagraph:P815:begin
\subsection*{Cook 41.12}
% mdparagraph:P815:end

% item: EX-XII-XIII-41-12
% mdparagraph:P816:begin
God wrote (āwrītan) the old law (ǣ) with his finger (finger, 43) on the stone tables (tabele, 53).
% mdparagraph:P816:end

\end{minipage}\par\addvspace{4pt}

% teaching: U17-SUP-AE-LAW-ACC-SG
% mdparagraph:P817:begin
\textbf{Supplied:} ǣ — law; accusative, feminine, singular.
% mdparagraph:P817:end

% teaching: U17-SUP-STANEN-DAT-F-PL
% mdparagraph:P818:begin
\textbf{Supplied:} stānenum — made of stone; dative, feminine, plural, weak.
% mdparagraph:P818:end

% output: EX-XII-XIII-41-12-O001
% mdparagraph:P819:begin
Write one complete Old English translation, keeping the printed controls and the supplied expressions available.
% mdparagraph:P819:end

\reserve{7\baselineskip}
% mdparagraph:P820:begin
\section*{Cook 42. Vision and figurative description}
% mdparagraph:P820:end

% mdparagraph:P821:begin
The longer vision sentences contain nested descriptions and source-specific number expressions. Their complete supplied phrases remain available. Separate the finite clauses from participles that describe perceived people.
% mdparagraph:P821:end

\par\begin{minipage}{\linewidth}\raggedright
% mdparagraph:P822:begin
\subsection*{Cook 42.1}
% mdparagraph:P822:end

% item: EX-XII-XIII-42-01
% mdparagraph:P823:begin
The young man fell in love with the girl (p. 175, l. 20; cf. mīn lufu befielð on ðē, ‘my love falls on thee,’ ‘I fall in love with thee’).
% mdparagraph:P823:end

\end{minipage}\par\addvspace{4pt}

% teaching: NEW-U17-GRAM-U17-H-42-01
% teaching: NEW-U17-LOCAL-SENSE-ON-U17-LOVE-GOAL
% teaching: NEW-U17-LOCAL-REFERENCE-BEFEALLAN-U17-QUOTE-PRES-3-SG-REFERENCE
% mdparagraph:P824:begin
\textbf{Construction.} The exact quotation supplies the idiom: love is the subject, and on introduces the girl toward whom it is directed. The supplied past expression retains this structure.
% mdparagraph:P824:end

% mdparagraph:P825:begin
\textbf{Supplied:} Ðæs geongan mannes lufu befēoll on ðæt mǣden. — The young man fell in love with the girl: literally, the young man’s love fell upon the girl. Love is the nominative subject, the young man its genitive possessor, and the girl the accusative complement of on. Use the whole supplied past expression; the quoted present befielð in the prompt remains reference help.
% mdparagraph:P825:end

% teaching: U17-SUP-LUFU-NOM-SG
% teaching: NEW-U17-VOC-BEFEALLAN-U17-SOURCE-SENSE
% mdparagraph:P826:begin
\textbf{Supplied:} lufu — love; nominative, feminine, singular.
% mdparagraph:P826:end

% mdparagraph:P827:begin
\textbf{Supplied:} befēoll — fell upon, with love as the subject and the girl as its goal; indicative, preterite, third person, singular.
% mdparagraph:P827:end

% mdparagraph:P828:begin
\textbf{Supplied:} on — on/\allowbreak{}upon, the target of love in the supplied fall-in-love expression.
% mdparagraph:P828:end

% output: EX-XII-XIII-42-01-O001
% mdparagraph:P829:begin
Write one complete Old English translation, keeping the printed controls and the supplied expressions available.
% mdparagraph:P829:end

\par\begin{minipage}{\linewidth}\raggedright
% mdparagraph:P830:begin
\subsection*{Cook 42.2}
% mdparagraph:P830:end

% item: EX-XII-XIII-42-02
% mdparagraph:P831:begin
I heard the voice of harpers harping with their harps; and they sung a new song before the throne, and before the four beasts; and no man could learn that song except the hundred and forty and four thousand.
% mdparagraph:P831:end

\end{minipage}\par\addvspace{4pt}

% teaching: NEW-U17-GRAM-U17-H-42-02-PART
% teaching: NEW-U17-GRAM-U17-H-42-02-NUMBER
% teaching: NEW-U17-LOCAL-REFERENCE-BUTAN-EXCEPT-U17-BUTON-REFERENCE
% teaching: NEW-U17-LOCAL-REFERENCE-HUNDTEONTIG-U17-REFERENCE
% teaching: NEW-U17-LOCAL-REFERENCE-THUSEND-U17-EXCLUSION-FINAL-O
% teaching: NEW-U17-GRAM-U17-SUPPLIED-144000-EXCLUSION
% mdparagraph:P832:begin
\textbf{Construction.} The voice is singular; the harpers are plural. Both the nested harper phrase and the entire exclusion of the 144,000 remain supplied. Retain nīwne sang and the later ðone sang with their visible support.
% mdparagraph:P832:end

% mdparagraph:P833:begin
\textbf{Supplied:} ðāra hearpera stefne mid hiera hearpum hearpiendra; — The voice of the harpers harping with their harps. Stefne is the accusative voice heard. Hearpera and hearpiendra are genitive plural; mid governs dative plural hearpum. The following they refers to the harpers. This complete nested description is supplied.
% mdparagraph:P833:end

% teaching: NEW-U17-VOC-HEARPIAN-U17-SOURCE-SENSE
% mdparagraph:P834:begin
\textbf{Supplied:} būton ðæt hundtēontig and fēowertig and fēower ðūsendo. — Except the one hundred and forty-four thousand. This entire exclusion phrase is supplied, including collective ðæt and final-o ðūsendo. It is not a dative chain. The number’s internal arrangement and unfamiliar forms require no independent recall.
% mdparagraph:P834:end

% mdparagraph:P835:begin
\textbf{Supplied:} hearpiendra — harping; describing the harpers; present participle, genitive, masculine, plural, strong.
% mdparagraph:P835:end

% teaching: U17-SUP-NIWE-ACC-M-SG
% mdparagraph:P836:begin
\textbf{Supplied:} nīwne — new; accusative, masculine, singular, strong.
% mdparagraph:P836:end

% mdparagraph:P837:begin
\textbf{Supplied:} būton — except, within the whole supplied exclusion phrase.
% mdparagraph:P837:end

% mdparagraph:P838:begin
\textbf{Supplied:} hundtēontig — hundred, within the whole supplied 144,000 expression.
% mdparagraph:P838:end

% mdparagraph:P839:begin
\textbf{Supplied:} fēowertig — forty.
% mdparagraph:P839:end

% mdparagraph:P840:begin
\textbf{Supplied:} ðūsendo — thousand, here within the whole supplied 144,000 expression.
% mdparagraph:P840:end

% output: EX-XII-XIII-42-02-O001
% mdparagraph:P841:begin
Write one complete Old English translation, keeping the printed controls and the supplied expressions available.
% mdparagraph:P841:end

\par\begin{minipage}{\linewidth}\raggedright
% mdparagraph:P842:begin
\subsection*{Cook 42.3}
% mdparagraph:P842:end

% item: EX-XII-XIII-42-03
% mdparagraph:P843:begin
I saw so great (om. a) company (męnigu), as no man can number (gerīman), of all nations and of every tribe (mǣgð), standing before the throne (ðrymsetl) of God, all clothed (gescrȳdan) with white garments (gegierela), and holding palm-twigs (palmtwīg, 47) in their hands.
% mdparagraph:P843:end

\end{minipage}\par\addvspace{4pt}

% teaching: NEW-U17-GRAM-U17-H-42-03
% mdparagraph:P844:begin
\textbf{Construction.} The supplied expression keeps the grammatical singular company and the following descriptions of its many members. The participles describe perceived people; they do not become separate progressive finite verbs.
% mdparagraph:P844:end

% teaching: U17-SUP-MENIGU-ACC-SG
% teaching: NEW-U17-VOC-GERIMAN-U17-SOURCE-SENSE
% mdparagraph:P845:begin
\textbf{Supplied:} swā micle męnigu swā nān mann gerīman ne mæg, of eallum ðēodum and of ǣlcre mǣgðe, standende beforan Godes ðrymsetle, ealle mid hwītum gegierelum gescrȳdde and palmtwīg on hiera handum healdende. — So great a company as no man can number, from all nations and every tribe, standing before God’s throne, all clothed in white garments and holding palm branches in their hands. Męnigu is grammatically feminine singular, while the later plural expressions describe its members. Standende and healdende are nonfinite descriptions, and gescrȳdde is a plural predicate past participle. Keep the whole supplied expression; this particular shift from collective to members is not a general agreement rule.
% mdparagraph:P845:end

% mdparagraph:P846:begin
\textbf{Supplied:} męnigu — company; crowd; large quantity; accusative, feminine, singular.
% mdparagraph:P846:end

% mdparagraph:P847:begin
\textbf{Supplied:} gerīman — number; plain infinitive under mæg.
% mdparagraph:P847:end

% mdparagraph:P848:begin
\textbf{Supplied:} gescrȳdde — clothed; plural predicate past participle describing the company’s members.
% mdparagraph:P848:end

% mdparagraph:P849:begin
\textbf{Supplied:} healdende — holding; present participle describing the company’s members.
% mdparagraph:P849:end

% output: EX-XII-XIII-42-03-O001
% mdparagraph:P850:begin
Write one complete Old English translation, keeping the printed controls and the supplied expressions available.
% mdparagraph:P850:end

\par\begin{minipage}{\linewidth}\raggedright
% mdparagraph:P851:begin
\subsection*{Cook 42.4}
% mdparagraph:P851:end

% item: EX-XII-XIII-42-04
% mdparagraph:P852:begin
And they sang with (om. a) loud (hlūd) voice, Be salvation (om. to) our God which sitteth upon (ofer, 166) his throne.
% mdparagraph:P852:end

% output: EX-XII-XIII-42-04-O001
% mdparagraph:P853:begin
Write one complete Old English translation, keeping the printed controls and the supplied expressions available.
% mdparagraph:P853:end

\end{minipage}\par\addvspace{4pt}

\par\begin{minipage}{\linewidth}\raggedright
% mdparagraph:P854:begin
\subsection*{Cook 42.5}
% mdparagraph:P854:end

% item: EX-XII-XIII-42-05
% mdparagraph:P855:begin
And all (om. the) angels stood round about [in circuit of] his throne, and bowed (ālūtan, 103) to God, thus saying,
% mdparagraph:P855:end

\end{minipage}\par\addvspace{4pt}

% teaching: NEW-U17-GRAM-U17-SUPPLIED-ACCOMPANYING-SPEECH
% mdparagraph:P856:begin
\textbf{Construction.} Omit the marked article and retain the supplied “in circuit of” arrangement. Possessive his is not an article.
% mdparagraph:P856:end

% mdparagraph:P857:begin
\textbf{Supplied:} ðus cweðende — Thus saying. The present participle accompanies the angels in 42.5 or the soldiers in 42.8 and introduces their words; it is not another finite verb. Keep this whole accompanying-speech expression available.
% mdparagraph:P857:end

% output: EX-XII-XIII-42-05-O001
% mdparagraph:P858:begin
Write one complete Old English translation, keeping the printed controls and the supplied expressions available.
% mdparagraph:P858:end

\par\begin{minipage}{\linewidth}\raggedright
% mdparagraph:P859:begin
\subsection*{Cook 42.6}
% mdparagraph:P859:end

% item: EX-XII-XIII-42-06
% mdparagraph:P860:begin
Be (om. to) our God blessing and brightness, wisdom and thanksgiving (ðancung), dignity and strength, for ever and ever.
% mdparagraph:P860:end

\end{minipage}\par\addvspace{4pt}

% teaching: NEW-U17-GRAM-U17-H-42-06
% mdparagraph:P861:begin
\textbf{Construction.} The entire singular attribution of blessing and praise is supplied. It is a local formula, not a rule that every coordinated subject takes a singular verb.
% mdparagraph:P861:end

% mdparagraph:P862:begin
\textbf{Supplied:} Ūrum Gode sīe blētsung and beorhtnes, wīsdōm and ðancung, weorðscipe and stręngðu, ā ā ā. — To our God be blessing and brightness, wisdom and thanksgiving, honour and strength, for ever. The recipient is dative and sīe is the singular optative. The whole coordinated formula, including ā ā ā, is supplied.
% mdparagraph:P862:end

% teaching: U17-SUP-STRENGTHU-NOM-SG
% mdparagraph:P863:begin
\textbf{Supplied:} stręngðu — strength; nominative, feminine, singular.
% mdparagraph:P863:end

% mdparagraph:P864:begin
\textbf{Supplied:} ā — ever; always.
% mdparagraph:P864:end

% mdparagraph:P865:begin
\textbf{Supplied:} ā — ever; always.
% mdparagraph:P865:end

% mdparagraph:P866:begin
\textbf{Supplied:} ā — ever; always.
% mdparagraph:P866:end

% output: EX-XII-XIII-42-06-O001
% mdparagraph:P867:begin
Write one complete Old English translation, keeping the printed controls and the supplied expressions available.
% mdparagraph:P867:end

\par\begin{minipage}{\linewidth}\raggedright
% mdparagraph:P868:begin
\subsection*{Cook 42.7}
% mdparagraph:P868:end

% item: EX-XII-XIII-42-07
% mdparagraph:P869:begin
Foxes (fox, 43) have holes (hol, 47), and the birds of heaven nests (nest, 47); but (sōðlice) the Son of Man hath not where (hwǣr) he (om. may) lay (āhyldan, 113) his head.
% mdparagraph:P869:end

% mdparagraph:P870:begin
\textbf{Construction.} Omit a separate may. The subjunctive in the indefinite relative clause expresses a place where one might lay one’s head.
% mdparagraph:P870:end

\end{minipage}\par\addvspace{4pt}

% teaching: U17-SUP-HWAER-INVARIANT
% mdparagraph:P871:begin
\textbf{Supplied:} hwǣr — where.
% mdparagraph:P871:end

% output: EX-XII-XIII-42-07-O001
% mdparagraph:P872:begin
Write one complete Old English translation, keeping the printed controls and the supplied expressions available.
% mdparagraph:P872:end

\par\begin{minipage}{\linewidth}\raggedright
% mdparagraph:P873:begin
\subsection*{Cook 42.8}
% mdparagraph:P873:end

% item: EX-XII-XIII-42-08
% mdparagraph:P874:begin
The soldiers platted (windan, 104) (om. a) crown of (of) thorns (ðorn, 43), and placed on (172) his head, and a reed (hrēod, 47) in his right (om. hand), and bowed their knee before him, thus saying, Hale be thou, King of the Jews (Iudea).
% mdparagraph:P874:end

% mdparagraph:P875:begin
\textbf{Construction.} Cook §172 specifies accusative complements in the marked on phrases, including placement on the head and spitting upon a person. Do not reduce all these meanings to one physical-motion rule. Omit the marked hand.
% mdparagraph:P875:end

\end{minipage}\par\addvspace{4pt}

% teaching: U17-SUP-IUDEA-GEN-PL
% teaching: U17-SUP-HAL-NOM-M-SG
% mdparagraph:P876:begin
\textbf{Supplied:} Hāl — hale; well, in a greeting; nominative, masculine, singular, strong.
% mdparagraph:P876:end

% mdparagraph:P877:begin
\textbf{Supplied:} Iudea — Jews, in the supplied genitive plural; genitive, plural.
% mdparagraph:P877:end

% mdparagraph:P878:begin
\textbf{Supplied:} ðus cweðende — Thus saying. The present participle accompanies the angels in 42.5 or the soldiers in 42.8 and introduces their words; it is not another finite verb. Keep this whole accompanying-speech expression available.
% mdparagraph:P878:end

% output: EX-XII-XIII-42-08-O001
% mdparagraph:P879:begin
Write one complete Old English translation, keeping the printed controls and the supplied expressions available.
% mdparagraph:P879:end

\par\begin{minipage}{\linewidth}\raggedright
% mdparagraph:P880:begin
\subsection*{Cook 42.9}
% mdparagraph:P880:end

% item: EX-XII-XIII-42-09
% mdparagraph:P881:begin
And they spit (spǣtan, 113) on (on, 172) him, and took (om. a) reed and smote (bēatan) his head.
% mdparagraph:P881:end

% mdparagraph:P882:begin
\textbf{Construction.} Keep the actions coordinated and identify their subject before choosing each past form.
% mdparagraph:P882:end

% output: EX-XII-XIII-42-09-O001
% mdparagraph:P883:begin
Write one complete Old English translation, keeping the printed controls and the supplied expressions available.
% mdparagraph:P883:end

\end{minipage}\par\addvspace{4pt}

\reserve{7\baselineskip}
% mdparagraph:P884:begin
\section*{Cook 43. Mixed composition}
% mdparagraph:P884:end

% mdparagraph:P885:begin
Apply the same distinctions to comparisons, requests, reports and descriptions. The change of subject does not change the need to identify each clause’s participants and complements.
% mdparagraph:P885:end

\par\begin{minipage}{\linewidth}\raggedright
% mdparagraph:P886:begin
\subsection*{Cook 43.1}
% mdparagraph:P886:end

% item: EX-XII-XIII-43-01
% mdparagraph:P887:begin
All [each] flesh is grass, and the glory (wuldor) of the flesh is like [as if] the blossom of herbs.
% mdparagraph:P887:end

\end{minipage}\par\addvspace{4pt}

% teaching: NEW-U17-GRAM-U17-H-43-01
% mdparagraph:P888:begin
\textbf{Construction.} The bracketed each and as if are part of the target. Preserve both images, then distinguish a copular identification from an elliptical comparison.
% mdparagraph:P888:end

% mdparagraph:P889:begin
\textbf{Supplied:} Ǣlc flǣsc is gærs, and ðæs flǣsces wuldor is swilce ðāra wyrta blōstma. — Every flesh is grass, and the glory of the flesh is as the blossom of plants. Retain the flesh/\allowbreak{}grass identification and the separate glory/\allowbreak{}blossom comparison. Swilce introduces an elliptical comparison here; no separate finite counterfactual verb is required.
% mdparagraph:P889:end

% teaching: U17-SUP-SWILCE-AS-IF
% mdparagraph:P890:begin
\textbf{Supplied:} swilce — like; as if.
% mdparagraph:P890:end

% output: EX-XII-XIII-43-01-O001
% mdparagraph:P891:begin
Write one complete Old English translation, keeping the printed controls and the supplied expressions available.
% mdparagraph:P891:end

\par\begin{minipage}{\linewidth}\raggedright
% mdparagraph:P892:begin
\subsection*{Cook 43.2}
% mdparagraph:P892:end

% item: EX-XII-XIII-43-02
% mdparagraph:P893:begin
No man heard (183) his voice in (om. the) streets.
% mdparagraph:P893:end

% mdparagraph:P894:begin
\textbf{Construction.} Use negative concord (§183) and omit the marked article.
% mdparagraph:P894:end

% output: EX-XII-XIII-43-02-O001
% mdparagraph:P895:begin
Write one complete Old English translation, keeping the printed controls and the supplied expressions available.
% mdparagraph:P895:end

\end{minipage}\par\addvspace{4pt}

\par\begin{minipage}{\linewidth}\raggedright
% mdparagraph:P896:begin
\subsection*{Cook 43.3}
% mdparagraph:P896:end

% item: EX-XII-XIII-43-03
% mdparagraph:P897:begin
Cry (cleopian), and cease thou not, lift up (āhębban) thy voice like (swā-swā) (om. a) trumpet.
% mdparagraph:P897:end

\end{minipage}\par\addvspace{4pt}

% teaching: NEW-U17-GRAM-U17-H-43-03
% mdparagraph:P898:begin
\textbf{Construction.} Use the command to the stated addressee. The supplied swā-swā phrase is an elliptical comparison; it does not make an actual trumpet the object being lifted.
% mdparagraph:P898:end

% mdparagraph:P899:begin
\textbf{Supplied:} swā-swā bieman. — As a trumpet: the accusative comparison imagines lifting a trumpet. It supplies a standard for the voice being raised, rather than an actual second object of the command. Retain the exact given swā-swā.
% mdparagraph:P899:end

% mdparagraph:P900:begin
\textbf{Supplied:} swā-swā — just as; like.
% mdparagraph:P900:end

% output: EX-XII-XIII-43-03-O001
% mdparagraph:P901:begin
Write one complete Old English translation, keeping the printed controls and the supplied expressions available.
% mdparagraph:P901:end

\par\begin{minipage}{\linewidth}\raggedright
% mdparagraph:P902:begin
\subsection*{Cook 43.4}
% mdparagraph:P902:end

% item: EX-XII-XIII-43-04
% mdparagraph:P903:begin
Lady, whither goest thou?
% mdparagraph:P903:end

% output: EX-XII-XIII-43-04-O001
% mdparagraph:P904:begin
Write one complete Old English translation, keeping the printed controls and the supplied expressions available.
% mdparagraph:P904:end

\end{minipage}\par\addvspace{4pt}

\par\begin{minipage}{\linewidth}\raggedright
% mdparagraph:P905:begin
\subsection*{Cook 43.5}
% mdparagraph:P905:end

% item: EX-XII-XIII-43-05
% mdparagraph:P906:begin
Thou doest ill (yfele).
% mdparagraph:P906:end

\end{minipage}\par\addvspace{4pt}

% teaching: U17-SUP-YFELE-ADVERB-INVARIANT
% mdparagraph:P907:begin
\textbf{Supplied:} yfele — badly; wrongly.
% mdparagraph:P907:end

% output: EX-XII-XIII-43-05-O001
% mdparagraph:P908:begin
Write one complete Old English translation, keeping the printed controls and the supplied expressions available.
% mdparagraph:P908:end

\par\begin{minipage}{\linewidth}\raggedright
% mdparagraph:P909:begin
\subsection*{Cook 43.6}
% mdparagraph:P909:end

% item: EX-XII-XIII-43-06
% mdparagraph:P910:begin
He did ill, because he censured the princess, the king’s daughter.
% mdparagraph:P910:end

% mdparagraph:P911:begin
\textbf{Supplied:} yfele — badly; wrongly.
% mdparagraph:P911:end

% output: EX-XII-XIII-43-06-O001
% mdparagraph:P912:begin
Write one complete Old English translation, keeping the printed controls and the supplied expressions available.
% mdparagraph:P912:end

\end{minipage}\par\addvspace{4pt}

\par\begin{minipage}{\linewidth}\raggedright
% mdparagraph:P913:begin
\subsection*{Cook 43.7}
% mdparagraph:P913:end

% item: EX-XII-XIII-43-07
% mdparagraph:P914:begin
They never learned arithmetic well.
% mdparagraph:P914:end

% output: EX-XII-XIII-43-07-O001
% mdparagraph:P915:begin
Write one complete Old English translation, keeping the printed controls and the supplied expressions available.
% mdparagraph:P915:end

\end{minipage}\par\addvspace{4pt}

\par\begin{minipage}{\linewidth}\raggedright
% mdparagraph:P916:begin
\subsection*{Cook 43.8}
% mdparagraph:P916:end

% item: EX-XII-XIII-43-08
% mdparagraph:P917:begin
You (plur.) will soon know what you do not yet know.
% mdparagraph:P917:end

% mdparagraph:P918:begin
\textbf{Construction.} The addressees are plural. Express their not knowing with ordinary ne and the appropriate finite form of witan.
% mdparagraph:P918:end

% output: EX-XII-XIII-43-08-O001
% mdparagraph:P919:begin
Write one complete Old English translation, keeping the printed controls and the supplied expressions available.
% mdparagraph:P919:end

\end{minipage}\par\addvspace{4pt}

\par\begin{minipage}{\linewidth}\raggedright
% mdparagraph:P920:begin
\subsection*{Cook 43.9}
% mdparagraph:P920:end

% item: EX-XII-XIII-43-09
% mdparagraph:P921:begin
Play (with) ball!
% mdparagraph:P921:end

% mdparagraph:P922:begin
\textbf{Construction.} The parenthetic with calls for mid with the dative instrument. Use the familiar ball paradigm.
% mdparagraph:P922:end

% output: EX-XII-XIII-43-09-O001
% mdparagraph:P923:begin
Write one complete Old English translation, keeping the printed controls and the supplied expressions available.
% mdparagraph:P923:end

\end{minipage}\par\addvspace{4pt}

\par\begin{minipage}{\linewidth}\raggedright
% mdparagraph:P924:begin
\subsection*{Cook 43.10}
% mdparagraph:P924:end

% item: EX-XII-XIII-43-10
% mdparagraph:P925:begin
I eat nothing, although other men eat and are merry.
% mdparagraph:P925:end

\end{minipage}\par\addvspace{4pt}

% teaching: NEW-U17-GRAM-U17-H-43-10
% teaching: NEW-U17-VOC-THE-U17-CONCESSIVE-MARKER
% mdparagraph:P926:begin
\textbf{Construction.} The supplied although-clause has plural present subjunctives. Its meaning is “eat and are merry”, a concession rather than a wish.
% mdparagraph:P926:end

% mdparagraph:P927:begin
\textbf{Supplied:} ðēah ðe ōðre męnn eten and blīðe sīen. — Although other people eat and are merry. Eten and sīen are third plural present subjunctives in the concessive clause. Ðe belongs to the connective ðēah ðe; it is not a relative who or which. The complete clause remains supplied.
% mdparagraph:P927:end

% teaching: U17-SUP-BLITHE-NOM-M-PL
% mdparagraph:P928:begin
\textbf{Supplied:} blīðe — merry; glad; nominative, masculine, plural, strong.
% mdparagraph:P928:end

% output: EX-XII-XIII-43-10-O001
% mdparagraph:P929:begin
Write one complete Old English translation, keeping the printed controls and the supplied expressions available.
% mdparagraph:P929:end

\par\begin{minipage}{\linewidth}\raggedright
% mdparagraph:P930:begin
\subsection*{Cook 43.11}
% mdparagraph:P930:end

% item: EX-XII-XIII-43-11
% mdparagraph:P931:begin
I desire to learn [that I may learn] from you the knowledge that you have.
% mdparagraph:P931:end

% mdparagraph:P932:begin
\textbf{Construction.} The bracket calls for a finite desired-content clause. Its first-person subject and the following plural you keep separate agreement; of with the dative names the source of learning.
% mdparagraph:P932:end

% output: EX-XII-XIII-43-11-O001
% mdparagraph:P933:begin
Write one complete Old English translation, keeping the printed controls and the supplied expressions available.
% mdparagraph:P933:end

\end{minipage}\par\addvspace{4pt}

\par\begin{minipage}{\linewidth}\raggedright
% mdparagraph:P934:begin
\subsection*{Cook 43.12}
% mdparagraph:P934:end

% item: EX-XII-XIII-43-12
% mdparagraph:P935:begin
Choose which of them you (sing.) will.
% mdparagraph:P935:end

% mdparagraph:P936:begin
\textbf{Construction.} Address one person. The relative clause is indefinite: the one you wish to choose. Use the appropriate subjunctive.
% mdparagraph:P936:end

% output: EX-XII-XIII-43-12-O001
% mdparagraph:P937:begin
Write one complete Old English translation, keeping the printed controls and the supplied expressions available.
% mdparagraph:P937:end

\end{minipage}\par\addvspace{4pt}

\par\begin{minipage}{\linewidth}\raggedright
% mdparagraph:P938:begin
\subsection*{Cook 43.13}
% mdparagraph:P938:end

% item: EX-XII-XIII-43-13
% mdparagraph:P939:begin
Within a few hours, two educated and high-born men came toward him.
% mdparagraph:P939:end

\end{minipage}\par\addvspace{4pt}

% teaching: NEW-U17-GRAM-U17-SUPPLIED-WITHIN-TIME
% mdparagraph:P940:begin
\textbf{Supplied:} Binnan fēawum tīdum — Within a few hours. The adjective and time noun are dative plural; keep the whole expression available.
% mdparagraph:P940:end

% teaching: U17-SUP-FEAWE-DAT-F-PL
% mdparagraph:P941:begin
\textbf{Supplied:} Binnan — within.
% mdparagraph:P941:end

% mdparagraph:P942:begin
\textbf{Supplied:} fēawum — few; dative, feminine, plural, strong.
% mdparagraph:P942:end

% teaching: U17-SUP-TOGEANES-INVARIANT
% mdparagraph:P943:begin
\textbf{Supplied:} him tōgēanes — Toward him; to meet him. Him is the dative complement, with tōgēanes placed after it. Keep the whole expression available.
% mdparagraph:P943:end

% output: EX-XII-XIII-43-13-O001
% mdparagraph:P944:begin
Write one complete Old English translation, keeping the printed controls and the supplied expressions available.
% mdparagraph:P944:end

\par\begin{minipage}{\linewidth}\raggedright
% mdparagraph:P945:begin
\subsection*{Cook 43.14}
% mdparagraph:P945:end

% item: EX-XII-XIII-43-14
% mdparagraph:P946:begin
I woo the illustrious maiden.
% mdparagraph:P946:end

% output: EX-XII-XIII-43-14-O001
% mdparagraph:P947:begin
Write one complete Old English translation, keeping the printed controls and the supplied expressions available.
% mdparagraph:P947:end

\end{minipage}\par\addvspace{4pt}

\par\begin{minipage}{\linewidth}\raggedright
% mdparagraph:P948:begin
\subsection*{Cook 43.15}
% mdparagraph:P948:end

% item: EX-XII-XIII-43-15
% mdparagraph:P949:begin
We all say that no one here is (194. b) so acceptable to God as he.
% mdparagraph:P949:end

% mdparagraph:P950:begin
\textbf{Construction.} The denied content takes the selected present subjunctive under Cook §194.b. Keep the local pleasing expression supplied.
% mdparagraph:P950:end

\end{minipage}\par\addvspace{4pt}

% teaching: U17-SUP-GECWEME-NOM-M-SG
% mdparagraph:P951:begin
\textbf{Supplied:} Gode swā gecwēme — So acceptable to God. Gode is dative, naming the person to whom acceptable; gecwēme is the strong masculine nominative singular predicate adjective. The degree word swā stands between them. Keep the whole expression available.
% mdparagraph:P951:end

% output: EX-XII-XIII-43-15-O001
% mdparagraph:P952:begin
Write one complete Old English translation, keeping the printed controls and the supplied expressions available.
% mdparagraph:P952:end

\par\begin{minipage}{\linewidth}\raggedright
% mdparagraph:P953:begin
\subsection*{Cook 43.16}
% mdparagraph:P953:end

% item: EX-XII-XIII-43-16
% mdparagraph:P954:begin
At first he thought that he should be slain [that one should slay him].
% mdparagraph:P954:end

\end{minipage}\par\addvspace{4pt}

% teaching: NEW-U17-GRAM-U17-H-43-16
% mdparagraph:P955:begin
\textbf{Construction.} The bracket supplies an active paraphrase: an indefinite person might slay him. Preserve the distinction between that grammatical subject and the man who held the expectation.
% mdparagraph:P955:end

% mdparagraph:P956:begin
\textbf{Supplied:} ðæt hine man slōge. — That someone should slay him. Man is an indefinite singular subject, and hine is the man who had the expectation. Past subjunctive slōge expresses the contemplated act; it does not assert that he died.
% mdparagraph:P956:end

% output: EX-XII-XIII-43-16-O001
% mdparagraph:P957:begin
Write one complete Old English translation, keeping the printed controls and the supplied expressions available.
% mdparagraph:P957:end

\par\begin{minipage}{\linewidth}\raggedright
% mdparagraph:P958:begin
\subsection*{Cook 43.17}
% mdparagraph:P958:end

% item: EX-XII-XIII-43-17
% mdparagraph:P959:begin
We all fall at (tō) your (sing.) feet.
% mdparagraph:P959:end

% mdparagraph:P960:begin
\textbf{Construction.} The possessor addressed by your is singular. We is the subject of fall and requires a plural indicative verb; the sentence is a statement, not a command.
% mdparagraph:P960:end

% output: EX-XII-XIII-43-17-O001
% mdparagraph:P961:begin
Write one complete Old English translation, keeping the printed controls and the supplied expressions available.
% mdparagraph:P961:end

\end{minipage}\par\addvspace{4pt}

\par\begin{minipage}{\linewidth}\raggedright
% mdparagraph:P962:begin
\subsection*{Cook 43.18}
% mdparagraph:P962:end

% item: EX-XII-XIII-43-18
% mdparagraph:P963:begin
I (om. do) not believe that you are (194. b) my wife.
% mdparagraph:P963:end

% mdparagraph:P964:begin
\textbf{Construction.} Omit the English supporting do. The reported content takes second singular present subjunctive with thou under Cook §194.b.
% mdparagraph:P964:end

% output: EX-XII-XIII-43-18-O001
% mdparagraph:P965:begin
Write one complete Old English translation, keeping the printed controls and the supplied expressions available.
% mdparagraph:P965:end

\end{minipage}\par\addvspace{4pt}

\par\begin{minipage}{\linewidth}\raggedright
% mdparagraph:P966:begin
\subsection*{Cook 43.19}
% mdparagraph:P966:end

% item: EX-XII-XIII-43-19
% mdparagraph:P967:begin
Where is my husband?
% mdparagraph:P967:end

% mdparagraph:P968:begin
\textbf{Supplied:} Hwǣr — where.
% mdparagraph:P968:end

% output: EX-XII-XIII-43-19-O001
% mdparagraph:P969:begin
Write one complete Old English translation, keeping the printed controls and the supplied expressions available.
% mdparagraph:P969:end

\end{minipage}\par\addvspace{4pt}

\par\begin{minipage}{\linewidth}\raggedright
% mdparagraph:P970:begin
\subsection*{Cook 43.20}
% mdparagraph:P970:end

% item: EX-XII-XIII-43-20
% mdparagraph:P971:begin
I lived eleven years in the land of the Egyptians.
% mdparagraph:P971:end

\end{minipage}\par\addvspace{4pt}

% teaching: U17-SUP-ENDLEOFAN-INVARIANT
% mdparagraph:P972:begin
\textbf{Supplied:} endleofan — eleven.
% mdparagraph:P972:end

% teaching: U17-SUP-EGYPTA-GEN-PL
% mdparagraph:P973:begin
\textbf{Supplied:} Ægypta — the Egyptians; genitive, plural.
% mdparagraph:P973:end

% output: EX-XII-XIII-43-20-O001
% mdparagraph:P974:begin
Write one complete Old English translation, keeping the printed controls and the supplied expressions available.
% mdparagraph:P974:end

\reserve{7\baselineskip}
% mdparagraph:P975:begin
\section*{7. Interpreting the images}
% mdparagraph:P975:end

% mdparagraph:P976:begin
Keep the wording of the completed translations beside the following questions. Explain the effect of an image after identifying its literal grammatical structure. Do not replace an entire translated clause with a general statement of its message.
% mdparagraph:P976:end

\reserve{8\baselineskip}
% mdparagraph:P977:begin
\subsection*{A05}
% mdparagraph:P977:end

% item: NEW-U17-A05
% mdparagraph:P978:begin
After translating Cook 43.1, analyse both images. Keep your complete translation beside the English prompt and the supplied vocabulary.
% mdparagraph:P978:end

% output: NEW-U17-A05-O001
% mdparagraph:P979:begin
\itempara{1}{Identify the two things joined by the first copular image and explain its contextual force.}
% mdparagraph:P979:end

% output: NEW-U17-A05-O002
% mdparagraph:P980:begin
\itempara{2}{Identify the second comparison and distinguish its structure from the first.}
% mdparagraph:P980:end

% output: NEW-U17-A05-O003
% mdparagraph:P981:begin
\itempara{3}{Explain why Human life is brief alone is not a complete translation of 43.1.}
% mdparagraph:P981:end

\reserve{8\baselineskip}
% mdparagraph:P982:begin
\subsection*{A06}
% mdparagraph:P982:end

% item: NEW-U17-A06
% mdparagraph:P983:begin
Using Cook 37.14 and 40.9 and the printed lexical help, explain how context narrows the meanings of horn and seed. Retain the actual images in translation.
% mdparagraph:P983:end

% output: NEW-U17-A06-O001
% mdparagraph:P984:begin
\itempara{1}{For 37.14, distinguish the horn image from a report about a literal animal.}
% mdparagraph:P984:end

% output: NEW-U17-A06-O002
% mdparagraph:P985:begin
\itempara{2}{For 40.9, explain what seed denotes and identify whose seed it is.}
% mdparagraph:P985:end

% output: NEW-U17-A06-O003
% mdparagraph:P986:begin
\itempara{3}{State what the familiar inflection alone cannot establish in these two examples.}
% mdparagraph:P986:end

\reserve{7\baselineskip}
% mdparagraph:P987:begin
\section*{8. A musical performance}
% mdparagraph:P987:end

% mdparagraph:P988:begin
After a shipwreck, Apollonius has reached the court of King Arcestrates. At the feast the king’s daughter asks about his misfortunes. The king then tells her to fetch her harp and ease the young man’s sorrow.
% mdparagraph:P988:end

% mdparagraph:P989:begin
The reading is from Cook’s First Book, Reader XII, printed pages 174–175, lines 174.7–175.17 (scan pages 190–191). It begins when the daughter goes out and ends with the audience’s approval. The line references count the physical prose lines, anchored by the original marginal numbers. The text preserves the source’s spelling; line-end word divisions are joined and the original note numbers are identified by page. All help remains available during the reading tasks.
% mdparagraph:P989:end

% source: FB-XII-APOLLONIUS-MUSIC-PERFORMANCE
% mdparagraph:P990:begin
\readingtext{Ðā ēode hēo ūt,\textsuperscript{P174N9} and hēt fęccean hiere hearpan. And sōna swā hēo hearpian ongan, hēo mid wynsumum sange gemęngde ðǣre hearpan swēg. Ðā ongunnon ealle ðā męnn hīe hęrian on hiere swēgcræfte; and Apollonius āna\textsuperscript{P174N10} swīgode. Ðā cwæð se cyning: “Apolloni, nū ðū dēst\textsuperscript{P174N11} yfele, for-ðǣm-ðe ealle męnn hęriað mīne dohtor on hiere swēgcræfte,\textsuperscript{P174N12} and ðū āna hīe, swīgende,\textsuperscript{P174N12a} tǣlst.\textsuperscript{P174N13}” Apollonius cwæð: “Ēalā, ðū gōda cyning, gif ðū mē gelīefst,\textsuperscript{P174N14} ic sęcge ðæt ic ongiete ðæt sōðlice ðīn dohtor gefēoll\textsuperscript{P174N15} on swēgcræft, ac hēo næfð hine nā wel geleornod; ac hāt mē\textsuperscript{P174N16} nū sęllan ðā hearpan, ðonne wāst\textsuperscript{P174N17} ðū nū ðæt ðū gīet nāst.\textsuperscript{P174N17}” Arcestrates se cyning cwæð: “Apolloni, ic oncnāwe sōðlice ðæt ðū eart\textsuperscript{P175N1} on eallum ðingum wel gelǣred.” Ðā hēt se cyning sęllan Apollonie ðā hearpan. Apollonius ðā ūtēode, and hine scrȳdde, and sętte ænne cynehelm upon his hēafod, and nōm ðā hearpan on his hand, and inēode, and swā stōd ðæt se cyning and ealle ðā ymbsittendan wēndon ðæt hē nǣre Apollonius, ac ðæt hē wǣre Apollines,\textsuperscript{P175N2} ðāra hǣðenra god. Ðā wearð stilnes and swīge\textsuperscript{P175N3} geworden innan ðǣre healle. And Apollonius his hearpenægl genōm, and hē ðā hearpestręngas mid cræfte āstyrian ongan, and ðǣre hearpan swēg mid wynsumum sange gemęngde.\textsuperscript{P175N4} And se cyning self, and ealle ðe ðǣr andwearde wǣron, micelre stefne cleopedon and hine hęredon. Æfter ðisum forlēt\textsuperscript{P175N5} Apollonius ðā hearpan, and\textsuperscript{P175N6} plegode, and fela fægerra ðinga\textsuperscript{P175N7} ðǣr forðtēah,\textsuperscript{P175N8} ðe ðǣm folce ungecnāwen wæs and ungewunelic. And him\textsuperscript{P175N9} eallum ðearle līcode ǣlc ðāra ðinga ðe hē forðtēah.}
% mdparagraph:P990:end

\par
\reserve{7\baselineskip}
% mdparagraph:P991:begin
\subsection*{Finding the cited lines}
% mdparagraph:P991:end

% mdparagraph:P992:begin
The passage is set continuously above. These short openings locate the physical source lines used by the questions; words divided at a line end are printed complete here.
% mdparagraph:P992:end

\begin{itemize}
% mdparagraph:P993:begin
\item 174.7 — Ðā ēode hēo ūt
% mdparagraph:P993:end
% mdparagraph:P994:begin
\item 174.8 — sōna swā hēo hearpian
% mdparagraph:P994:end
% mdparagraph:P995:begin
\item 174.9 — gemęngde ðǣre hearpan swēg.
% mdparagraph:P995:end
% mdparagraph:P996:begin
\item 174.10 — męnn hīe hęrian on
% mdparagraph:P996:end
% mdparagraph:P997:begin
\item 174.11 — āna swīgode. Ðā cwæð
% mdparagraph:P997:end
% mdparagraph:P998:begin
\item 174.12 — dēst yfele, for-ðǣm-ðe ealle
% mdparagraph:P998:end
% mdparagraph:P999:begin
\item 174.13 — on hiere swēgcræfte, and
% mdparagraph:P999:end
% mdparagraph:P1000:begin
\item 174.14 — Apollonius cwæð: “Ēalā, ðū
% mdparagraph:P1000:end
% mdparagraph:P1001:begin
\item 174.15 — gelīefst, ic sęcge ðæt
% mdparagraph:P1001:end
% mdparagraph:P1002:begin
\item 174.16 — gefēoll on swēgcræft, ac
% mdparagraph:P1002:end
% mdparagraph:P1003:begin
\item 174.17 — ac hāt mē nū
% mdparagraph:P1003:end
% mdparagraph:P1004:begin
\item 174.18 — ðū gīet nāst.” Arcestrates
% mdparagraph:P1004:end
% mdparagraph:P1005:begin
\item 175.1 — ic oncnāwe sōðlice ðæt
% mdparagraph:P1005:end
% mdparagraph:P1006:begin
\item 175.2 — gelǣred.”
% mdparagraph:P1006:end
% mdparagraph:P1007:begin
\item 175.3 — Ðā hēt se cyning
% mdparagraph:P1007:end
% mdparagraph:P1008:begin
\item 175.4 — Apollonius ðā ūtēode, and
% mdparagraph:P1008:end
% mdparagraph:P1009:begin
\item 175.5 — cynehelm upon his hēafod
% mdparagraph:P1009:end
% mdparagraph:P1010:begin
\item 175.6 — hand, and inēode, and
% mdparagraph:P1010:end
% mdparagraph:P1011:begin
\item 175.7 — ðā ymbsittendan wēndon ðæt
% mdparagraph:P1011:end
% mdparagraph:P1012:begin
\item 175.8 — hē wǣre Apollines, ðāra
% mdparagraph:P1012:end
% mdparagraph:P1013:begin
\item 175.9 — and swīge geworden innan
% mdparagraph:P1013:end
% mdparagraph:P1014:begin
\item 175.10 — his hearpenægl genōm, and
% mdparagraph:P1014:end
% mdparagraph:P1015:begin
\item 175.11 — cræfte āstyrian ongan, and
% mdparagraph:P1015:end
% mdparagraph:P1016:begin
\item 175.12 — wynsumum sange gemęngde. And
% mdparagraph:P1016:end
% mdparagraph:P1017:begin
\item 175.13 — ðe ðǣr andwearde wǣron
% mdparagraph:P1017:end
% mdparagraph:P1018:begin
\item 175.14 — hine hęredon. Æfter ðisum
% mdparagraph:P1018:end
% mdparagraph:P1019:begin
\item 175.15 — and plegode, and fela
% mdparagraph:P1019:end
% mdparagraph:P1020:begin
\item 175.16 — ðǣm folce ungecnāwen wæs
% mdparagraph:P1020:end
% mdparagraph:P1021:begin
\item 175.17 — eallum ðearle līcode ǣlc
% mdparagraph:P1021:end
\end{itemize}

\reserve{7\baselineskip}
% mdparagraph:P1022:begin
\subsection*{Modern help for the reading}
% mdparagraph:P1022:end

% mdparagraph:P1023:begin
The following explanations supply the unfamiliar constructions and contextual meanings. Consult them throughout the reading and its questions. The original notes are retained separately at the end; their Latin parallels require no Latin exercise.
% mdparagraph:P1023:end

\reserve{7\baselineskip}
% mdparagraph:P1024:begin
\subsection*{174.7: the harp is fetched}
% mdparagraph:P1024:end

% teaching: NEW-U17-GRAM-U17-RSUP-01
% teaching: NEW-U17-VOC-FECCEAN-U17-READING
% mdparagraph:P1025:begin
The daughter goes out and has her harp fetched. Hēt governs fęccean, fetch, with the performer of the fetching left unexpressed; hiere gives the harp’s possessor. Do not make the daughter herself the necessary performer of fetch.
% mdparagraph:P1025:end

\reserve{7\baselineskip}
% mdparagraph:P1026:begin
\subsection*{174.8–10: playing, singing and praise}
% mdparagraph:P1026:end

% teaching: NEW-U17-GRAM-U17-RSUP-02
% teaching: U17-SUP-SONA-TEMPORAL-CORRELATIVE
% teaching: U17-SUP-SWA-TEMPORAL-CORRELATIVE
% teaching: NEW-U17-VOC-HEARPIAN-U17-READING
% teaching: NEW-U17-VOC-ONGINNAN-U17-READING
% teaching: NEW-U17-VOC-GEMENGAN-U17-READING
% teaching: U17-SUP-SWEG-ACC-SG
% teaching: U17-SUP-SWEGCRAEFT-DAT-SG
% teaching: NEW-U17-GRAM-U17-READ-SONA-SWA
% mdparagraph:P1027:begin
Sōna swā means as soon as and introduces the complete time clause. Hearpian is the infinitive after began. Gemęngde means mixed: the harp’s sound is accusative and the song follows mid with the dative. On hiere swēgcræfte means with respect to her musical skill. The men’s praise concerns that skill.
% mdparagraph:P1027:end

\reserve{7\baselineskip}
% mdparagraph:P1028:begin
\subsection*{174.11–13: silence and blame}
% mdparagraph:P1028:end

% teaching: NEW-U17-GRAM-U17-RSUP-03
% teaching: U17-SUP-APOLLONIUS-READING-NOM
% teaching: U17-SUP-AN-ALONE
% teaching: NEW-U17-VOC-SWIGIAN-U17-READING
% teaching: U17-SUP-APOLLONIUS-READING-VOC
% teaching: U17-SUP-FOR-THAM-THE-SOURCE-INVARIANT
% mdparagraph:P1029:begin
Āna means alone. Swīgode means remained silent; the source participle swīgende is explained by Cook’s note 12a as swīgiende, being silent. It accompanies the addressed you in the accusation. Dēst yfele means you do wrong: dēst is second singular present indicative of dōn, and yfele is an adverb. Tǣlst means you blame or censure. Keep the source spelling of the participle.
% mdparagraph:P1029:end

\reserve{7\baselineskip}
% mdparagraph:P1030:begin
\subsection*{174.14–16: permission and competence}
% mdparagraph:P1030:end

% teaching: NEW-U17-GRAM-U17-RSUP-04
% teaching: NEW-U17-VOC-GELIEFAN-PERMIT-U17-SOURCE-SENSE
% teaching: NEW-U17-VOC-GEFEALLAN-U17-READING
% teaching: U17-SUP-ON-CHANCED-DOMAIN
% teaching: U17-SUP-SWEGCRAEFT-ACC-SG
% teaching: NEW-U17-VOC-GELEORNIAN-U17-READING
% teaching: NEW-U17-GRAM-U17-READ-ON-MUSIC
% mdparagraph:P1031:begin
Gif ðū mē gelīefst means if you permit me to speak. Gelīefst is second singular present indicative; mē is the dative person granted permission. Speaking is understood from the following ic sęcge, I will say; no speaking-infinitive occurs in the printed clause. This is the permission verb, distinct from believe. Ic ongiete introduces what Apollonius perceives. Gefēoll on swēgcræft means has chanced upon musical skill: on has an accusative complement in this supplied idiom. Masculine hine resumes swēgcræft, and næfð … nā wel geleornod says she has not learned or mastered it well.
% mdparagraph:P1031:end

\reserve{7\baselineskip}
% mdparagraph:P1032:begin
\subsection*{174.17–18: giving and knowing}
% mdparagraph:P1032:end

% teaching: NEW-U17-GRAM-U17-RSUP-05
% teaching: U17-SUP-THONNE-THEN-CONSEQUENCE
% teaching: NEW-U17-VOC-NYTAN-U17-READING
% mdparagraph:P1033:begin
Hāt mē nū sęllan ðā hearpan asks for the harp to be given to the speaker. Mē is the dative recipient of giving, as original note 16 directs; it is not the person commanded to give. Ðonne means then in the consequence. Wāst is you know and nāst is the contracted you do not know; the latter is supplied here without a new conjugation. The present wāst refers to the knowledge that will follow. Ðæt ðū gīet nāst means what you do not yet know.
% mdparagraph:P1033:end

\reserve{7\baselineskip}
% mdparagraph:P1034:begin
\subsection*{175.1–2: the king’s acknowledgment}
% mdparagraph:P1034:end

% teaching: NEW-U17-GRAM-U17-RSUP-06
% teaching: NEW-U17-VOC-ONCNAWAN-U17-READING
% mdparagraph:P1035:begin
Oncnāwe means acknowledge or recognize. The king’s content clause asserts that the addressee is well educated, using indicative eart. Original note 1 points to the certainty exception; the presence of ðæt does not mechanically demand a subjunctive. On eallum ðingum means in all matters, and gelǣred is the familiar adjective learned or educated.
% mdparagraph:P1035:end

\reserve{7\baselineskip}
% mdparagraph:P1036:begin
\subsection*{175.3–6: preparing the demonstration}
% mdparagraph:P1036:end

% teaching: NEW-U17-GRAM-U17-RSUP-07
% teaching: U17-SUP-APOLLONIUS-READING-DAT
% teaching: NEW-U17-VOC-UTGAN-U17-READING
% teaching: NEW-U17-VOC-SCRYDAN-U17-READING
% teaching: U17-SUP-AN-SOURCE-ACC-M-SG
% teaching: U17-SUP-UPON-INVARIANT
% mdparagraph:P1037:begin
The king has the harp given to dative Apollonie. Apollonius then goes out (ūtēode), dresses himself (hine scrȳdde), puts a crown upon his head, takes the harp into his hand and enters (inēode). The supplied source form ænne is masculine accusative singular; retain its unmarked vowel. Upon his hēafod and on his hand describe the destinations of placement and taking.
% mdparagraph:P1037:end

\reserve{7\baselineskip}
% mdparagraph:P1038:begin
\subsection*{175.6–8: appearance and supposition}
% mdparagraph:P1038:end

% teaching: NEW-U17-GRAM-U17-RSUP-08
% teaching: NEW-U17-VOC-YMBSITTAN-U17-READING
% teaching: NEW-U17-VOC-WESAN-U17-READING
% teaching: U17-SUP-APOLLINES-READING-NOM
% teaching: U17-SUP-HAETHEN-GEN-PL
% teaching: U17-SUP-GOD-COMMON-NOM-SG
% mdparagraph:P1039:begin
Swā … ðæt introduces a result: he stood in such a way that the king and those sitting around supposed something about him. Ðā ymbsittendan is a weak plural description naming those people. The coordinated subject takes plural wēndon. The two content clauses have past subjunctives nǣre and wǣre; they present the spectators’ supposition. Apollines is the source’s name-form for Apollo, the god of the heathens. Preserve it without confusing him with Apollonius.
% mdparagraph:P1039:end

\reserve{7\baselineskip}
% mdparagraph:P1040:begin
\subsection*{175.8–9: silence}
% mdparagraph:P1040:end

% teaching: NEW-U17-GRAM-U17-RSUP-09
% teaching: NEW-U17-VOC-WEORTHAN-U17-READING
% teaching: U17-SUP-STILNES-NOM-SG
% teaching: U17-SUP-SWIGE-NOM-SG
% teaching: NEW-U17-VOC-GEWEORTHAN-U17-READING
% teaching: U17-SUP-INNAN-PREPOSITION
% teaching: U17-SUP-HEALL-DAT-SG
% mdparagraph:P1041:begin
Wearð … geworden presents quietness and silence coming about as a whole result; it is not a transitive passive. Retain singular wearð in this supplied expression. Stilnes and swīge are feminine nouns, and innan ðǣre healle means within the hall, with a dative complement.
% mdparagraph:P1041:end

\reserve{7\baselineskip}
% mdparagraph:P1042:begin
\subsection*{175.10–12: the performance}
% mdparagraph:P1042:end

% teaching: NEW-U17-GRAM-U17-RSUP-10
% teaching: U17-SUP-HEARPENAEGL-ACC-SG
% teaching: NEW-U17-VOC-GENIMAN-U17-READING
% teaching: U17-SUP-HEARPESTRENG-ACC-PL
% teaching: U17-SUP-CRAEFT-DAT-SG-SKILL
% teaching: NEW-U17-VOC-CRAEFT-U17-SKILL
% teaching: NEW-U17-VOC-ASTYRIAN-U17-READING
% mdparagraph:P1043:begin
Genōm means took; the object is his hearpenægl, his plectrum or string-plucking implement. Āstyrian ongan means began to set in motion, with the harp strings as accusative object. Mid cræfte means skilfully, with skill. Gemęngde again combines harp sound with song. These supplied senses explain the performance without creating new finite systems for all its verbs.
% mdparagraph:P1043:end

\reserve{7\baselineskip}
% mdparagraph:P1044:begin
\subsection*{175.12–14: the response}
% mdparagraph:P1044:end

% teaching: NEW-U17-GRAM-U17-RSUP-11
% teaching: U17-SUP-ANDWEARD-NOM-M-PL
% teaching: U17-SUP-MICEL-DAT-F-SG-LOUD
% mdparagraph:P1045:begin
The king himself and all those present are the plural actors of cleopedon, cried out, and hęredon, praised. The invariant relative ðe represents the people who were present. Micelre stefne means with a loud voice; it is a dative expression of manner without a printed mid. Hine is the person praised.
% mdparagraph:P1045:end

\reserve{7\baselineskip}
% mdparagraph:P1046:begin
\subsection*{175.14–17: further accomplishments}
% mdparagraph:P1046:end

% teaching: NEW-U17-GRAM-U17-RSUP-12
% teaching: U17-SUP-THES-DAT-N-SG-SOURCE
% teaching: NEW-U17-VOC-FORLAETAN-U17-READING
% teaching: NEW-U17-VOC-PLEGIAN-U17-READING
% teaching: NEW-U17-VOC-FORTHTEON-U17-READING
% teaching: NEW-U17-VOC-UNGECNAWEN-U17-READING
% teaching: U17-SUP-UNGEWUNELIC-NOM-N-SG
% teaching: U17-SUP-THEARLE-INVARIANT
% teaching: NEW-U17-VOC-LICIAN-U17-READING
% mdparagraph:P1047:begin
Forlēt means set aside the harp; plegode means acted or performed; forðtēah means presented or performed the accomplishments. These performance meanings remain supplied. Fela governs genitive plural fægerra ðinga. The following relative construes what was unfamiliar collectively and has singular ungecnāwen wæs and ungewunelic; preserve this source construction without making it a general rule for agreement. In the last sentence, ǣlc is singular subject of līcode, pleased, and him eallum gives the dative people pleased. Ðāra ðinga is the partitive genitive after each; ðearle means greatly.
% mdparagraph:P1047:end

\reserve{7\baselineskip}
% mdparagraph:P1048:begin
\subsection*{Exact word help}
% mdparagraph:P1048:end

% mdparagraph:P1049:begin
These exact source forms can be consulted. The noun and verb endings below identify their particular occurrence; they do not require an additional complete paradigm from memory.
% mdparagraph:P1049:end

% mdparagraph:P1050:begin
\textbf{fęccean} — fetch; plain infinitive after hēt.
% mdparagraph:P1050:end

% mdparagraph:P1051:begin
\textbf{sōna} — as soon as, with swā.
% mdparagraph:P1051:end

% mdparagraph:P1052:begin
\textbf{swā} — as, in sōna swā.
% mdparagraph:P1052:end

% mdparagraph:P1053:begin
\textbf{hearpian} — play the harp; plain infinitive after began.
% mdparagraph:P1053:end

% mdparagraph:P1054:begin
\textbf{ongan} — begin, with the explicitly supplied playing/\allowbreak{}praising action; indicative, preterite, third person, singular.
% mdparagraph:P1054:end

% mdparagraph:P1055:begin
\textbf{gemęngde} — mix sound with song; indicative, preterite, third person, singular.
% mdparagraph:P1055:end

% mdparagraph:P1056:begin
\textbf{swēg} — music; sound; accusative, masculine, singular.
% mdparagraph:P1056:end

% mdparagraph:P1057:begin
\textbf{ongunnon} — begin, with the explicitly supplied playing/\allowbreak{}praising action; indicative, preterite, third person, plural.
% mdparagraph:P1057:end

% mdparagraph:P1058:begin
\textbf{hęrian} — praise; plain infinitive after began.
% mdparagraph:P1058:end

% mdparagraph:P1059:begin
\textbf{swēgcræfte} — musical skill; dative, masculine, singular.
% mdparagraph:P1059:end

% mdparagraph:P1060:begin
\textbf{Apollonius} — Apollonius, the visitor; nominative, masculine, singular.
% mdparagraph:P1060:end

% mdparagraph:P1061:begin
\textbf{āna} — alone; nominative, masculine, singular, weak.
% mdparagraph:P1061:end

% mdparagraph:P1062:begin
\textbf{swīgode} — be silent; indicative, preterite, third person, singular.
% mdparagraph:P1062:end

% mdparagraph:P1063:begin
\textbf{Apolloni} — Apollonius, the visitor; vocative, masculine, singular.
% mdparagraph:P1063:end

% mdparagraph:P1064:begin
\textbf{yfele} — badly; wrongly.
% mdparagraph:P1064:end

% mdparagraph:P1065:begin
\textbf{for-ðǣm-ðe} — because.
% mdparagraph:P1065:end

% mdparagraph:P1066:begin
\textbf{swīgende} — being silent; the supplied present participle.
% mdparagraph:P1066:end

% mdparagraph:P1067:begin
\textbf{gelīefst} — permit; allow; indicative, present, second person, singular.
% mdparagraph:P1067:end

% mdparagraph:P1068:begin
\textbf{gefēoll} — chance upon; indicative, preterite, third person, singular.
% mdparagraph:P1068:end

% mdparagraph:P1069:begin
\textbf{on} — on/\allowbreak{}into, in the chanced-upon idiom.
% mdparagraph:P1069:end

% mdparagraph:P1070:begin
\textbf{swēgcræft} — musical skill; accusative, masculine, singular.
% mdparagraph:P1070:end

% mdparagraph:P1071:begin
\textbf{geleornod} — learned; mastered; past participle.
% mdparagraph:P1071:end

% mdparagraph:P1072:begin
\textbf{ðonne} — then; at that point.
% mdparagraph:P1072:end

% mdparagraph:P1073:begin
\textbf{nāst} — not know; indicative, present, second person, singular.
% mdparagraph:P1073:end

% teaching: U17-SUP-ARCESTRATES-READING-NOM
% mdparagraph:P1074:begin
\textbf{Arcestrates} — Arcestrates, the king; nominative, masculine, singular.
% mdparagraph:P1074:end

% mdparagraph:P1075:begin
\textbf{oncnāwe} — recognize; acknowledge; indicative, present, first person, singular.
% mdparagraph:P1075:end

% mdparagraph:P1076:begin
\textbf{Apollonie} — Apollonius, the visitor; dative, masculine, singular.
% mdparagraph:P1076:end

% mdparagraph:P1077:begin
\textbf{ūtēode} — go out; indicative, preterite, third person, singular.
% mdparagraph:P1077:end

% mdparagraph:P1078:begin
\textbf{scrȳdde} — dress oneself; indicative, preterite, third person, singular.
% mdparagraph:P1078:end

% mdparagraph:P1079:begin
\textbf{ænne} — a; one; accusative, masculine, singular, strong.
% mdparagraph:P1079:end

% mdparagraph:P1080:begin
\textbf{upon} — on; onto.
% mdparagraph:P1080:end

% mdparagraph:P1081:begin
\textbf{ymbsittendan} — people sitting around; present participle, nominative, masculine, plural, weak.
% mdparagraph:P1081:end

% mdparagraph:P1082:begin
\textbf{nǣre} — was not, in the reported supposition; subjunctive, preterite, third person, singular.
% mdparagraph:P1082:end

% mdparagraph:P1083:begin
\textbf{Apollines} — Apollo, the god; not Apollonius; nominative, masculine, singular.
% mdparagraph:P1083:end

% mdparagraph:P1084:begin
\textbf{hǣðenra} — heathen people; genitive, masculine, plural, weak.
% mdparagraph:P1084:end

% mdparagraph:P1085:begin
\textbf{god} — a god; deity; nominative, masculine, singular.
% mdparagraph:P1085:end

% mdparagraph:P1086:begin
\textbf{wearð} — came about, with geworden in the supplied result expression; indicative, preterite, third person, singular.
% mdparagraph:P1086:end

% mdparagraph:P1087:begin
\textbf{stilnes} — stillness; quietness; nominative, feminine, singular.
% mdparagraph:P1087:end

% mdparagraph:P1088:begin
\textbf{swīge} — silence; nominative, feminine, singular.
% mdparagraph:P1088:end

% mdparagraph:P1089:begin
\textbf{geworden} — come about; past participle in the supplied result expression.
% mdparagraph:P1089:end

% mdparagraph:P1090:begin
\textbf{innan} — within; inside.
% mdparagraph:P1090:end

% mdparagraph:P1091:begin
\textbf{healle} — hall; dative, feminine, singular.
% mdparagraph:P1091:end

% mdparagraph:P1092:begin
\textbf{hearpenægl} — plectrum; harp pick; accusative, masculine, singular.
% mdparagraph:P1092:end

% mdparagraph:P1093:begin
\textbf{genōm} — take; indicative, preterite, third person, singular.
% mdparagraph:P1093:end

% mdparagraph:P1094:begin
\textbf{hearpestręngas} — harp strings; accusative, masculine, plural.
% mdparagraph:P1094:end

% mdparagraph:P1095:begin
\textbf{mid} — with; in the manner of the accompanying quality or mental disposition.
% mdparagraph:P1095:end

% mdparagraph:P1096:begin
\textbf{cræfte} — skill; accomplishment; dative, masculine, singular.
% mdparagraph:P1096:end

% mdparagraph:P1097:begin
\textbf{āstyrian} — move; set in motion; plain infinitive after began.
% mdparagraph:P1097:end

% mdparagraph:P1098:begin
\textbf{ongan} — begin, with the explicitly supplied playing/\allowbreak{}praising action; indicative, preterite, third person, singular.
% mdparagraph:P1098:end

% mdparagraph:P1099:begin
\textbf{ðe} — that; which; who.
% mdparagraph:P1099:end

% mdparagraph:P1100:begin
\textbf{andwearde} — present; nominative, masculine, plural, strong.
% mdparagraph:P1100:end

% mdparagraph:P1101:begin
\textbf{micelre} — great; contextually loud; dative, feminine, singular, strong.
% mdparagraph:P1101:end

% mdparagraph:P1102:begin
\textbf{cleopedon} — cry aloud; indicative, preterite, third person, plural.
% mdparagraph:P1102:end

% mdparagraph:P1103:begin
\textbf{ðisum} — this; this event; dative, neuter, singular.
% mdparagraph:P1103:end

% mdparagraph:P1104:begin
\textbf{forlēt} — set aside; indicative, preterite, third person, singular.
% mdparagraph:P1104:end

% mdparagraph:P1105:begin
\textbf{plegode} — act; perform; indicative, preterite, third person, singular.
% mdparagraph:P1105:end

% mdparagraph:P1106:begin
\textbf{forðtēah} — present; perform; indicative, preterite, third person, singular.
% mdparagraph:P1106:end

% mdparagraph:P1107:begin
\textbf{ungecnāwen} — unknown; singular predicate in the supplied collective expression.
% mdparagraph:P1107:end

% mdparagraph:P1108:begin
\textbf{ungewunelic} — unusual; unfamiliar; nominative, neuter, singular, strong.
% mdparagraph:P1108:end

% mdparagraph:P1109:begin
\textbf{ðearle} — greatly; exceedingly.
% mdparagraph:P1109:end

% mdparagraph:P1110:begin
\textbf{līcode} — please; indicative, preterite, third person, singular.
% mdparagraph:P1110:end

% mdparagraph:P1111:begin
\textbf{forðtēah} — present; perform; indicative, preterite, third person, singular.
% mdparagraph:P1111:end

% mdparagraph:P1112:begin
\textbf{Names:} Apollonius is the visitor; Apolloni addresses him, and Apollonie is the dative form for the recipient. Arcestrates is the king. Apollines is Apollo, the god. These exact name-forms are supplied.
% mdparagraph:P1112:end

% mdparagraph:P1113:begin
\textbf{Source distinctions:} for-ðǣm-ðe means because, with the source’s internal hyphens and ǣ. Āna means alone, rather than cardinal one. The lowercase common noun god names a deity here. Hǣðenra names the heathens in the genitive plural after ðāra. Mid cræfte describes manner, whereas mid wynsumum sange supplies the accompanying song. Æfter ðisum means after this, with the source’s single s.
% mdparagraph:P1113:end

\reserve{7\baselineskip}
% mdparagraph:P1114:begin
\subsection*{Reading tasks}
% mdparagraph:P1114:end

% mdparagraph:P1115:begin
Keep the exact passage, the original notes and all modern help open. These questions ask you to relate the supplied forms to one another and to the episode’s argument; they do not test independent recall of source-only forms.
% mdparagraph:P1115:end

\reserve{8\baselineskip}
% mdparagraph:P1116:begin
\subsection*{R01}
% mdparagraph:P1116:end

% item: NEW-U17-R01
% mdparagraph:P1117:begin
Translate the complete assigned reading in the three stated divisions, keeping all local help available.
% mdparagraph:P1117:end

% output: NEW-U17-R01-O001
% mdparagraph:P1118:begin
\itempara{1}{Translate 174.7–174.13, through tǣlst.}
% mdparagraph:P1118:end

% output: NEW-U17-R01-O002
% mdparagraph:P1119:begin
\itempara{2}{Translate 174.14–175.2, through wel gelǣred.}
% mdparagraph:P1119:end

% output: NEW-U17-R01-O003
% mdparagraph:P1120:begin
\itempara{3}{Translate 175.3–175.17, through the final forðtēah.}
% mdparagraph:P1120:end

\reserve{8\baselineskip}
% mdparagraph:P1121:begin
\subsection*{R02}
% mdparagraph:P1121:end

% item: NEW-U17-R02
% mdparagraph:P1122:begin
Identify each referent and the grammatical evidence. Use the supplied source notes; say when context is needed.
% mdparagraph:P1122:end

% output: NEW-U17-R02-O001
% mdparagraph:P1123:begin
\itempara{1}{Give the referent of hēo in 174.7.}
% mdparagraph:P1123:end

% output: NEW-U17-R02-O002
% mdparagraph:P1124:begin
\itempara{2}{Give the referent of hine in 174.16 and its role.}
% mdparagraph:P1124:end

% output: NEW-U17-R02-O003
% mdparagraph:P1125:begin
\itempara{3}{Give the referent of hē in the reported clauses 175.7–8.}
% mdparagraph:P1125:end

% output: NEW-U17-R02-O004
% mdparagraph:P1126:begin
\itempara{4}{Give the referent and case of him eallum in 175.16–17.}
% mdparagraph:P1126:end

\reserve{8\baselineskip}
% mdparagraph:P1127:begin
\subsection*{R03}
% mdparagraph:P1127:end

% item: NEW-U17-R03
% mdparagraph:P1128:begin
Identify the clause relationship in each cited expression and distinguish it from a similar construction elsewhere.
% mdparagraph:P1128:end

% output: NEW-U17-R03-O001
% mdparagraph:P1129:begin
\itempara{1}{Classify sōna swā 174.8 and explain its time relation.}
% mdparagraph:P1129:end

% output: NEW-U17-R03-O002
% mdparagraph:P1130:begin
\itempara{2}{Classify gif 174.14–15 in its context.}
% mdparagraph:P1130:end

% output: NEW-U17-R03-O003
% mdparagraph:P1131:begin
\itempara{3}{Classify swā … ðæt 175.6–7.}
% mdparagraph:P1131:end

% output: NEW-U17-R03-O004
% mdparagraph:P1132:begin
\itempara{4}{Explain why the two ðæt clauses 175.7–8 do not assert that Apollonius actually is Apollo.}
% mdparagraph:P1132:end

\reserve{8\baselineskip}
% mdparagraph:P1133:begin
\subsection*{R04}
% mdparagraph:P1133:end

% item: NEW-U17-R04
% mdparagraph:P1134:begin
Analyse the supplied participial and predicate expressions without changing the source spelling.
% mdparagraph:P1134:end

% output: NEW-U17-R04-O001
% mdparagraph:P1135:begin
\itempara{1}{Explain swīgende 174.13 using original note 12a.}
% mdparagraph:P1135:end

% output: NEW-U17-R04-O002
% mdparagraph:P1136:begin
\itempara{2}{Explain the referent of ðā ymbsittendan 175.7.}
% mdparagraph:P1136:end

% output: NEW-U17-R04-O003
% mdparagraph:P1137:begin
\itempara{3}{Account for singular ungecnāwen wæs and ungewunelic 175.16.}
% mdparagraph:P1137:end

% output: NEW-U17-R04-O004
% mdparagraph:P1138:begin
\itempara{4}{Identify the subject of līcode 175.17 and distinguish it from him eallum.}
% mdparagraph:P1138:end

\reserve{8\baselineskip}
% mdparagraph:P1139:begin
\subsection*{R05}
% mdparagraph:P1139:end

% item: NEW-U17-R05
% mdparagraph:P1140:begin
Separate literal wording from the contextual performance sense in each expression, using the local help and original notes.
% mdparagraph:P1140:end

% output: NEW-U17-R05-O001
% mdparagraph:P1141:begin
\itempara{1}{Explain gefēoll on swēgcræft 174.16.}
% mdparagraph:P1141:end

% output: NEW-U17-R05-O002
% mdparagraph:P1142:begin
\itempara{2}{Explain plegode and forðtēah 175.15.}
% mdparagraph:P1142:end

% output: NEW-U17-R05-O003
% mdparagraph:P1143:begin
\itempara{3}{Explain the function of Apollines 175.8.}
% mdparagraph:P1143:end

\reserve{8\baselineskip}
% mdparagraph:P1144:begin
\subsection*{R06}
% mdparagraph:P1144:end

% item: NEW-U17-R06
% mdparagraph:P1145:begin
Explain the argument of the episode from the actual Old English.
% mdparagraph:P1145:end

% output: NEW-U17-R06-O001
% mdparagraph:P1146:begin
\itempara{1}{State the difference between the king’s reading of silence and Apollonius’s explanation.}
% mdparagraph:P1146:end

% output: NEW-U17-R06-O002
% mdparagraph:P1147:begin
\itempara{2}{Identify two kinds of evidence by which Apollonius answers the criticism of his conduct.}
% mdparagraph:P1147:end

% output: NEW-U17-R06-O003
% mdparagraph:P1148:begin
\itempara{3}{Explain how the final sentence combines novelty with approval.}
% mdparagraph:P1148:end

\reserve{8\baselineskip}
% mdparagraph:P1149:begin
\section*{Source and editorial notes}
% mdparagraph:P1149:end

\begingroup\fontsize{10}{12.8}\selectfont
\setlength{\parskip}{1.5pt}
% mdparagraph:P1150:begin
The English exercises are Cook, Exercises in Old English, joint Group XII, XIII, Exercises 37–43, printed pages 22–28 (scan pages 30–36). All original prompts and controls are retained. The Old English reading is Cook, A First Book in Old English, second edition, Reader XII, printed pages 174–175. Original reading note numbers restart on page 175; P174 and P175 below identify their page. The modern explanations and the lettered exercises are separate teaching material.
% mdparagraph:P1150:end

% mdparagraph:P1151:begin
The supplied exclusion phrase for the 144,000 follows the parallel recorded by Bosworth–Toller under þūsend, printed page 1082. Acute length marks are presented as macrons and a divided word is joined; final-o ðūsendo is retained. The supplied nīwne follows the long nīwe adjective’s wilde pattern described by Wright, pages 49, 213–214. In the although-expression, the dictionary’s ðēah-ðe is spaced as ðēah ðe. These limited presentation choices do not change the original English prompts or the reading text.
% mdparagraph:P1151:end

% mdparagraph:P1152:begin
\textsuperscript{P174N9}: Not in the Latin.
% mdparagraph:P1152:end

% mdparagraph:P1153:begin
\textsuperscript{P174N10}: See 79.
% mdparagraph:P1153:end

% mdparagraph:P1154:begin
\textsuperscript{P174N11}: See 140.
% mdparagraph:P1154:end

% mdparagraph:P1155:begin
\textsuperscript{P174N12}: Lat. arte musica.
% mdparagraph:P1155:end

% mdparagraph:P1156:begin
\textsuperscript{P174N12a}: For swīgiende.
% mdparagraph:P1156:end

% mdparagraph:P1157:begin
\textsuperscript{P174N13}: Lat. vituperas.
% mdparagraph:P1157:end

% mdparagraph:P1158:begin
\textsuperscript{P174N14}: See 196. d.
% mdparagraph:P1158:end

% mdparagraph:P1159:begin
\textsuperscript{P174N15}: Lat. incidit. Translate, has chanced.
% mdparagraph:P1159:end

% mdparagraph:P1160:begin
\textsuperscript{P174N16}: See 164. a.
% mdparagraph:P1160:end

% mdparagraph:P1161:begin
\textsuperscript{P174N17}: See 126.
% mdparagraph:P1161:end

% mdparagraph:P1162:begin
\textsuperscript{P175N1}: See 194, note.
% mdparagraph:P1162:end

% mdparagraph:P1163:begin
\textsuperscript{P175N2}: Apollo.
% mdparagraph:P1163:end

% mdparagraph:P1164:begin
\textsuperscript{P175N3}: We are reminded of Æn. II. 1, Conticuere omnes.
% mdparagraph:P1164:end

% mdparagraph:P1165:begin
\textsuperscript{P175N4}: To this sentence there corresponds in the Latin: — arripuit plectrum, animumque accommodat arti; cum chordis miscetur vox cantu modulata.
% mdparagraph:P1165:end

% mdparagraph:P1166:begin
\textsuperscript{P175N5}: Lat. deponens.
% mdparagraph:P1166:end

% mdparagraph:P1167:begin
\textsuperscript{P175N6}: The rest of this sentence paraphrases: induit statum comicum et inauditas actiones expressit, deinde tragicum.
% mdparagraph:P1167:end

% mdparagraph:P1168:begin
\textsuperscript{P175N7}: See 154. a, b.
% mdparagraph:P1168:end

% mdparagraph:P1169:begin
\textsuperscript{P175N8}: Lat. expressit.
% mdparagraph:P1169:end

% mdparagraph:P1170:begin
\textsuperscript{P175N9}: See 164. k.
% mdparagraph:P1170:end

% mdparagraph:P1171:begin
The short Latin glosses identify musical art (174 note 12), censure (174 note 13), chancing upon something (174 note 15), setting aside (175 note 5), and representing or performing (175 note 8). The allusion in 175 note 3 recalls everyone falling silent. The verse in note 4 describes taking the plectrum and mingling song with the strings; note 6 names comic and tragic performance. They illuminate the Old English scene without replacing its actual wording.
% mdparagraph:P1171:end

\endgroup
\end{document}
